\documentclass[11pt]{article}
\usepackage[margin=1in]{geometry}
\usepackage{amsmath,amsthm,mathtools}
\usepackage{amssymb}
\usepackage{hyperref}
\usepackage{enumitem}
\usepackage{booktabs}

\newtheorem{theorem}{Theorem}[section]
\newtheorem{lemma}[theorem]{Lemma}
\newtheorem{proposition}[theorem]{Proposition}
\newtheorem{corollary}[theorem]{Corollary}
\theoremstyle{definition}
\newtheorem{definition}[theorem]{Definition}
\newtheorem{remark}[theorem]{Remark}
\newtheorem{conjecture}[theorem]{Conjecture}

\newcommand{\R}{\mathbb{R}}
\newcommand{\C}{\mathbb{C}}
\newcommand{\T}{\mathbb{T}}
\newcommand{\Z}{\mathbb{Z}}
\newcommand{\E}{\mathbb{E}}
\newcommand{\Hdot}[1]{\dot{H}^{#1}}
\newcommand{\cS}{\mathcal{S}}
\newcommand{\cC}{\mathcal{C}}
\newcommand{\cB}{\mathcal{B}}
\newcommand{\cK}{\mathcal{K}}
\DeclareMathOperator{\re}{Re}
\DeclareMathOperator{\im}{Im}
\DeclareMathOperator{\corr}{corr}
\DeclareMathOperator{\sgn}{sgn}

\title{Phase-Conjugation Antisymmetry and Angular\\
Spreading in the 3D Navier--Stokes Nonlinearity}

\author{Ramiro Fontes\\
Quijotic Research\\
\texttt{ramirofontes@gmail.com}}

\date{March 2026}

\begin{document}
\maketitle

\begin{abstract}
We study the angular distribution of $\Hdot{1/2}$ energy
in Fourier space for the 3D incompressible Navier--Stokes
equations on $\T^3$.

\smallskip\noindent\textbf{Exact results.}
We prove that the per-mode nonlinear $\Hdot{1/2}$ energy
transfer $T(k;u)$ is (i)~exactly antisymmetric under
Fourier phase conjugation $\hat u(k)\mapsto\overline{\hat u(k)}$,
(ii)~independent of the Leray projector (pressure drops out),
and (iii)~degree-one homogeneous in $u$.
We also prove that for any random divergence-free
field whose distribution is invariant under
Fourier phase conjugation ($\hat u(k)\mapsto
\overline{\hat u(k)}$), both $\E[T(k;u)]$ and
$\E[\Sigma(u)]$ vanish.
All four results extend to the $\Hdot{s}$ spreading
functional $\Sigma^{(s)}$ for every $s\geq 0$
via the scaling identity
$T^{(s)}(k)=|k|^{2(s-1/2)}T(k)$, which also shows
$\sgn\,T^{(s)}(k) = \sgn\,T(k)$: the per-mode transfer
direction is Sobolev-level independent.
For fields with purely imaginary Fourier coefficients
($\hat u(k) = iV(k)$, $V\in\R^3$), we prove a real
transfer formula that eliminates all complex analysis
from the spreading problem.
We also derive an exact decomposition of $\Sigma$ in
terms of the sector excess $\varepsilon_j = p_j - 1/J$,
giving a rigorous basis for the numerically observed
constitutive law.

\smallskip\noindent\textbf{Taylor--Green computation.}
For the Taylor--Green vortex, we show by explicit
enumeration of all contributing triads and exact rational
arithmetic that the second-generation nonlinear forcing
at every original mode satisfies
$V(k)\cdot\widetilde R(k)=-1/512<0$, giving
$\cos\alpha(k) = -1$ (perfect anti-alignment).  The
computation exploits the purely imaginary Fourier structure
$\hat u(k)=iV(k)$ and proceeds by direct enumeration of
all 8 modes.

\smallskip\noindent\textbf{Numerical findings.}
Pseudo-spectral simulations at resolutions $128^3$--$256^3$
with six different initial conditions show that developed
Taylor--Green turbulence exhibits persistent angular
spreading ($\Sigma < 0$) at all Sobolev levels
$s\in[0,3/2]$, including the enstrophy ($s=1$).
The spreading is viscosity-independent to four significant
figures, dominated by cross-sector transfer
($|\Sigma_{\rm cross}|/|\Sigma_{\rm within}|\approx 830$),
and robust to 50\% phase destruction.
Triad phases lock to $\pm\pi/2$ in the NS solution ---
a direct consequence of the preserved purely imaginary
Fourier structure --- while random-phase fields and
concentrated jets produce uniform phase distributions.
All three random broadband initial conditions exhibit
enstrophy-level spreading ($\Sigma^{(1)}<0$) at every
snapshot (30/30), providing evidence for generality
beyond the TG symmetry class.
Surgically concentrating a developed NS solution
(boosting $\cC^{(1)}$ to $0.50$ while preserving
phase structure) produces monotone relaxation toward
isotropy at every snapshot (44/44), with the
stretching--spreading ratio $\mathcal{R}\in[0.6,2.2]$
and a near-universal constitutive law
$\Sigma^{(1)}\approx -0.47\,(\cC^{(1)}-1/J)$.
Randomization tests (same amplitudes, random
directions) produce $\Sigma^{(1)}<0$ at only
$48\%$--$50\%$, indistinguishable from the
conjugation-symmetric null: spreading is not a
property of the amplitude spectrum, the purely
imaginary structure, or the divergence-free
constraint, but of inter-scale correlations
created by the NS dynamics.
Evolving 30 random initial conditions under NS,
the fraction with $\Sigma^{(1)}<0$ rises from
$50\%$ to $100\%$ within one time unit.
Conditioning on initial sign reveals a contraction
asymmetry: all 15 seeds starting with $\Sigma^{(1)}>0$
flip to $\Sigma^{(1)}<0$ by $t=1$, while seeds
starting negative persist at $\geq 93\%$ through
$t=10$.  The creation result is confirmed at
$256^3$ (10~seeds), ruling out resolution artifacts,
and the creation timescale is independent of
viscosity across $\nu\in[0.002,0.02]$ (40~seeds),
confirming a purely nonlinear mechanism.

\smallskip\noindent\textbf{Heuristic mechanism and
regularity program.}
We propose a mechanism --- the ``$-i$ rotation'' in the
advection operator --- that explains both the exact TG
result and the numerical observations, and formulate a
conjecture at the enstrophy level ($s=1$) with a concrete
connection to the regularity problem: the NS nonlinearity
simultaneously spreads angular enstrophy and limits
vortex stretching through the same phase mechanism,
suggesting a self-regulating structure that may prevent
blow-up.  A homogeneity analysis reveals
a structural obstruction: the constitutive coefficient
$g^{(1)}=-\Sigma^{(1)}/(\cC^{(1)}-1/J)$ is degree~1
in $u$ (not dimensionless), so the constitutive law
controls the \emph{direction} of nonlinear transfer
but not the \emph{magnitude} of vortex stretching.
We resolve this by introducing the \emph{affine ratio}
$c^{(1)} = \sum\varepsilon_j T_j\,/\,
(T_{\rm tot}\sum\varepsilon_j^2)$,
a degree-0 (scale-invariant) quantity satisfying
$\Sigma^{(1)} = 2(c^{(1)}-1)\,T_{\rm tot}\,
\sum\varepsilon_j^2/E$, so that spreading with
$T_{\rm tot}>0$ is equivalent to $c^{(1)}<1$.
The homogeneity obstruction resides entirely in the
degree-1 prefactor $T_{\rm tot}/E$, not in $c^{(1)}$.
Numerical experiments show: $c^{(1)}<1$ along the TG
trajectory for $t\geq 1$ (values $0.18$--$0.69$);
$c^{(1)}$ remains $<1$ under surgical concentration up to
$\cC^{(1)}=0.124$; random divergence-free fields have
$c^{(1)}<1$ in only $12\%$ of cases with $T_{\rm tot}>0$
(the bound is dynamical, not kinematic); and the basin
$\{c^{(1)}<1\}$ is absorbing under NS evolution
($5/5$ seeds in each direction within one turnover time).
We prove that $c_{\rm cen}<1$ (the centroid-based
affine ratio) holds for all smooth NS solutions
with anisotropic initial data, for a time window
of order $1/(\nu\kappa_L)$ (the viscous timescale
of the energy-containing range).  The proof uses
only the NS equation, smoothness, and the Gronwall
inequality: viscous damping at rate $2\nu|k|^2$
suppresses the high-$k$ angular anisotropy of the
field $\kappa_H/\kappa_L$ times faster than the
low-$k$ anisotropy, creating spectral separation.
A concentration--spreading identity then gives
$\Sigma^{(1)} = d\sigma/dt
- 4\nu\sigma D(1-c_{\rm cen})/E$, where
$\sigma = \cC^{(1)}-1/J$ is the angular variance.
Since $c_{\rm cen}<1$ is proved, this yields:
(i)~$\Sigma^{(1)}<0$ at every concentration
equilibrium ($d\sigma/dt=0$);
(ii)~time-averaged $\Sigma^{(1)}<0$; and
(iii)~$c^{(1)}\geq 1$ cannot persist (any
excursion is self-correcting).
High-resolution sampling confirms that transient
excursions $c^{(1)}>1$ do occur during cascade
development ($\max c=1.23$ at $\nu=0.005$) but
are confined and self-limiting; the proved results
establish spreading as the dominant regime.
We identify the alignment deviation
$D = (\sum(\cos^2\theta_i-1/3)^2)^{1/2}$
between vorticity and strain eigendirections
and prove $|\alpha| \leq |S|\,D$.
Numerical experiments falsify the conjectured
transfer inequality ($D$ vs.\ $\cC^{(1)}$) but
reveal the actual depletion mechanism:
\emph{eigenvalue selection}.  NS turbulence
aligns vorticity with the intermediate strain
eigenvector ($\cos^2\theta_2\to 0.55$), which
carries the smallest eigenvalue, depleting
$81\%$ of the potential stretching
($\alpha/|S|\approx 0.19$).  Under surgical
concentration, $\alpha/|S|\approx
0.14\sqrt{1-\cC^{(1)}}$: stretching efficiency
decreases monotonically with angular
concentration.
\end{abstract}

\tableofcontents
\newpage

%======================================================================
\section{Introduction}
%======================================================================

The 3D incompressible Navier--Stokes equations on
$\T^3 = [0,2\pi]^3$,
\begin{equation}\label{eq:NS}
\partial_t u + (u\cdot\nabla)u + \nabla p = \nu\Delta u,
\quad \nabla\cdot u = 0, \quad u(0) = u_0,
\end{equation}
remain the subject of a central open problem: whether
smooth initial data produce globally smooth solutions
\cite{Fefferman,Leray}.  Partial results constrain
potential singularities \cite{ESS,BKM}, and Tao~\cite{Tao}
has shown that blow-up can occur for averaged versions of
the equations, suggesting that any regularity proof must
use specific structural properties of the nonlinearity.

Standard approaches to regularity treat the nonlinear
term $(u\cdot\nabla)u$ as the adversary.  In this paper
we present analytical and numerical evidence for a
different picture: the nonlinear term acts as a
\emph{regulator} of angular concentration in Fourier
space, spreading energy from concentrated to dispersed
angular directions.

The paper is organized into three layers of decreasing
rigor:
\begin{itemize}[leftmargin=2em]
\item \textbf{Sections~\ref{sec:general}--\ref{sec:TG}:}
  exact algebraic identities and a finite computation,
  all fully proved.
\item \textbf{Section~\ref{sec:results}:} numerical
  findings from pseudo-spectral simulations.
\item \textbf{Sections~\ref{sec:mechanism}--\ref{sec:conjecture}:}
  a heuristic mechanism, a conditional sign criterion,
  and a conjecture.
\end{itemize}
We make every effort to keep these layers clearly
separated.

\subsection{Setup and definitions}

We decompose $S^2$ (the sphere of wavevector directions
$\hat k = k/|k|$) into $J$ approximately equal-area
sectors $\{\Omega_j\}_{j=1}^J$ using a Fibonacci spiral
partition (Section~\ref{sec:method}).

\begin{definition}[Sector energies and concentration]
\label{def:sector}
For a smooth divergence-free field $u$ on $\T^3$
with Fourier coefficients $\hat u(k)$, $k\in\Z^3$:
\begin{align}
E_j &= \sum_{\hat k\in\Omega_j}|k|\,|\hat u(k)|^2,
\qquad E = \sum_j E_j, \qquad p_j = E_j/E,\\
\cC(u) &= \sum_j p_j^2 = \|p\|_{\ell^2}^2
\quad\text{(angular concentration)},\\
\cS(u) &= -\sum_j p_j\log p_j
\quad\text{(angular entropy)}.
\end{align}
Here $E_j$ is the $\Hdot{1/2}$ energy in sector
$\Omega_j$, and $p_j$ is the fraction of total
$\Hdot{1/2}$ energy in that sector.  Note $\cC\in[1/J,1]$
with $\cC=1/J$ for isotropic fields.
\end{definition}

\begin{definition}[Nonlinear energy transfer]
\label{def:transfer}
Let $B(u,u) = -P[(u\cdot\nabla)u]$ denote the
Leray-projected advection, where
$P_{jm}(k) = \delta_{jm}-k_jk_m/|k|^2$ is the Leray
projector in Fourier space.  The Fourier representation is
\begin{equation}\label{eq:B-fourier}
\widehat{B(u,u)}_j(k)
= -P_{jm}(k)\sum_{\substack{p+q=k\\p,q\in\Z^3}}
iq_l\,\hat u_l(p)\,\hat u_m(q),
\end{equation}
with repeated indices summed ($j,l,m\in\{1,2,3\}$).
The per-mode $\Hdot{1/2}$ energy transfer is
\begin{equation}\label{eq:Tk}
T(k;u) = 2\,\re\bigl[|k|\,\hat u(k)^*
\cdot\widehat{B(u,u)}(k)\bigr].
\end{equation}
The angular spreading functional is:
\begin{equation}\label{eq:Sigma}
\Sigma(u) = \frac{d\cC}{dt}\bigg|_{\rm NL}
= \frac{2}{E}\sum_j p_j\Bigl(\sum_{\hat k\in\Omega_j}T(k;u)
- p_j\sum_k T(k;u)\Bigr).
\end{equation}
Spreading corresponds to $\Sigma < 0$ (concentration
decreasing); concentration to $\Sigma > 0$.
\end{definition}

\subsection{Summary of results}

\noindent\textbf{Exact algebraic results}
(Section~\ref{sec:general}):
\begin{enumerate}[leftmargin=2em,label=(A\arabic*)]
\item Phase-conjugation antisymmetry:
  $T(k;\bar u) = -T(k;u)$
  (Theorem~\ref{thm:antisym}).
\item Conjugation-symmetric vanishing: $\E[T(k;u)] = 0$
  and $\E[\Sigma(u)]=0$
  (Theorem~\ref{thm:gaussian}).
\item Pressure invariance
  (Proposition~\ref{prop:pressure}).
\item Homogeneity: $\Sigma(\lambda u)=\lambda\Sigma(u)$
  (Proposition~\ref{prop:homog}).
\item Sobolev scaling:
  $T^{(s)}(k)=|k|^{2(s-1/2)}T(k)$ and
  $\sgn\,T^{(s)}(k)=\sgn\,T(k)$; all of
  A1--A4 extend to $\Sigma^{(s)}$ for all
  $s\geq 0$
  (Proposition~\ref{prop:sobolev-scaling},
  Corollary~\ref{cor:universal-s}).
\item Real transfer formula: for purely imaginary
  fields $\hat u=iV$, $T(k)$ is a real sum with
  no complex quantities
  (Proposition~\ref{prop:real-transfer}).
\item Linear response: exact decomposition
  $\Sigma = (2/E)\sum\varepsilon_j T_j
  - (2T_{\rm tot}/E)\sum\varepsilon_j^2$
  (Proposition~\ref{prop:linear-response}).
\end{enumerate}

\noindent\textbf{Taylor--Green computation}
(Section~\ref{sec:TG}):
\begin{enumerate}[leftmargin=2em,label=(B\arabic*)]
\item $V(k)\cdot\widetilde R(k)=-1/512 < 0$ at all
  TG modes (Theorem~\ref{thm:TG}).
\item Along the entire TG orbit ($t\geq 0$), the
  transfer $T^{(s)}(k)$ is a real sum at every
  Sobolev level
  (Corollary~\ref{cor:TG-real-orbit}).
\end{enumerate}

\noindent\textbf{Numerical findings}
(Section~\ref{sec:results}):
\begin{enumerate}[leftmargin=2em,label=(N\arabic*)]
\item $\Sigma(u) < 0$ throughout TG-generated turbulence.
\item Spreading rate independent of $\nu$.
\item Phase conjugation flips $\Sigma$ exactly.
\item Random div-free fields: 49\% spreading.
\item Spreading survives 50\% phase randomization.
\item Spreading established within one time step.
\item Cross-sector transfer dominates by $\mathbf{830\times}$.
\item High-$p$ sectors have $\cos\alpha < 0$ (depletion).
\item Multiple ICs: TG spreads; random fields $\Sigma\approx 0$;
  concentrated jet concentrates despite mode-level depletion.
\item Triad phases lock to $\pm\pi/2$ (purely imaginary)
  in the NS solution, with a $4\!:\!1$ bias toward $+\pi/2$.
  Random-phase and jet fields show uniform distributions.
\item $\Sigma/(\cC-1/J)$ independent of $\nu$ to four
  significant figures across $\nu\in[0.002,0.02]$.
\item $\Sigma^{(s)}<0$ for all $s\in[0,3/2]$: the nonlinear
  term spreads angular concentration at every Sobolev level,
  including the enstrophy ($s=1$).
\item Vortex stretching weakly correlated with enstrophy
  concentration ($\corr=+0.14$), but full nonlinear transfer
  produces net spreading ($\corr=-0.43$): cross-sector
  redistribution dominates.
\item $\Sigma^{(1)}/(\cC^{(1)}-1/J)$ independent of $\nu$:
  enstrophy-level constitutive relation is purely nonlinear.
\item All three random broadband ICs have $\Sigma^{(1)}<0$ at
  every snapshot (30/30): enstrophy spreading is more universal
  than $\Hdot{1/2}$ spreading.
\item Mode-level depletion $\kappa_{\rm align}$ undergoes a
  sharp sign transition at $\cC^{(1)}\approx 0.15$: the
  spreading mechanism engages only above a concentration
  threshold.
\item Enstrophy spreading survives 50\% phase randomization
  ($\Sigma^{(1)}<0$ at $\geq 10/11$ concentration levels
  for $f\leq 0.5$).
\item Surgically concentrated fields ($\cC^{(1)}$ boosted
  to $0.20$--$0.50$ while preserving NS phase structure)
  exhibit $\Sigma^{(1)}<0$ at all 44 snapshots ($256^3$),
  with monotone relaxation toward isotropy.
\item The stretching--spreading ratio
  $\mathcal{R}=|\mathcal{W}|/(|\Sigma^{(1)}|
  \cdot\|\omega\|^2)\in[0.6,2.2]$ in the surgical runs,
  \emph{decreasing} with concentration.
\item The constitutive coefficient $g^{(1)} =
  -\Sigma^{(1)}/(\cC^{(1)}-1/J) \approx 0.47\pm 0.04$
  across $\cC^{(1)}\in[0.13, 0.50]$, identifying the
  analytical target for the regularity program.
\item Randomization tests: random $V$ (same $|V(k)|$,
  random directions), isotropized $V$, and permuted $V$
  all give $\Sigma^{(1)}<0$ at $48\%$--$50\%$ (50~seeds
  each), indistinguishable from the
  conjugation-symmetric null.  Spreading is not a
  static property of the amplitude spectrum.
\item The NS cascade creates inter-scale
  anti-correlation ($C(s_1,s_2)<0$ for $s_1\neq s_2$,
  ratio $-263\times$); random fields show the opposite
  sign ($+144\times$).  This correlation fingerprint is
  the dynamical origin of spreading.
\item Evolving 30 random ICs under NS ($128^3$): the fraction with
  $\Sigma^{(1)}<0$ rises from $50\%$ ($t=0$) to $100\%$
  ($t=1$) and stabilizes at $90\%$+ ($t=10$).  The NS
  flow creates spreading from generic initial data within
  one nonlinear turnover time.
\item Contraction asymmetry: all 15 seeds starting with
  $\Sigma^{(1)}>0$ flip to $\Sigma^{(1)}<0$ by $t=1$;
  seeds starting negative persist at $\geq 93\%$
  through $t=10$.  The basin $\{\Sigma^{(1)}<0\}$ is
  nearly absorbing under NS evolution.
\item High-resolution creation ($256^3$, 10~seeds):
  the $50\%\to 100\%$ transition replicates at full
  resolution, with mean $g^{(1)}$ peaking at $0.43$
  ($t=2$).  The creation phenomenon is not a
  resolution artifact.
\item Viscosity independence of creation ($128^3$,
  10~seeds~$\times$~4~viscosities,
  $\nu\in\{0.002,0.005,0.01,0.02\}$):
  all four viscosities reach $\geq 90\%$ spreading
  by $t=1$, and the creation timescale is $O(1)$
  regardless of $\nu$.  Spreading creation is a purely
  nonlinear phenomenon.
\item Viscous--nonlinear balance: the viscous
  contribution to $d\cC^{(1)}/dt$ is positive
  (concentrating) and comparable in magnitude to
  $|\Sigma^{(1)}|$ at all times
  ($|\text{visc}/\text{NL}|\approx 1.0$).
  The constitutive coefficient $g^{(1)}$ is degree~1
  in $u$, not dimensionless: it varies by $2\times$
  across viscosities and $27\%$ between $128^3$ and
  $256^3$ at the same $\nu$.
\item The affine ratio
  $c^{(1)}=\sum\varepsilon_jT_j/(T_{\rm tot}\sum\varepsilon_j^2)$
  is degree~0 (confirmed to machine precision under
  $u\mapsto\lambda u$) and satisfies $c^{(1)}<1$ along
  the TG trajectory for $t\geq 1$
  (Section~\ref{sec:affine-experiments}).
\item $c^{(1)}$ increases with surgical concentration
  ($0.62$ at $\cC^{(1)}=0.037$ to $0.96$ at
  $\cC^{(1)}=0.124$) but remains $<1$ at all tested
  levels.  At extreme single-sector concentration
  ($p_{\max}>0.99$), $c^{(1)}$ can exceed~$1$, but
  NS evolution corrects this within $t=0.25$.
\item The bound $c^{(1)}<1$ is dynamical, not kinematic:
  only $12\%$ of random divergence-free fields with
  $T_{\rm tot}>0$ have $c^{(1)}<1$ (50~seeds).
\item The basin $\{c^{(1)}<1\}$ is absorbing under NS
  dynamics: $5/5$ seeds starting with $c^{(1)}>1$ drop
  below~$1$ by $t=0.5$; $5/5$ seeds starting $c^{(1)}<1$
  persist through $t=1.5$.
\item Spectral separation from viscous balance
  (Theorem~\ref{thm:viscous-sep}): $c_{\rm cen}<1$
  for all smooth NS solutions with anisotropic
  initial data, for $t\leq O(1/(\nu\kappa_L))$.
  The proof uses only the NS equation, Gronwall,
  and the two-shell algebraic identity.
\item Concentration--spreading identity
  (Theorem~\ref{thm:conc-spread}):
  $\Sigma^{(1)} = d\sigma/dt
  - 4\nu\sigma D(1-c_{\rm cen})/E$.
  With $c_{\rm cen}<1$ proved: $\Sigma^{(1)}<0$ at
  every concentration equilibrium
  (Theorem~\ref{thm:equil}), time-averaged
  $\Sigma^{(1)}<0$ (Theorem~\ref{thm:time-avg}),
  and $c^{(1)}\geq 1$ cannot persist
  (Remark~\ref{rem:self-reinforce}).
\end{enumerate}

\noindent\textbf{Heuristic mechanism and conjecture}
(Sections~\ref{sec:mechanism}--\ref{sec:conjecture}):
\begin{enumerate}[leftmargin=2em,label=(H\arabic*)]
\item Cross-sector depletion via the $-i$ rotation
  (Section~\ref{sec:mechanism}).
\item A conditional sign criterion linking the forward
  cascade phase structure to angular spreading
  (Section~\ref{sec:conditional}).
\item Conjecture~\ref{conj:enstrophy}: $\Sigma^{(1)}(u)<0$
  for all smooth NS solutions with $\cC^{(1)} > 1/J$
  (enstrophy-level spreading).
\item The affine ratio $c^{(1)}$: proved $c_{\rm cen}<1$
  always, $c^{(1)}<1$ at concentration equilibria,
  time-averaged $\Sigma^{(1)}<0$, and $c\geq 1$
  self-correcting.  Transient excursions $c>1$ during
  cascade development (N36: max $c=1.23$ at $\nu=0.005$)
  are confined and do not persist
  (Remark~\ref{conj:affine},
  Section~\ref{sec:concentration-dynamics}).
\item A regularity program based on the coupled
  stretching--spreading mechanism
  (Remark~\ref{rem:regularity-final}).
\item Homogeneity obstruction: the constitutive
  coefficient is degree~1, not dimensionless,
  preventing direct reduction to regularity
  (Section~\ref{sec:homogeneity}).
\item Structural dictionary between the NS spreading
  conjecture and Bakry--\'Emery curvature, with
  the key disanalogy at homogeneity
  (Section~\ref{sec:bakry-emery}).
\item Alignment deviation identity and the transfer
  inequality conjecture --- tested and falsified;
  the actual depletion mechanism is eigenvalue
  selection ($\alpha/|S|\approx 0.19$, N27--N28)
  (Section~\ref{sec:alignment-experiment}).
\end{enumerate}

%======================================================================
\section{Exact Algebraic Results}
\label{sec:general}
%======================================================================

\subsection{Phase-conjugation antisymmetry}

\begin{definition}[Phase conjugation]
\label{def:conj}
For a divergence-free field $u$ on $\T^3$, define
$\bar u$ by $\hat{\bar u}(k) = \overline{\hat u(k)}$.
For a real field this gives $\bar u(x)=u(-x)$ (spatial
parity inversion).  All $\Hdot{s}$ norms, sector energies,
and probabilities $p_j$ are unchanged.
\end{definition}

\begin{remark}
This operation is Fourier phase conjugation (equivalently,
spatial parity), \emph{not} dynamical time reversal.
For the Euler equations, the standard time-reversal
involves $u(x,t)\mapsto -u(x,-t)$; for Navier--Stokes,
viscosity already breaks time-reversal symmetry at the
PDE level.  The identity below is therefore an algebraic
oddness property of the trilinear transfer functional
under parity, not a statement about physical reversibility.
\end{remark}

\begin{theorem}[Phase-conjugation antisymmetry]
\label{thm:antisym}
For any smooth divergence-free $u$ on $\T^3$ and
$k\in\Z^3\setminus\{0\}$:
\begin{equation}
T(k;\bar u) = -T(k;u).
\end{equation}
Consequently $\Sigma(\bar u) = -\Sigma(u)$.
\end{theorem}

\begin{proof}
From~\eqref{eq:Tk} and~\eqref{eq:B-fourier}:
\begin{equation}
T(k;u) = 2|k|\re\Bigl[\hat u(k)^*\cdot
\Bigl(-P(k)\sum_{p+q=k}iq_l\hat u_l(p)\hat u(q)
\Bigr)\Bigr].
\end{equation}
Since $\re[\bar w\cdot(-iz)] = \im[\bar w\cdot z]$,
we may write $T(k;u) = 2|k|\,\im[Z(k;u)]$ where
\begin{equation}\label{eq:Z}
Z(k;u) = \sum_{j,l,m}P_{jm}(k)\sum_{p+q=k}
q_l\,\hat u_j(k)^*\hat u_l(p)\hat u_m(q).
\end{equation}
Under $\hat u(k)\to\overline{\hat u(k)}$: the conjugation
reverses $\hat u_j(k)^*\to\hat u_j(k)$ and converts
$\hat u_l(p)\to\overline{\hat u_l(p)}$,
$\hat u_m(q)\to\overline{\hat u_m(q)}$.  Since $P_{jm}(k)$
and $q_l$ are real:
\begin{equation}
Z(k;\bar u) = \overline{Z(k;u)},
\end{equation}
whence $\im[Z(k;\bar u)]=-\im[Z(k;u)]$ and
$T(k;\bar u)=-T(k;u)$.

For $\Sigma$: since $p_j(\bar u) = p_j(u)$ and
$E(\bar u) = E(u)$, the formula~\eqref{eq:Sigma}
gives $\Sigma(\bar u) = -\Sigma(u)$.
\end{proof}

\begin{corollary}[Pure oddness]
\label{cor:odd}
The transfer is purely odd under phase conjugation:
$T^+ := (T(k;u)+T(k;\bar u))/2 \equiv 0$.
\end{corollary}

\subsection{Vanishing for conjugation-symmetric fields}

\begin{theorem}[Vanishing under conjugation symmetry]
\label{thm:gaussian}
Let $u$ be a random smooth divergence-free field on $\T^3$
with $\E[|T(k;u)|]<\infty$ and $\E[|\Sigma(u)|]<\infty$,
whose distribution is invariant under phase conjugation:
$\bar u \stackrel{d}{=} u$, where
$\hat{\bar u}(k) = \overline{\hat u(k)}$.

Then $\E[T(k;u)]=0$ for every $k\neq 0$, and
$\E[\Sigma(u)]=0$.
\end{theorem}

\begin{proof}
By Theorem~\ref{thm:antisym}: $T(k;\bar u) = -T(k;u)$
for every realization.  Since $\bar u\stackrel{d}{=}u$:
\begin{equation}
\E[T(k;u)] = \E[T(k;\bar u)] = \E[-T(k;u)]
= -\E[T(k;u)].
\end{equation}
Hence $\E[T(k;u)]=0$.

For the spreading functional: under phase conjugation,
$E(\bar u) = E(u)$, $p_j(\bar u)=p_j(u)$, and
$T(k;\bar u) = -T(k;u)$.  From the
definition~\eqref{eq:Sigma}, $\Sigma$ is linear in the
values $\{T(k)\}$ with coefficients depending only on
$\{p_j\}$ and $E$.  Since all coefficients are unchanged
under conjugation while each $T(k)$ flips sign:
$\Sigma(\bar u) = -\Sigma(u)$.  By the same distributional
argument, $\E[\Sigma]=0$.
\end{proof}

\begin{remark}[Applicable models]
The hypothesis $\bar u\stackrel{d}{=}u$ holds for any
isotropic Gaussian field, and more generally for any
random field whose Fourier amplitudes are deterministic
and whose phases (on a half-lattice) are i.i.d.\ uniform
or invariant under $\phi\mapsto -\phi$.  It fails for
fields with systematic phase bias, such as developed
NS solutions --- which is precisely why those solutions
can have $\Sigma\neq 0$.
\end{remark}

\subsection{Pressure invariance}

\begin{proposition}[Pressure drops out]
\label{prop:pressure}
$T(k;u) = 2|k|\re[\hat u(k)^*\cdot
(-\widehat{(u\cdot\nabla)u}(k))]$.
That is, $T(k)$ is the same whether computed with the
Leray-projected forcing $\hat B(k)$ or with the raw
advection $-\widehat{(u\cdot\nabla)u}(k)$.
\end{proposition}

\begin{proof}
The Leray projection removes the component parallel
to $k$:
$\hat B(k) = -\widehat{(u\cdot\nabla)u}(k)
+ k\,\beta(k)$
for some scalar $\beta(k)$.  Since $k\cdot\hat u(k)=0$
(divergence-free), the correction $k\,\beta$ is
orthogonal to $\hat u(k)^*$ and drops out of the
inner product.
\end{proof}

\subsection{Factored form and homogeneity}

\begin{proposition}[Factored form]
\label{prop:factored}
\begin{equation}\label{eq:factored}
T(k;u) = 2|k|\,\im\Bigl[
\sum_{p+q=k}(q\cdot\hat u(p))(\hat u(k)^*\cdot\hat u(q))
\Bigr].
\end{equation}
\end{proposition}

\begin{proof}
By Proposition~\ref{prop:pressure}, we may use the raw
advection.  Its $j$-th Fourier component is
$\widehat{(u\cdot\nabla)u}_j(k)
= \sum_{p+q=k}\sum_l iq_l\hat u_l(p)\hat u_j(q)$
(no Leray projector, replacing $P_{jm}$ with $\delta_{jm}$).
Computing $T = 2|k|\im[\cdots]$ and contracting:
$\sum_j \hat u_j(k)^*\hat u_j(q) = \hat u(k)^*\cdot\hat u(q)$
and
$\sum_l q_l\hat u_l(p) = q\cdot\hat u(p)$.
\end{proof}

\begin{proposition}[Homogeneity]
\label{prop:homog}
$\Sigma(\lambda u) = \lambda\,\Sigma(u)$ for $\lambda>0$.
\end{proposition}

\begin{proof}
$T(k;\lambda u) = \lambda^3 T(k;u)$ (cubic in $\hat u$).
$E(\lambda u) = \lambda^2 E(u)$.  $p_j(\lambda u)=p_j(u)$.
From~\eqref{eq:Sigma}: $\Sigma$ picks up $\lambda^3$
from $T$ and $\lambda^{-2}$ from $1/E$, giving
net degree $3-2=1$.
\end{proof}

\subsection{Sobolev extension}

The $\Hdot{s}$ transfer $T^{(s)}(k;u) = 2|k|^{2s}\re[\hat u(k)^*\cdot\hat B(k)]$ (Definition~\ref{def:Hs-spreading}) differs from $T(k;u) = T^{(1/2)}(k;u)$ only by the positive real weight $|k|^{2s}$.  This gives an immediate scaling identity.

\begin{proposition}[Sobolev scaling]
\label{prop:sobolev-scaling}
For all $s\geq 0$ and $k\in\Z^3\setminus\{0\}$:
\begin{equation}\label{eq:Ts-scaling}
T^{(s)}(k;u) = |k|^{2(s-1/2)}\,T(k;u).
\end{equation}
In particular, $\sgn\, T^{(s)}(k) = \sgn\, T(k)$ for all $s$.
\end{proposition}

\begin{proof}
$T^{(s)}(k) = 2|k|^{2s}\re[\hat u^*\cdot\hat B]
= |k|^{2s-1}\cdot 2|k|\re[\hat u^*\cdot\hat B]
= |k|^{2(s-1/2)}\,T(k)$.
\end{proof}

\begin{corollary}[Universal algebraic results]
\label{cor:universal-s}
The following hold for all $s\geq 0$:
\begin{enumerate}[leftmargin=2em,label=(\roman*)]
\item Phase-conjugation antisymmetry:
  $T^{(s)}(k;\bar u) = -T^{(s)}(k;u)$
  and $\Sigma^{(s)}(\bar u) = -\Sigma^{(s)}(u)$.
\item Conjugation-symmetric vanishing:
  $\E[T^{(s)}(k)]=0$ and $\E[\Sigma^{(s)}]=0$.
\item Homogeneity: $\Sigma^{(s)}(\lambda u)=\lambda\,\Sigma^{(s)}(u)$.
\item Factored form:
  $T^{(s)}(k;u) = 2|k|^{2s}\im\bigl[\sum_{p+q=k}
  (q\cdot\hat u(p))(\hat u(k)^*\cdot\hat u(q))\bigr]$.
\end{enumerate}
\end{corollary}

\begin{proof}
Each follows from the base result
(Theorems~\ref{thm:antisym}--\ref{thm:gaussian},
Propositions~\ref{prop:factored}--\ref{prop:homog})
via the scaling identity~\eqref{eq:Ts-scaling}:
the factor $|k|^{2(s-1/2)}$ is positive, real,
unchanged by conjugation, and homogeneous of
degree $0$ in $u$.  The conjugation-symmetric
vanishing for $\Sigma^{(s)}$ follows because
$p_j^{(s)}(\bar u)=p_j^{(s)}(u)$ and
$E^{(s)}(\bar u)=E^{(s)}(u)$.
\end{proof}

\begin{remark}
Corollary~\ref{cor:universal-s} means the exact
null hypothesis ($\E[\Sigma^{(s)}]=0$ for
conjugation-symmetric ensembles) and the algebraic
oddness structure hold at every regularity level,
including the enstrophy ($s=1$).  This is the
algebraic foundation for
Conjecture~\ref{conj:enstrophy}.
\end{remark}

\subsection{Real transfer formula for purely imaginary fields}

\begin{proposition}[Real transfer formula]
\label{prop:real-transfer}
If $\hat u(k) = iV(k)$ with $V(k)\in\R^3$ for all $k$,
then
\begin{equation}\label{eq:real-transfer}
T(k;u) = 2|k|\sum_{p+q=k}
(q\cdot V(p))\,(V(k)\cdot V(q)),
\end{equation}
and more generally
\begin{equation}\label{eq:real-transfer-s}
T^{(s)}(k;u) = 2|k|^{2s}\sum_{p+q=k}
(q\cdot V(p))\,(V(k)\cdot V(q)).
\end{equation}
The transfer is a purely real sum with no
complex-valued quantities.
\end{proposition}

\begin{proof}
Substitute $\hat u(p)=iV(p)$,
$\hat u(k)^*=-iV(k)$, $\hat u(q)=iV(q)$
into the factored form
(Proposition~\ref{prop:factored}):
\begin{align}
(q\cdot\hat u(p))(\hat u(k)^*\cdot\hat u(q))
&= \bigl(iq\cdot V(p)\bigr)
  \bigl(-iV(k)\cdot iV(q)\bigr) \notag\\
&= i\,(q\cdot V(p))\cdot(-i^2)(V(k)\cdot V(q))
\notag\\
&= i\,(q\cdot V(p))(V(k)\cdot V(q)).
\end{align}
The sum over triads is
$i\sum_{p+q=k}(q\cdot V(p))(V(k)\cdot V(q))$,
which is purely imaginary (the sum is real).
Taking $\im[\cdot]$ extracts the real sum,
giving~\eqref{eq:real-transfer}.
The $\Hdot{s}$ version follows from
Proposition~\ref{prop:sobolev-scaling}.
\end{proof}

\begin{remark}[Reduction to real analysis]
\label{rem:real-analysis}
By Corollary~\ref{cor:persistence}, the TG
solution satisfies $\hat u(k,t)=iV(k,t)$ with
$V(k,t)\in\R^3$ for all $k$ and all $t\geq 0$.
Along the entire TG orbit, the spreading question
reduces to the sign of the real
sum~\eqref{eq:real-transfer}
aggregated by sector.  All complex analysis is
eliminated: the problem is entirely about the
geometry of the real vectors $V(k,t)$ in the
triad structure $p+q=k$.
\end{remark}

\subsection{Structure of $\Sigma$ near isotropy}

\begin{proposition}[Linear response]
\label{prop:linear-response}
Let $\varepsilon_j = p_j - 1/J$ denote the sector
excess, so $\sum_j\varepsilon_j = 0$ and
$\cC - 1/J = \sum_j\varepsilon_j^2 + (2/J)\sum_j\varepsilon_j
= \sum_j\varepsilon_j^2$.
Then:
\begin{equation}\label{eq:Sigma-linear}
\Sigma(u) = \frac{2}{E}\sum_j \varepsilon_j\, T_j
\;-\; \frac{2\,T_{\rm tot}}{E}\sum_j\varepsilon_j^2,
\end{equation}
where $T_j = \sum_{\hat k\in\Omega_j}T(k;u)$ is the
sector transfer and
$T_{\rm tot} = \sum_j T_j$.  In particular, if
$|\varepsilon_j|\ll 1/J$ (near isotropy):
\begin{equation}\label{eq:Sigma-leading}
\Sigma(u) \approx \frac{2}{E}
\sum_j \varepsilon_j\, T_j.
\end{equation}
Spreading ($\Sigma<0$) thus requires the sector transfer
$T_j$ to be negatively correlated with the sector excess
$\varepsilon_j$: sectors with above-average energy must
receive below-average nonlinear transfer.
\end{proposition}

\begin{proof}
From~\eqref{eq:Sigma}:
\begin{align}
\frac{E}{2}\Sigma
&= \sum_j p_j\bigl(T_j - p_j\,T_{\rm tot}\bigr)
= \sum_j p_j T_j - T_{\rm tot}\sum_j p_j^2
\notag\\
&= \sum_j (1/J+\varepsilon_j)T_j
  - T_{\rm tot}(1/J+\varepsilon_j)^2 \notag\\
&= \frac{T_{\rm tot}}{J} + \sum_j\varepsilon_j T_j
  - T_{\rm tot}\Bigl(\frac{1}{J}
  +\sum_j\varepsilon_j^2\Bigr) \notag\\
&= \sum_j\varepsilon_j T_j
  - T_{\rm tot}\sum_j\varepsilon_j^2,
\end{align}
using $\sum_j 1/J = 1$, $\sum_j\varepsilon_j = 0$,
and $\sum_j p_j^2 = 1/J + \sum_j\varepsilon_j^2$.
\end{proof}

\begin{remark}
\label{rem:constitutive-form}
Equation~\eqref{eq:Sigma-linear} is exact, not
perturbative.  The leading-order
approximation~\eqref{eq:Sigma-leading} drops the
term $T_{\rm tot}\sum\varepsilon_j^2$, which is
second-order in the sector excess.
This gives a rigorous basis for the constitutive
law $\Sigma \approx -g(\cC-1/J)$ observed
numerically (N11, N14, N19): the law holds whenever
$T_j$ is approximately linear in $\varepsilon_j$
with negative slope, i.e.,
$T_j \approx T_{\rm tot}/J - g_0\,\varepsilon_j$
for some $g_0>0$.
\end{remark}

\begin{definition}[Affine ratio]
\label{def:affine-ratio}
For a divergence-free field $u$ with $T_{\rm tot}\neq 0$
and $\cC^{(s)}\neq 1/J$, define the \emph{affine ratio}
\begin{equation}\label{eq:affine-ratio}
c^{(s)}(u) := \frac{\sum_j \varepsilon_j^{(s)}\, T_j^{(s)}}
{T_{\rm tot}^{(s)}\,\sum_j (\varepsilon_j^{(s)})^2}.
\end{equation}
Proposition~\ref{prop:linear-response} gives the exact identity
\begin{equation}\label{eq:Sigma-via-c}
\frac{E^{(s)}}{2}\,\Sigma^{(s)}
= \bigl(c^{(s)} - 1\bigr)\,T_{\rm tot}^{(s)}\,
\sum_j(\varepsilon_j^{(s)})^2.
\end{equation}
When $T_{\rm tot}^{(s)}>0$ (net enstrophy production
at Sobolev level~$s$), spreading $\Sigma^{(s)}<0$ is
equivalent to $c^{(s)}<1$.
\end{definition}

\begin{proposition}[Degree-0 property of the affine ratio]
\label{prop:c-degree-zero}
Under $u\mapsto\lambda u$ ($\lambda>0$):
$c^{(s)}(\lambda u) = c^{(s)}(u)$.
That is, $c^{(s)}$ is degree~$0$ (scale-invariant).
\end{proposition}

\begin{proof}
Both the numerator $\sum_j\varepsilon_j T_j$ and the
denominator $T_{\rm tot}\sum_j\varepsilon_j^2$ are
degree~$3$ in $u$ (the transfer is trilinear;
$\varepsilon_j$ is degree~$0$ since $p_j$ is).
The ratio cancels $\lambda^3$.
\end{proof}

\begin{remark}[Affine interpretation]
\label{rem:affine-interp}
If the sector transfer is approximately affine in the
sector fraction, $T_j \approx b + m\,p_j$ for
constants $b$ (intercept) and $m$ (slope), then
$\sum_j\varepsilon_j T_j = m\sum_j\varepsilon_j^2$
(using $\sum_j\varepsilon_j=0$), so $c = m/T_{\rm tot}$.
Since $T_{\rm tot} = Jb + m$ (summing over $J$ sectors
with $\sum p_j = 1$), $c<1$ is equivalent to $b>0$:
every sector receives a positive baseline of nonlinear
transfer from the external field, regardless of its own
enstrophy content.  The quantity $b/T_{\rm tot}$
measures the strength of this baseline relative to the
total transfer.
\end{remark}

%======================================================================
\section{The Taylor--Green Computation}
\label{sec:TG}
%======================================================================

\subsection{Fourier structure}

The Taylor--Green vortex
$u_0 = (\sin x_1\cos x_2\cos x_3,\;
-\cos x_1\sin x_2\cos x_3,\;0)$
has Fourier support on $\cK = \{(\pm 1,\pm 1,\pm 1)\}$
(8 wavevectors):
\begin{equation}\label{eq:TG-fourier}
\hat u_1(k) = \frac{-ik_1}{8},\qquad
\hat u_2(k) = \frac{ik_2}{8},\qquad
\hat u_3(k) = 0,\qquad k\in\cK.
\end{equation}
All coefficients are purely imaginary:
$\hat u(k) = iV(k)$ where
$V(k) = (-k_1/8,\,k_2/8,\,0) \in \R^3$.

\begin{lemma}[Properties of $V$]
\label{lem:TG-imag}
For all $k\in\cK$: (i)~$k\cdot V(k) = 0$
(divergence-free, since $k_1^2=k_2^2=1$);
(ii)~$V(-k)=-V(k)$ (reality);
(iii)~$|V(k)|^2 = 1/32$ (uniform amplitude).
\end{lemma}

\subsection{Phase preservation by the nonlinearity}

\begin{proposition}
\label{prop:phase-preserve}
If $\hat u(k)=iV(k)$ with $V(k)\in\R^3$ for all $k$,
then $\hat B(k)=iR(k)$ with $R(k)\in\R^3$.
\end{proposition}

\begin{proof}
The raw advection is
$\widehat{(u\cdot\nabla)u}_j(k')
= \sum_{p+q=k'}\sum_l iq_l\cdot iV_l(p)\cdot iV_j(q)
= i^3\sum_{p+q=k'}\sum_l q_lV_l(p)V_j(q)$.
Since $i^3=-i$, this equals $-iS_j(k')$ with
$S_j(k') = \sum_{p+q=k'}(q\cdot V(p))V_j(q)\in\R$.
The Leray projector $P(k')$ has real entries, and
$B = -P(\text{adv})$, giving
$\hat B_j(k') = -P_{jm}(-iS_m) = iP_{jm}S_m = iR_j$.
\end{proof}

\begin{corollary}\label{cor:persistence}
The purely imaginary structure is preserved by the full NS
evolution: nonlinearity preserves it
(Proposition~\ref{prop:phase-preserve}), and viscosity
$-\nu|k|^2\hat u = -\nu|k|^2\cdot iV$ preserves it by
linearity.
\end{corollary}

\subsection{First generation: no self-transfer}

\begin{proposition}\label{prop:gen1}
$B(u_0,u_0)$ has Fourier support on
$\cK'=\{k\in\Z^3\setminus\{0\}: k_i\in\{0,\pm 2\}\}$,
which is disjoint from $\cK$.  Therefore $T(k;u_0)=0$
for all $k\in\cK$.
\end{proposition}

\begin{proof}
For $p,q\in\cK$: each component of $p+q$ lies in
$\{-2,0,2\}$, never in $\{\pm 1\}$.  So the advection
creates no modes at the original wavevectors.
Since there is no forcing at $k\in\cK$,
$\hat B(k) = 0$ and $T(k) = 0$.
\end{proof}

\subsection{First-generation modes}

Of the 26 nonzero wavevectors in $\cK'$, exactly 8
receive nonzero forcing (the others vanish by
cancellation).  By explicit computation:

\begin{proposition}\label{prop:gen1-modes}
The nonzero first-generation modes are
$\hat B(k') = iR(k')$ with:
\begin{center}
\begin{tabular}{@{}ll@{}}
\toprule
$k'$ & $R(k')$ \\
\midrule
$(2,0,2)$   & $(1/32,\,0,\,-1/32)$ \\
$(2,0,-2)$  & $(1/32,\,0,\,1/32)$ \\
$(-2,0,2)$  & $(-1/32,\,0,\,-1/32)$ \\
$(-2,0,-2)$ & $(-1/32,\,0,\,1/32)$ \\
$(0,2,2)$   & $(0,\,1/32,\,-1/32)$ \\
$(0,2,-2)$  & $(0,\,1/32,\,1/32)$ \\
$(0,-2,2)$  & $(0,\,-1/32,\,-1/32)$ \\
$(0,-2,-2)$ & $(0,\,-1/32,\,1/32)$ \\
\bottomrule
\end{tabular}
\end{center}
\end{proposition}

\begin{proof}
For each $k'\in\cK'$, compute
$S_j(k')=\sum_{p+q=k'}(q\cdot V(p))V_j(q)$ by
enumerating all pairs $p,q\in\cK$ with $p+q=k'$,
then apply $R = P(k')S$.

We show the computation for $k'=(2,0,2)$.  The
constraint $p+q=(2,0,2)$ with $p,q\in\{\pm 1\}^3$
forces $p_1=q_1=1$, $p_3=q_3=1$, and $p_2=-q_2$.
The two contributing pairs are
$(p,q)\in\{((1,s,1),(1,-s,1)):s=\pm 1\}$.

For $s=1$: $q\cdot V(p) = (1,-1,1)\cdot(-1/8,1/8,0)
= -1/4$.  $V(q) = V(1,-1,1) = (-1/8,-1/8,0)$.
Contribution: $(-1/4)(-1/8,-1/8,0) = (1/32,1/32,0)$.

For $s=-1$: $q\cdot V(p) = (1,1,1)\cdot(-1/8,-1/8,0)
= -1/4$.  $V(q) = V(1,1,1) = (-1/8,1/8,0)$.
Contribution: $(-1/4)(-1/8,1/8,0) = (1/32,-1/32,0)$.

Sum: $S(2,0,2) = (1/16, 0, 0)$.
Leray projection: $|k'|^2=8$, $k'\cdot S = 2/16 = 1/8$.
$R = S - k'(k'\cdot S)/|k'|^2
= (1/16,0,0) - (2,0,2)(1/8)/8
= (1/16-1/32,\; 0,\; -1/32)
= (1/32, 0, -1/32)$.

The remaining 7 modes are computed by the same procedure
and verified by the companion script
\texttt{tg\_exact\_symbolic.py} using exact rational
arithmetic.  All 8 values appear in the table in
Proposition~\ref{prop:gen1-modes}.

A companion script (\texttt{tg\_exact\_symbolic.py})
verifies all 8 modes using Python's \texttt{Fraction}
arithmetic with zero floating-point error.
\end{proof}

\subsection{Second generation: exact anti-alignment}

After one time step $\Delta t$, the field has modes
at both $\cK$ (amplitude $O(1)$) and $\cK'$ (amplitude
$O(\Delta t)$).  Second-generation forcing at $k\in\cK$
comes from triads with one leg in $\cK'$ and one in $\cK$.

\begin{definition}[Second-generation forcing]
\label{def:gen2}
For $k\in\cK$, define:
\begin{equation}\label{eq:Rtilde}
\widetilde R_j(k) = \sum_mP_{jm}(k)\sum_{\substack{
q\in\cK\\k-q\in\cK'}}
\bigl[(q\cdot R(k\!-\!q))V_m(q)
+((k\!-\!q)\cdot V(q))R_m(k\!-\!q)\bigr],
\end{equation}
where the first term is the Type-A contribution
(first-generation mode advects, original is advected)
and the second is Type-B (original advects,
first-generation mode is advected).
The second-generation forcing is
$\hat B^{(2)}(k) = i\Delta t\,\widetilde R(k)$.
\end{definition}

\begin{theorem}[Taylor--Green anti-alignment]
\label{thm:TG}
For all $k\in\cK$:
\begin{equation}
\widetilde R(k) = -\frac{1}{16}V(k),
\end{equation}
and therefore:
\begin{equation}
V(k)\cdot\widetilde R(k) = -\frac{|V(k)|^2}{16}
= -\frac{1}{512} < 0.
\end{equation}
The forcing is exactly anti-parallel to the velocity:
$\cos\alpha(k) = -1$ at every TG mode.
\end{theorem}

\begin{proof}
The computation is finite: for each $k\in\cK$, we
enumerate all $q\in\cK$ with $k-q$ among the 8 nonzero
first-generation modes (Proposition~\ref{prop:gen1-modes}),
evaluate the Type-A and Type-B contributions
in~\eqref{eq:Rtilde}, apply the Leray projector, and sum.
We give the full computation for $k_0=(1,1,1)$ and
tabulate all 8 results.

\smallskip\noindent\textbf{Computation at $k=(1,1,1)$.}
The constraint $q\in\cK$, $k-q\in\cK'$ with $k-q$
among the 8 nonzero first-generation modes yields
exactly two contributing triads:
$q=(-1,1,-1)$ with $p'=(2,0,2)$, and
$q=(1,-1,-1)$ with $p'=(0,2,2)$.
(All other $q\in\cK$ give $k-q$ outside the nonzero
first-generation set.)

For $q=(-1,1,-1)$, $p'=(2,0,2)$:
Type~A: $(q\cdot R(p'))V(q)
= ((-1,1,-1)\cdot(1/32,0,-1/32))\,V(-1,1,-1)
= (-1/32+1/32)\,V(-1,1,-1) = 0$.
Type~B: $(p'\cdot V(q))\,R(p')
= ((2,0,2)\cdot(1/8,1/8,0))\,(1/32,0,-1/32)
= (1/4)(1/32,0,-1/32)
= (1/128,0,-1/128)$.

For $q=(1,-1,-1)$, $p'=(0,2,2)$:
Type~A: $(q\cdot R(p'))V(q)
= ((1,-1,-1)\cdot(0,1/32,-1/32))\,V(1,-1,-1)
= (-1/32+1/32)\,V(1,-1,-1) = 0$.
Type~B: $(p'\cdot V(q))\,R(p')
= ((0,2,2)\cdot(-1/8,-1/8,0))\,(0,1/32,-1/32)
= (-1/4)(0,1/32,-1/32)
= (0,-1/128,1/128)$.

Sum before Leray projection:
$(1/128,-1/128,-1/128+1/128)=(1/128,-1/128,0)$.
Leray projection at $k=(1,1,1)$:
$|k|^2=3$, $k\cdot(1/128,-1/128,0)=0$, so
$\widetilde R(1,1,1)=(1/128,-1/128,0)$.

Check: $-V(1,1,1)/16 = -(-1/8,1/8,0)/16
= (1/128,-1/128,0)$. \checkmark

Inner product: $V\cdot\widetilde R
= (-1/8)(1/128)+(1/8)(-1/128) = -1/512$. \checkmark

Cosine: $(V\cdot\widetilde R)^2 = 1/262144
= |V|^2|\widetilde R|^2$, confirming
$\cos^2\alpha=1$, hence $\cos\alpha=-1$.

\smallskip\noindent\textbf{All 8 modes.}
The same enumeration procedure applied to each
$k\in\cK$ gives:
\begin{center}
\begin{tabular}{@{}lll@{}}
\toprule
$k$ & $\widetilde R(k)$ & $-V(k)/16$ \\
\midrule
$(1,1,1)$    & $(1/128,\,-1/128,\,0)$ & $(1/128,\,-1/128,\,0)$ \\
$(1,1,-1)$   & $(1/128,\,-1/128,\,0)$ & $(1/128,\,-1/128,\,0)$ \\
$(1,-1,1)$   & $(1/128,\,1/128,\,0)$  & $(1/128,\,1/128,\,0)$  \\
$(1,-1,-1)$  & $(1/128,\,1/128,\,0)$  & $(1/128,\,1/128,\,0)$  \\
$(-1,1,1)$   & $(-1/128,\,-1/128,\,0)$ & $(-1/128,\,-1/128,\,0)$ \\
$(-1,1,-1)$  & $(-1/128,\,-1/128,\,0)$ & $(-1/128,\,-1/128,\,0)$ \\
$(-1,-1,1)$  & $(-1/128,\,1/128,\,0)$ & $(-1/128,\,1/128,\,0)$ \\
$(-1,-1,-1)$ & $(-1/128,\,1/128,\,0)$ & $(-1/128,\,1/128,\,0)$ \\
\bottomrule
\end{tabular}
\end{center}
In every case $\widetilde R(k)=-V(k)/16$, giving
$V\cdot\widetilde R = -1/512$ and $\cos\alpha=-1$.

The companion script \texttt{tg\_exact\_symbolic.py}
reproduces the full 8-mode enumeration using Python's
\texttt{Fraction} arithmetic with zero floating-point
error, independently verifying every entry in both
tables (Proposition~\ref{prop:gen1-modes} and above).
\end{proof}

\begin{remark}[The $i^2=-1$ mechanism]
\label{rem:i-squared}
The anti-alignment has a transparent algebraic origin.
The TG coefficients carry one factor of $i$
($\hat u = iV$).  The first-generation forcing picks up
$i^3 = -i$ from two coefficients and the advection
operator's $iq_l$, then the sign in $B=-P(\text{adv})$
restores $+i$ (Proposition~\ref{prop:phase-preserve}).
When first-generation modes feed back through a second
advection, another factor of $i$ accumulates: the total
phase for two-step feedback is $i\cdot i = -1$ relative
to the original velocity.
\end{remark}

\begin{remark}[Persistence]
By Corollary~\ref{cor:persistence}, $\hat u(k,t) = iV(k,t)$
for all $t\geq 0$.  This gives $\cos^2\alpha = 1$ at
every mode (the inner product $V\cdot\widetilde R$ is
always real).  The sign $\cos\alpha = -1$ at the original
TG modes is proved at leading order in $\Delta t$ by
Theorem~\ref{thm:TG}; numerical simulations confirm it
persists throughout the evolution.
\end{remark}

\begin{corollary}[Real transfer along the TG orbit]
\label{cor:TG-real-orbit}
Along the TG evolution, $\hat u(k,t)=iV(k,t)$ for
all $t\geq 0$
(Corollary~\ref{cor:persistence}).
By Proposition~\ref{prop:real-transfer}, the transfer
at every mode and every Sobolev level is the real sum
\begin{equation}\label{eq:TG-orbit-transfer}
T^{(s)}(k;u(t)) = 2|k|^{2s}\sum_{p+q=k}
(q\cdot V(p,t))\,(V(k,t)\cdot V(q,t))
\end{equation}
for all $t\geq 0$, $s\geq 0$, $k\neq 0$.
The spreading question for the TG orbit at any
Sobolev level reduces to the sign of this real sum
aggregated by angular sector.
\end{corollary}

%======================================================================
\section{Numerical Methods}
\label{sec:method}
%======================================================================

\subsection{Solver}

We solve~\eqref{eq:NS} in Fourier space on the
$N^3$ grid ($N=128,256$) using an integrating-factor
Runge--Kutta scheme (IF-RK4).  The linear viscous
operator $-\nu|k|^2$ is integrated exactly via the
factor $e^{-\nu|k|^2\Delta t}$, and the nonlinear term
$\hat B(u,u)$ is advanced by classical RK4 in the
integrating-factor frame.  Aliasing is removed by the
$2/3$ rule (zeroing modes with $|k|>N/3$).  The time
step is $\Delta t = 10^{-3}$ for all runs reported here;
stability was verified by halving $\Delta t$ and checking
that all diagnostics change by less than $10^{-6}$.

\subsection{Initial conditions}

We test six initial conditions spanning different symmetry
classes and concentration levels:
\begin{enumerate}[leftmargin=2em]
\item \textbf{Taylor--Green} (TG):
  $u_0 = (\sin x_1\cos x_2\cos x_3,\,
  -\cos x_1\sin x_2\cos x_3,\,0)$.
  Purely imaginary Fourier structure; 8 active modes.
\item \textbf{ABC flow}:
  $u_0 = (A\sin x_3+C\cos x_2,\,B\sin x_1+A\cos x_3,\,
  C\sin x_2+B\cos x_1)$ with $(A,B,C)=(1,0.8,0.6)$.
  A Beltrami field ($\nabla\times u = u$); 12 active modes.
\item \textbf{Random broadband} (3 seeds):
  isotropic energy spectrum $\propto k^4e^{-k^2/32}$
  peaked at $|k|\approx 4$, independent random phases,
  Leray-projected.  Concentration $\cC\approx 0.07$
  (nearly isotropic).
\item \textbf{Concentrated jet}:
  energy in a narrow cone (half-angle $0.3$\,rad) around
  the $x_3$-axis at $|k|\approx 3$, random phases,
  Leray-projected.  Concentration $\cC\approx 0.48$
  (strongly anisotropic).
\end{enumerate}
Each is evolved to $t=5$ ($5{,}000$ time steps) with
diagnostics recorded every $\Delta t=0.5$.

\subsection{Angular sectors}

We use a Fibonacci spiral partition of $S^2$ into $J$
equal-area sectors with centers
$\hat n_j = (\sin\vartheta_j\cos\psi_j,
\sin\vartheta_j\sin\psi_j,\cos\vartheta_j)$,
$\vartheta_j = \arccos(1-2(j+1/2)/J)$,
$\psi_j = 2\pi j/\varphi$, where
$\varphi = (1+\sqrt{5})/2$ is the golden ratio.
Each wavevector $k\neq 0$ is assigned to its nearest
sector center.  Results are reported for
$J\in\{16,32,64,128\}$; the spreading sign $\Sigma<0$
is independent of $J$.

\subsection{Diagnostics}

\begin{definition}[Phase alignment]
\label{def:phase-align}
For a mode $k$ with $\hat u(k)\neq 0$ and
$\hat B(k)\neq 0$:
\begin{equation}
\cos\alpha(k) = \frac{\re[\hat u(k)^*\cdot\hat B(k)]}
{|\hat u(k)|\,|\hat B(k)|}.
\end{equation}
Since $T(k) = 2|k||\hat u||\hat B|\cos\alpha$:
$\cos\alpha>0$ means the forcing amplifies mode $k$,
$\cos\alpha<0$ means it depletes.
\end{definition}

At each diagnostic snapshot we compute: sector energies
$E_j$ and probabilities $p_j$; nonlinear forcing
$\hat B(k)$; per-mode transfer $T(k)$ via~\eqref{eq:Tk};
per-sector transfer $dE_j/dt|_{\rm NL}
=\sum_{\hat k\in\Omega_j}T(k)$;
spreading functional $\Sigma$ via~\eqref{eq:Sigma};
Pearson correlation $\corr(p_j, dE_j|_{\rm NL})$;
per-sector mean phase alignment
$\langle\cos\alpha\rangle_j$; and the within-sector /
cross-sector transfer decomposition (described below).

\subsection{Cross-sector decomposition}

The nonlinear forcing at mode $k\in\Omega_j$ is
decomposed as $\hat B(k) = \hat B^{\rm within}(k)
+ \hat B^{\rm cross}(k)$, where
$\hat B^{\rm within}(k)$ includes only triads $p+q=k$
with the advecting velocity $\hat u(p)$ in the same
sector as $k$ ($\hat p\in\Omega_j$), and
$\hat B^{\rm cross}(k)$ includes the rest.  In
practice, we compute $\hat B^{\rm within}$ by masking
the advecting field to sector $j$ before evaluating
the pseudospectral convolution, then subtracting from
the full $\hat B$.  The induced spreading decomposition
is $\Sigma = \Sigma_{\rm within}+\Sigma_{\rm cross}$,
where $\Sigma_{\rm within}$ (resp.\ $\Sigma_{\rm cross}$)
is computed from~\eqref{eq:Sigma} using only the
within-sector (resp.\ cross-sector) transfers
$T^{\rm within}(k) = 2|k|\re[\hat u^*\cdot
\hat B^{\rm within}]$ (resp.\ $T^{\rm cross}$).

\subsection{Phase surgery}

We test the role of Fourier phase structure by modifying
a developed NS solution:
full phase randomization (random i.i.d.\ uniform
phases, Leray re-projected);
partial randomization (fraction $f$ of modes
randomized);
phase conjugation ($\hat u\mapsto\overline{\hat u}$).
Each modified field is Leray-projected before computing
$\Sigma$.

\subsection{Validation}

The solver conserves $L^2$ energy to machine precision
in the inviscid limit ($\nu=0$, short time).  The TG
solution at $t=0$ reproduces the exact Fourier coefficients
from~\eqref{eq:TG-fourier}, and the first-generation
forcing matches Proposition~\ref{prop:gen1-modes} to 15
digits.  The antisymmetry $T(k;\bar u)=-T(k;u)$ holds to
$10^{-16}$ relative error for all developed solutions.

%======================================================================
\section{Numerical Results}
\label{sec:results}
%======================================================================

All results below use $N=128$, $\nu=0.01$, $J=32$
unless otherwise stated, with the TG vortex evolved to
$t=5$.

\subsection{Persistent spreading along TG orbit (N1)}

Along the TG trajectory at $N=128$, $\nu=0.01$:
at all 1{,}001 diagnostic snapshots from $t=0.01$ to
$t=20$, $\Sigma < 0$ with
$\corr(p_j,dE_j|_{\rm NL}) = -0.56\pm 0.03$.

\subsection{Viscosity independence (N2)}

\begin{table}[h]
\centering
\caption{Spreading diagnostics at $t=5$, $N=128$, $J=64$,
TG initial data.  The nonlinear correlation is independent
of $\nu$ to three significant figures.}
\label{tab:nu}
\begin{tabular}{@{}ccr@{}}
\toprule
$\nu$ & $\corr(p_j,dE_j|_{\rm NL})$
  & $|{\rm NL}/{\rm viscous}|$ \\
\midrule
$10^{-2}$          & $-0.562$ & $610\times$ \\
$5\times 10^{-3}$  & $-0.562$ & $2{,}640\times$ \\
$2\times 10^{-3}$  & $-0.563$ & $13{,}200\times$ \\
\bottomrule
\end{tabular}
\end{table}

\subsection{Phase-conjugation antisymmetry (N3)}

\begin{center}
\begin{tabular}{@{}lcc@{}}
\toprule
& $\corr$ & $\Sigma$ \\
\midrule
NS solution     & $-0.6022$ & $-1.199\times 10^{-3}$ \\
Conjugated      & $+0.6022$ & $+1.199\times 10^{-3}$ \\
\bottomrule
\end{tabular}
\end{center}
Confirmation of Theorem~\ref{thm:antisym} to machine
precision.

\subsection{Non-universality for random fields (N4)}

Random div-free fields with the same energy spectrum:
22/45 spread (49\%), $\corr\in[-0.7,+0.9]$,
mean $+0.07$.  Consistent with
Theorem~\ref{thm:gaussian}: conjugation-symmetric fields
have $\E[\Sigma]=0$.

\subsection{Robustness (N5)}

Starting from a developed NS solution and randomizing
a fraction $f$ of Fourier phases:
\begin{center}
\begin{tabular}{@{}ccc@{}}
\toprule
$f$ & mean $\corr$ & spread \\
\midrule
$0\%$ & $-0.60$ & 5/5 \\
$30\%$ & $-0.29$ & 5/5 \\
$50\%$ & $-0.34$ & 5/5 \\
$70\%$ & $-0.04$ & 3/5 \\
$100\%$ & $+0.16$ & 1/5 \\
\bottomrule
\end{tabular}
\end{center}

\subsection{Onset (N6)}

$\corr = -0.61$ at $t = 0.01$ (one time step from TG),
stable to $t = 20$.

\subsection{Cross-sector dominance (N7)}

$\Sigma_{\rm within} = -1.4\times 10^{-6}$,
$\Sigma_{\rm cross} = -1.2\times 10^{-3}$
(ratio $\mathbf{830\times}$).  Within-sector self-interaction
\emph{concentrates} ($\corr(p_j, dE^{\rm within}_j)=+0.74$);
cross-sector interaction \emph{spreads}
($\corr(p_j, dE^{\rm cross}_j) = -0.69$).

\subsection{Phase alignment (N8)}

$\corr(p_j,\langle\cos\alpha\rangle_j) = -1.000$.
For this TG-based solution at $t=5$, the four high-$p$
sectors have $\cos\alpha = -1.000$ and low-$p$ sectors
have $\cos\alpha = +1.000$ (to numerical precision).
Theorem~\ref{thm:TG} proves exact anti-alignment at
leading order in $\Delta t$; the persistence to $t=5$
is a numerical observation, consistent with the preserved
imaginary structure (Corollary~\ref{cor:persistence})
but not directly proved by Theorem~\ref{thm:TG}.
For non-TG initial data, partial anti-alignment
($\cos\alpha<0$ on average in high-$p$ sectors)
is expected.

\subsection{Sector count independence}

\begin{center}
\begin{tabular}{@{}ccc@{}}
\toprule
$J$ & $\corr$ & $\Sigma<0$? \\
\midrule
$16$  & $-0.93$ & Yes \\
$32$  & $-0.61$ & Yes \\
$64$  & $-0.56$ & Yes \\
$128$ & $-0.54$ & Yes \\
\bottomrule
\end{tabular}
\end{center}

\subsection{Multiple initial conditions (N9)}
\label{sec:multi-ic}

To test whether spreading is specific to TG symmetry,
we repeat the spreading diagnostic for all six initial
conditions (Section~\ref{sec:method}), each evolved to
$t=5$ at $N=128$, $\nu=0.01$, $J=32$.

\begin{table}[h]
\centering
\caption{Spreading across initial conditions.
$n_{\rm spread}$: snapshots with $\Sigma<0$ out of 11.}
\label{tab:multi-ic}
\begin{tabular}{@{}lcccl@{}}
\toprule
IC & $\Sigma(t\!=\!5)$ & $\corr$ & $n_{\rm spread}$
  & status \\
\midrule
Taylor--Green    & $-1.20\times 10^{-3}$ & $-0.69$ & 11/11
  & spreads \\
ABC flow         & $\approx 0$     & $+1.00$ & 2/11
  & degenerate \\
Random (seed 42) & $+9.1\times 10^{-12}$ & $+0.13$ & 0/11
  & concentrates \\
Random (seed 99) & $-1.3\times 10^{-12}$ & $-0.18$ & 11/11
  & spreads \\
Random (seed 7)  & $-8.4\times 10^{-12}$ & $-0.01$ & 11/11
  & spreads \\
Concentrated jet & $+4.1\times 10^{-12}$ & $+0.92$ & 0/11
  & concentrates \\
\bottomrule
\end{tabular}
\end{table}

Three qualitatively different regimes emerge:

\medskip\noindent\textbf{Strong spreading (TG).}
The TG orbit shows $\Sigma \sim -10^{-3}$, orders of
magnitude larger than all other ICs, with perfect
anti-alignment $\cos\alpha=-1$ in high-$p$ sectors.
This is the regime where the $-i$ mechanism operates
at full strength, supported by the purely imaginary
phase structure.

\medskip\noindent\textbf{Noise-level (random broadband).}
The three random ICs have $|\Sigma|\sim 10^{-12}$,
nine orders of magnitude smaller than TG, and the sign
varies by seed (2/3 spread, 1/3 concentrate).  These
fields are nearly isotropic ($\cC\approx 0.07$, close
to $1/J=0.03$) and have conjugation-symmetric phases,
so $\E[\Sigma]=0$ by Theorem~\ref{thm:gaussian}.
The observed values are consistent with zero-mean
fluctuations.

\medskip\noindent\textbf{Concentration (jet).}
The concentrated jet shows $\Sigma > 0$ at every snapshot,
with concentration \emph{increasing} over time
($\cC: 0.475\to 0.478$).  Remarkably, the jet still
exhibits the $-i$ mechanism at the mode level:
$\corr(p_j,\langle\cos\alpha\rangle_j) = -0.70$
(the forcing depletes high-$p$ modes).  But $\Sigma>0$
despite this mode-level depletion.  The explanation:
the jet's random phases lack the cascade phase structure
needed for the cross-sector transfer to dominate.
The nonlinear term creates new angular structure
(at new scales and directions) faster than it depletes
the existing concentration.

\medskip\noindent\textbf{Degenerate (ABC).}
The ABC flow is a Beltrami field satisfying
$\nabla\times u=u$, making it a stationary Euler solution.
The nonlinear term $(u\cdot\nabla)u = (\nabla\times u)
\times u + \nabla(|u|^2/2)$ reduces to a pressure gradient
for Beltrami fields, so $B(u,u)\approx 0$ and $\Sigma$
is at machine zero.  This is not a counterexample but
a degenerate case where the nonlinearity barely acts.

\medskip
The jet result is the most informative: it shows that
the $-i$ mechanism (mode-level depletion) is necessary
but not sufficient for spreading.  Angular spreading
$\Sigma<0$ additionally requires the NS cascade phase
structure, which is present in developed TG turbulence
but absent in random-phase concentrated fields.


\subsection{Triad phase locking (N10)}
\label{sec:triad-phase}

The factored form (Proposition~\ref{prop:factored}) writes
the per-mode transfer as
$T(k) = 2|k|\,\im\bigl[\sum_{p+q=k}z(k;p,q)\bigr]$,
where the triad integrand is
\begin{equation}\label{eq:triad-integrand}
z(k;p,q) = (q\cdot\hat u(p))\,
(\hat u(k)^*\cdot\hat u(q)).
\end{equation}
We sample $2\times 10^5$ random triads from the set of
active modes (those with $|\hat u(k)|^2$ above $10^{-12}$
of the peak) and compute the magnitude-weighted phase
histogram of $\arg z$.

For the TG-generated NS solution at $t=10$,
$\nu=0.002$, $N=128$ (648 active modes,
$1.8\times 10^5$ nontrivial triads), the triad phases
concentrate sharply at $\arg z \approx +\pi/2$ (81\%
of total weight) and $\arg z\approx -\pi/2$ (18\%),
with all other bins essentially empty
(Table~\ref{tab:triad-phase}).  For comparison, the
random-phase field (same amplitudes) and the
concentrated jet both produce uniform histograms
indistinguishable from i.i.d.\ phases
($\approx 4\%$ per bin).

\begin{table}[h]
\centering
\caption{Triad phase statistics.
$\langle\sin\theta\rangle$ is the magnitude-weighted
mean of $\sin(\arg z)$; positive means net amplification
of total $\Hdot{1/2}$ energy by the nonlinear term.}
\label{tab:triad-phase}
\begin{tabular}{@{}lccc@{}}
\toprule
Field & triads & $\langle\sin\theta\rangle$
  & phase structure \\
\midrule
NS solution    & $1.8\times 10^5$ & $+0.63$
  & locked to $\pm\pi/2$ \\
Random phases  & $1.8\times 10^5$ & $-0.08$
  & uniform \\
Conc.\ jet     & $4.7\times 10^5$ & $-0.00$
  & uniform \\
\bottomrule
\end{tabular}
\end{table}

The phase locking has a transparent algebraic origin.
By Corollary~\ref{cor:persistence}, $\hat u(k)=iV(k)$
with $V(k)\in\R^3$ persists along the TG orbit.
Substituting into~\eqref{eq:triad-integrand}:
\begin{equation}\label{eq:triad-imaginary}
z(k;p,q) = (q\cdot iV(p))\,(-iV(k)\cdot iV(q))
= i\,(q\cdot V(p))(V(k)\cdot V(q)),
\end{equation}
which is purely imaginary, with
$\arg z = +\pi/2$ when $(q\cdot V(p))(V(k)\cdot V(q))>0$
and $\arg z = -\pi/2$ otherwise.  The observed $4\!:\!1$
bias toward $+\pi/2$ reflects the fact that the dominant
triads have both factors positive, so the aggregate
nonlinear term amplifies total $\Hdot{1/2}$ energy
($\langle\sin\theta\rangle>0$).

The key point for spreading is that this amplification
is \emph{anti-correlated with angular concentration}:
modes in low-$p_j$ sectors receive more nonlinear
amplification than those in high-$p_j$ sectors
(consistent with $\corr(p_j,dE_j|_{\rm NL})<0$
in~N1), producing $\Sigma<0$ even though the total
energy transfer is positive.

The uniform phase distribution for the random-phase
field and the concentrated jet confirms that the
$\pm\pi/2$ locking is a consequence of the
NS-preserved purely imaginary structure, not a
generic property of the trilinear form.  This
provides a positive characterization of the
``cascade phase structure'' invoked in N9 and
Conjecture~\ref{conj:enstrophy}: it is the persistence
of $\hat u(k)\approx iV(k)$, which locks triad
phases to the imaginary axis.


\subsection{Viscosity-independent constitutive relation (N11)}
\label{sec:constitutive}

We evolve the TG vortex from $t=0$ to $t=10$ at four
viscosities $\nu\in\{0.02, 0.01, 0.005, 0.002\}$
($N=128$, $J=32$), recording $\cC(u)$ and $\Sigma(u)$
at intervals $\Delta t=0.1$.

\begin{table}[h]
\centering
\caption{Mean spreading diagnostics over
$t\in[5,10]$.  The ratio $\Sigma/(\cC-1/J)$ is
independent of $\nu$ to four significant figures.}
\label{tab:constitutive}
\begin{tabular}{@{}cccc@{}}
\toprule
$\nu$ & $\langle\cC\rangle$
  & $\langle\Sigma/(\cC\!-\!1/J)\rangle$
  & std \\
\midrule
$0.02$  & $0.243$ & $-8.245\times 10^{-3}$ & $1.50\times 10^{-3}$ \\
$0.01$  & $0.243$ & $-8.256\times 10^{-3}$ & $1.50\times 10^{-3}$ \\
$0.005$ & $0.243$ & $-8.261\times 10^{-3}$ & $1.50\times 10^{-3}$ \\
$0.002$ & $0.243$ & $-8.264\times 10^{-3}$ & $1.51\times 10^{-3}$ \\
\bottomrule
\end{tabular}
\end{table}

The ratio $\Sigma/(\cC-1/J)$ agrees across all four
viscosities to four significant figures at every time
snapshot (Table~\ref{tab:constitutive}), confirming
that the spreading is purely nonlinear: viscous
dissipation contributes negligibly.  However, the ratio
evolves from $-2.2\times 10^{-3}$ at $t=2$ to
$-1.1\times 10^{-2}$ at $t=10$ (a factor of~5), while
$(\cC-1/J)$ itself decreases by only 5\%, indicating
that $\Sigma$ depends on an additional state variable
beyond~$\cC$.  Identifying the correct constitutive law
\begin{equation}\label{eq:constitutive}
\Sigma(u) = -g(u)\,(\cC(u)-1/J),
\end{equation}
where $g(u)>0$ is a viscosity-independent functional
of $u$, is a key open problem for the regularity
program.


\subsection{$\Hdot{s}$ spreading for $s\in[0,3/2]$ (N12)}
\label{sec:Hs-spreading}

\begin{definition}[$\Hdot{s}$ sector energies and spreading]
\label{def:Hs-spreading}
For $s\geq 0$, define the $\Hdot{s}$ sector energies
\begin{equation}
E_j^{(s)} = \sum_{\hat k\in\Omega_j}
|k|^{2s}\,|\hat u(k)|^2,
\qquad p_j^{(s)} = E_j^{(s)}/E^{(s)},
\qquad \cC^{(s)}(u) = \sum_j (p_j^{(s)})^2,
\end{equation}
and the $\Hdot{s}$ transfer and spreading functional
\begin{equation}
T^{(s)}(k;u) = 2|k|^{2s}\,\re[\hat u(k)^*\cdot
\hat B(k)],
\qquad
\Sigma^{(s)}(u) = \frac{d\cC^{(s)}}{dt}
\bigg|_{\rm NL}.
\end{equation}
The case $s=1/2$ recovers the original definitions.
For $s=1$, the sector energies $E_j^{(1)}$ are the
angular components of the enstrophy.
\end{definition}

We evolve the TG vortex at $N=128$, $\nu=0.005$
from $t=0$ to $t=10$, computing $\Sigma^{(s)}$ for
$s\in\{0,\,1/4,\,1/2,\,3/4,\,1,\,3/2\}$.

\begin{table}[h]
\centering
\caption{$\Hdot{s}$ spreading diagnostics, averaged
over $t\in[5,10]$, $J=32$.  The spreading strengthens
with increasing~$s$.}
\label{tab:Hs}
\begin{tabular}{@{}ccccl@{}}
\toprule
$s$ & $\langle\cC^{(s)}\rangle$
  & $\langle\Sigma^{(s)}\rangle$
  & $\langle\corr\rangle$
  & $\Sigma^{(s)}<0$\;always? \\
\midrule
$0$   & $0.245$ & $-1.11\times 10^{-3}$ & $-0.75$ & Yes \\
$1/4$ & $0.244$ & $-1.41\times 10^{-3}$ & $-0.68$ & Yes \\
$1/2$ & $0.243$ & $-1.79\times 10^{-3}$ & $-0.60$ & Yes \\
$3/4$ & $0.241$ & $-2.25\times 10^{-3}$ & $-0.51$ & Yes \\
$1$   & $0.238$ & $-2.82\times 10^{-3}$ & $-0.43$ & Yes \\
$3/2$ & $0.231$ & $-4.37\times 10^{-3}$ & $-0.30$ & Yes \\
\bottomrule
\end{tabular}
\end{table}

$\Sigma^{(s)}<0$ holds at every snapshot and every~$s$
(Table~\ref{tab:Hs}), and the spreading intensifies
with~$s$: $|\Sigma^{(3/2)}|/|\Sigma^{(1/2)}|\approx 2.4$.
The $s$-independence of the sign is a consequence of
Proposition~\ref{prop:sobolev-scaling}:
$\sgn\, T^{(s)}(k) = \sgn\, T(k)$ for all $s$, so the
$\pm\pi/2$ triad phase locking of N10 drives
$\Sigma^{(s)}<0$ at every Sobolev level simultaneously.
The full algebraic apparatus ---
antisymmetry, vanishing, factored form, and
homogeneity --- extends to all~$s$ by
Corollary~\ref{cor:universal-s}.
For the TG orbit, the real transfer
formula~\eqref{eq:TG-orbit-transfer}
gives $T^{(s)}(k)$ as a real sum at every~$s$.

The case $s=1$ is of particular significance for the
regularity problem: $\Sigma^{(1)}(u) < 0$ means the
nonlinearity spreads enstrophy angularly, opposing the
concentration that would be needed for vortex-sheet or
vortex-tube singularities.


\subsection{Vortex-stretching decomposition (N13)}
\label{sec:stretching}

The enstrophy equation for 3D Navier--Stokes reads
\begin{equation}\label{eq:enstrophy}
\frac{1}{2}\frac{d}{dt}\|\omega\|_{L^2}^2
= -\nu\|\nabla\omega\|_{L^2}^2
+ \int_{\T^3}\omega\cdot(\omega\cdot\nabla)u\,dx,
\end{equation}
where the vortex-stretching term
$\mathcal{W}=\int\omega\cdot(\omega\cdot\nabla)u\,dx$
is the sole mechanism for potential enstrophy blow-up.
We decompose $\mathcal{W}$ by angular sector.

Over $t\in[5,10]$ at $N=128$, $\nu=0.005$:
\begin{center}
\renewcommand{\arraystretch}{1.2}
\begin{tabular}{@{}lcc@{}}
\toprule
\textbf{Mechanism} & \textbf{corr with $p_j^{(1)}$}
  & \textbf{Effect} \\
\midrule
Full NL transfer $dE_j^{(1)}|_{\rm NL}$
  & $-0.43$ & net spreading \\
Vortex stretching $\mathcal{W}_j$
  & $+0.14$ & slight concentration \\
\bottomrule
\end{tabular}
\end{center}

The full nonlinear transfer produces net spreading
($\corr=-0.43$), while the stretching component alone
produces weak concentration ($\corr=+0.14$).  The
spreading is dominated by cross-sector transfer, which
overwhelms the concentrating tendency of stretching by
a factor of $\approx 3$ in correlation magnitude.
Along the TG orbit, $\cC^{(1)}$ is monotonically
decreasing: spreading wins.


\subsection{Viscosity independence at $s=1$ (N14)}
\label{sec:enst-visc}

We repeat the viscosity-independence test at the
enstrophy level.

\begin{table}[h]
\centering
\caption{Enstrophy-level constitutive relation.
$\Sigma^{(1)}/(\cC^{(1)}-1/J)$ is independent of
$\nu$ to four significant figures.}
\label{tab:enst-visc}
\begin{tabular}{@{}cccc@{}}
\toprule
$\nu$ & $\langle\cC^{(1)}\rangle$
  & $\langle\Sigma^{(1)}/(\cC^{(1)}\!-\!1/J)\rangle$
  & std \\
\midrule
$0.02$  & $0.238$ & $-1.367\times 10^{-2}$
  & $2.36\times 10^{-3}$ \\
$0.01$  & $0.238$ & $-1.369\times 10^{-2}$
  & $2.37\times 10^{-3}$ \\
$0.005$ & $0.238$ & $-1.370\times 10^{-2}$
  & $2.37\times 10^{-3}$ \\
$0.002$ & $0.238$ & $-1.370\times 10^{-2}$
  & $2.37\times 10^{-3}$ \\
\bottomrule
\end{tabular}
\end{table}

The constitutive relation
$\Sigma^{(1)}(u) = -g^{(1)}(u)\,(\cC^{(1)}(u)-1/J)$
holds with $g^{(1)}>0$ independent of $\nu$
(Table~\ref{tab:enst-visc}), confirming that
enstrophy-level spreading is entirely a property of
the nonlinear term.


\subsection{Multi-IC enstrophy spreading (N15)}
\label{sec:enst-multi-ic}

\begin{table}[h]
\centering
\caption{Enstrophy spreading ($s=1$) across initial
conditions at $N=128$, $\nu=0.01$, $J=32$.
Compare with Table~\ref{tab:multi-ic} ($s=1/2$).}
\label{tab:enst-multi-ic}
\begin{tabular}{@{}lcccl@{}}
\toprule
IC & $\Sigma^{(1)}(t\!=\!5)$ & $\corr$
  & $n_{\rm spread}$ & status \\
\midrule
Taylor--Green
  & $-1.93\times 10^{-3}$ & $-0.45$ & 10/10
  & spreads \\
ABC flow
  & $\approx 0$ & $+0.99$ & 2/10
  & degenerate \\
Random (seed 42)
  & $-4.2\times 10^{-11}$ & $-0.17$ & 10/10
  & spreads \\
Random (seed 99)
  & $-8.9\times 10^{-13}$ & $-0.00$ & 10/10
  & spreads \\
Random (seed 7)
  & $-1.6\times 10^{-11}$ & $-0.07$ & 10/10
  & spreads \\
Conc.\ jet
  & $+3.5\times 10^{-11}$ & $-0.53$ & 0/10
  & concentrates \\
\bottomrule
\end{tabular}
\end{table}

At $s=1/2$ (Table~\ref{tab:multi-ic}), the three
random broadband ICs gave mixed results: two seeds
spread and one concentrated.  At $s=1$
(Table~\ref{tab:enst-multi-ic}), \emph{all three
random seeds spread at every snapshot} (30/30 negative
values).  The probability of 30 consecutive negative
values from a distribution symmetric around zero
(Theorem~\ref{thm:gaussian}) is
$2^{-30}\approx 10^{-9}$.  Enstrophy-level spreading
is therefore more universal than $\Hdot{1/2}$ spreading.

The jet remains the only counterexample, with
$\Sigma^{(1)}>0$ at every snapshot, and the ABC flow
remains degenerate.


\subsection{Stretching landscape and the depletion
transition (N16)}
\label{sec:Phi-landscape}

To determine how vortex stretching depends on
angular concentration and phase structure jointly,
we construct synthetic fields by independently
modifying the concentration (sector amplitude
rescaling, parameter $\gamma$) and the phase
coherence (partial phase randomization, fraction $f$)
of a developed NS solution.

At full phase coherence ($f=0$), the mode-level
depletion metric
$\kappa_{\rm align} = -\sum_k w_k
\cos(\arg[\hat u(k)^*\cdot\hat B(k)])$
(enstrophy-weighted) undergoes a sharp sign
transition:
\begin{center}
\begin{tabular}{@{}ccc@{}}
\toprule
$\cC^{(1)}$ & $\kappa_{\rm align}$ & regime \\
\midrule
$0.04$ & $+1.00$ & amplification \\
$0.09$ & $+0.99$ & amplification \\
$0.15$ & $+0.83$ & transition \\
$0.21$ & $-0.56$ & depletion \\
$0.25$ & $-1.00$ & full depletion \\
\bottomrule
\end{tabular}
\end{center}

For nearly isotropic fields, the nonlinear term
\emph{amplifies} all modes: there is no dominant
sector to deplete.  As concentration increases past
$\cC^{(1)}\approx 0.15$, the cross-sector transfer
engages and the $-i$ mechanism begins depleting the
concentrated sectors.  This demonstrates that the
hypothesis ``$\cC^{(1)}>1/J$'' in the spreading
conjecture is structurally necessary.

A simple bound $\Phi\leq h(\cC^{(1)})$ with
$h(1/J)=0$ does not hold: the normalized stretching
$\Phi = |\mathcal{W}|/(\|\omega\|_{L^2}^2
\|\nabla u\|_{L^2})$ depends on the full phase
structure, not just angular concentration.
A multivariate regression of $\log\Phi$ against
$(\log(\cC^{(1)}-1/J),\,\kappa_{\rm imag})$ yields
$R^2=0.16$: the stretching landscape is not
well-described by simple scalar parameters.


\subsection{Spreading robustness under phase
destruction (N17)}
\label{sec:robustness}

\begin{center}
\begin{tabular}{@{}cccc@{}}
\toprule
$f$ & $\Sigma^{(1)}<0$ & $\langle\Sigma^{(1)}\rangle$
  & $\langle\kappa_{\rm align}\rangle$ \\
\midrule
$0.0$ & 11/11 & $-1.74\times 10^{-3}$ & $-0.43$ \\
$0.1$ & 10/11 & $-1.27\times 10^{-3}$ & $-0.36$ \\
$0.3$ & 10/11 & $-8.30\times 10^{-4}$ & $-0.30$ \\
$0.5$ & 10/11 & $-3.04\times 10^{-4}$ & $-0.04$ \\
$0.7$ &  5/11 & $-1.93\times 10^{-4}$ & $-0.04$ \\
$1.0$ &  7/11 & $+9.25\times 10^{-6}$ & $+0.10$ \\
\bottomrule
\end{tabular}
\end{center}

Enstrophy spreading persists at nearly all
concentration levels ($\geq 10/11$) up to 50\%
phase randomization.  The breakdown around $f=0.7$
coincides with $\kappa_{\rm align}$ crossing zero.
Combined with N15 (where random broadband ICs develop
$\Sigma^{(1)}<0$ under NS evolution), this suggests
the NS dynamics drive sufficient phase coherence
for spreading even from incoherent initial data.


\subsection{Surgical concentration and self-regulation (N18)}
\label{sec:surgical}

The preceding tests evolve NS solutions from prescribed
initial conditions and measure spreading along the
resulting trajectories.  A stronger test of the
regularity program is to \emph{artificially concentrate}
a developed NS solution---preserving its phase
structure---and ask whether the dynamics restore
isotropy.

\medskip\noindent\textbf{Construction.}
We evolve the TG vortex at $N=256$, $\nu=0.01$ to
$t=5$, producing a developed turbulent field with
$\cC^{(1)}\approx 0.042$.  We then surgically
increase enstrophy concentration by amplifying the
dominant sector by a factor $\gamma>1$ and suppressing
all others by $\gamma^{-1/2}$, with $\gamma$ chosen
by bisection to reach target concentrations
$\cC^{(1)}\in\{0.20, 0.30, 0.40, 0.50\}$.  The
modified field is Leray-projected to restore
divergence-freeness.  Crucially, this procedure
preserves the Fourier phase structure generated by
NS evolution: only the sector amplitudes change.
Each surgically concentrated field is then evolved
under NS for $t\in[0,5]$ at the same $\nu=0.01$.

\medskip\noindent\textbf{Results.}
At all 44 diagnostic snapshots across all four
concentration levels, $\Sigma^{(1)}<0$ without
exception.  Enstrophy concentration decreases
monotonically from its surgically imposed value:

\begin{table}[h]
\centering
\caption{Surgical concentration: NS relaxation at
$N=256$, $\nu=0.01$, $J=32$.  $\Sigma^{(1)}<0$ at
every snapshot (44/44).}
\label{tab:surgical}
\begin{tabular}{@{}cccccc@{}}
\toprule
$\cC^{(1)}(0)$ & $\cC^{(1)}(5)$
  & peak $|\Sigma^{(1)}|$
  & $|W|$ range
  & $\mathcal{R}$ range
  & $g^{(1)}$ range \\
\midrule
$0.20$ & $0.13$ & $7.95\times 10^{-2}$
  & $11$--$34$ & $1.8$--$2.2$ & $0.40$--$0.56$ \\
$0.30$ & $0.17$ & $1.16\times 10^{-1}$
  & $11$--$30$ & $1.1$--$1.4$ & $0.33$--$0.52$ \\
$0.40$ & $0.21$ & $1.47\times 10^{-1}$
  & $11$--$29$ & $0.8$--$1.0$ & $0.27$--$0.49$ \\
$0.50$ & $0.23$ & $1.74\times 10^{-1}$
  & $12$--$27$ & $0.6$--$0.8$ & $0.22$--$0.48$ \\
\midrule
$0.042$ (base) & $0.040$ & $2.87\times 10^{-3}$
  & $17$--$78$ & $68$--$295$ & $0.08$--$0.31$ \\
\bottomrule
\end{tabular}
\end{table}

Three features of the data are significant:

\medskip\noindent\textbf{(i) Monotone relaxation.}
Enstrophy concentration $\cC^{(1)}(t)$ decreases at
every snapshot in every surgical run.  The NS
nonlinearity does not permit sustained angular
concentration: when concentration is artificially
elevated, the dynamics pull it back toward isotropy
with a rate proportional to the excess
$\cC^{(1)}-1/J$.

\medskip\noindent\textbf{(ii) Spreading scales with
concentration; stretching does not.}
The peak spreading rate $|\Sigma^{(1)}|$ increases
from $0.080$ to $0.174$ as the initial concentration
rises from $0.20$ to $0.50$ (a factor of $2.2\times$).
By contrast, vortex stretching $|W|$ remains in the
range $11$--$34$ across all surgical runs---essentially
the same range as the unmodified base trajectory.
This decoupling is the quantitative content of
self-regulation: the spreading mechanism responds
to concentration, while the stretching mechanism
does not.

\medskip\noindent\textbf{(iii) The ratio
$\mathcal{R}$ decreases with concentration.}
Define the stretching--spreading ratio
\begin{equation}\label{eq:R-ratio}
\mathcal{R}(t) = \frac{|\mathcal{W}(t)|}
{|\Sigma^{(1)}(t)|\cdot\|\omega(t)\|_{L^2}^2},
\end{equation}
which measures how efficiently the nonlinear
spreading functional $\Sigma^{(1)}$ counteracts
vortex stretching $\mathcal{W}$, normalized by
enstrophy.  In the surgical runs,
$\mathcal{R}\in[0.6, 2.2]$ and \emph{decreases}
as concentration increases
(Table~\ref{tab:surgical}).  This is the opposite
of the blow-up scenario: the higher the
concentration, the more effectively spreading
dominates stretching.

The large values $\mathcal{R}\approx 70$--$295$
in the base run are harmless: they occur at
$\cC^{(1)}\approx 0.04 \approx 1/J$, where the
field is already isotropic and $|\Sigma^{(1)}|$ is
small because there is no angular concentration to
spread.  In this regime, neither stretching nor
spreading drives the dynamics: viscous dissipation
dominates.


\subsection{The constitutive law (N19)}
\label{sec:constitutive-surgical}

Sections~\ref{sec:constitutive}
and~\ref{sec:enst-visc} established the constitutive
relation $\Sigma^{(1)} = -g^{(1)}(u)\,
(\cC^{(1)}-1/J)$ along single TG trajectories, where
$g^{(1)}$ is viscosity-independent but varies in
time by a factor of~5.  The earlier regression of
$g^{(1)}$ against scalar parameters
(Section~\ref{sec:Phi-landscape}) gave $R^2=0.16$.
The surgical data resolve this.

\medskip\noindent\textbf{A nearly universal
constant.}
By pooling all snapshots from the four surgical runs
at $t\geq 1.5$ (after the initial transient), the
constitutive coefficient
\begin{equation}\label{eq:g1-value}
g^{(1)} = \frac{-\Sigma^{(1)}}{\cC^{(1)}-1/J}
\end{equation}
stabilizes at
\begin{equation}\label{eq:g1-numerical}
g^{(1)} \approx 0.47 \pm 0.04
\end{equation}
across the full range
$\cC^{(1)}\in[0.13, 0.50]$
(Table~\ref{tab:surgical}).  The coefficient of
variation is under $10\%$.

\begin{remark}[Why the single-trajectory regression
failed]
Along a single TG trajectory, $\cC^{(1)}$ and the
cascade age (measured by $k_*$ or $t$) are perfectly
correlated: both evolve together.  The regression
cannot distinguish their effects.  The surgical
construction decouples them: all runs share the same
cascade age (turbulence developed to $t=5$) but have
different concentrations.  The resulting near-constant
$g^{(1)}\approx 0.47$ shows that the constitutive
coefficient depends on the cascade phase structure
(which is shared), not on $\cC^{(1)}$ itself.
\end{remark}

\begin{remark}[Analytical target]
\label{rem:analytical-target}
The numerical constitutive law
\begin{equation}\label{eq:constitutive-target}
\Sigma^{(1)}(u) \approx -0.47\,
(\cC^{(1)}(u)-1/J)
\end{equation}
suggests a static lower bound
$g^{(1)}(u) \geq c > 0$.  The randomization
tests of N20 below show this bound cannot hold
on any static function class: it requires
the NS dynamics.  Accordingly, the analytical
target is a \emph{dynamical} statement:
the NS flow preserves (or creates)
inter-scale correlations in $V(k)$ that
maintain $g^{(1)}>0$.
\end{remark}

\begin{remark}[Connection to Step~(ii)]
\label{rem:step-ii}
The boundedness of $\mathcal{R}$ in the surgical
runs (Table~\ref{tab:surgical}) is a direct
numerical verification of Step~(ii) in
Remark~\ref{rem:regularity-final}: the
stretching--spreading ratio remains bounded along
NS trajectories, even under adversarial
concentration.  Combined with Step~(i) (the
enstrophy spreading conjecture, supported by
N12--N17 and the 44/44 surgical snapshots), this
closes the numerical side of the regularity program.
The remaining challenge is analytical: proving
that the NS dynamics maintain $g^{(1)}>0$
(see N20).
\end{remark}


\subsection{The dynamical origin of spreading (N20)}
\label{sec:dynamical}

The constitutive law $g^{(1)}\approx 0.47$ was
established for NS solutions with developed cascade
phase structure (N19).  A natural question is whether
this positivity follows from any static property of
the velocity field---its amplitude spectrum, its
purely imaginary Fourier structure, or the
divergence-free constraint---or whether it requires
the specific correlations created by the NS dynamics.

\medskip\noindent\textbf{Randomization tests.}
Starting from the developed TG solution at $t=10$
($N=128$, $\nu=0.005$), we construct four families
of modified fields, each preserving progressively
less structure, and compute $\Sigma^{(1)}$ for
50 independent realizations of each:

\begin{enumerate}[leftmargin=2em]
\item \textbf{Random $V$}: same $|V(k)|$ at every
  mode, random direction $\perp k$, reality
  preserved.  Destroys all inter-mode correlations.
\item \textbf{Isotropized $V$}: same energy spectrum
  $E(|k|)$ per shell, uniform distribution within
  shells.  Destroys angular structure and mode-level
  amplitude variation.
\item \textbf{Permuted $V$}: swap directions between
  modes within each shell.  Preserves amplitudes and
  the set of polarizations per shell, destroys which
  mode has which direction.
\item \textbf{Concentrated random $V$}: random $V$
  with the dominant sector amplified by
  $\gamma\in\{1,2,3,5\}$.  Tests whether
  concentration alone activates spreading in
  non-NS fields.
\end{enumerate}

\begin{table}[h]
\centering
\caption{Randomization tests ($N=128$, $\nu=0.005$,
$J=32$, 50 seeds each).  Only the NS solution
produces robust spreading; all randomizations give
results indistinguishable from the
conjugation-symmetric null
(Theorem~\ref{thm:gaussian}).}
\label{tab:randomization}
\begin{tabular}{@{}lcccc@{}}
\toprule
\textbf{Field} & $\Sigma^{(1)}<0$
  & mean $g^{(1)}$ & std $g^{(1)}$
  & status \\
\midrule
NS solution & 1/1 & $+0.018$ & ---
  & spreads \\
Random $V$ & 24/50 & $-0.002$ & $0.13$
  & null \\
Isotropized $V$ & 25/50 & $+0.021$ & $0.46$
  & null \\
Permuted $V$ & 20/50 & $-0.027$ & $0.13$
  & null \\
\bottomrule
\end{tabular}
\end{table}

\begin{table}[h]
\centering
\caption{Concentrated random $V$: boosting the
dominant sector in non-NS fields (30 seeds each).}
\label{tab:conc-random}
\begin{tabular}{@{}cccc@{}}
\toprule
$\gamma$ & $\Sigma^{(1)}<0$ & mean $g^{(1)}$
  & std $g^{(1)}$ \\
\midrule
$1$ & 15/30 & $+0.004$ & $0.12$ \\
$2$ & 15/30 & $+0.008$ & $0.10$ \\
$3$ & 18/30 & $+0.010$ & $0.08$ \\
$5$ & 19/30 & $+0.010$ & $0.05$ \\
\bottomrule
\end{tabular}
\end{table}

Every randomization produces $\Sigma^{(1)}<0$
at a rate indistinguishable from 50\%
(Tables~\ref{tab:randomization}--\ref{tab:conc-random}),
consistent with the
conjugation-symmetric vanishing theorem
(Theorem~\ref{thm:gaussian}).
Neither the amplitude spectrum, the energy
distribution across shells, nor the purely
imaginary structure ($\hat u = iV$) is sufficient
for spreading.  Concentrated random fields show a
mild trend (15/30 at $\gamma=1$ to 19/30 at
$\gamma=5$), but this is not statistically
significant.

\medskip\noindent\textbf{Shell decomposition.}
To identify \emph{which} inter-scale correlations
drive spreading, we decompose $\Sigma^{(1)}$ by
the spectral shell of the advecting mode.  The
first two shells exhibit massive cancellation:
shell~0 ($|p|\sim 3$, the energy-containing range)
contributes $\Sigma_0 = -1.85\times 10^{-3}$
(spreading), while shell~1 ($|p|\sim 7$) contributes
$\Sigma_1 = +1.93\times 10^{-3}$ (concentrating).
The net $\Sigma^{(1)} = -8.0\times 10^{-5}$ is
the $4\%$ residual of a $96\%$ cancellation.
Shells~2--9 all contribute negatively but are
individually small.

\medskip\noindent\textbf{The inter-scale
anti-correlation.}
The NS cascade creates a distinctive correlation
structure between spectral shells.  Define the
shell correlation
\begin{equation}
C(s_1,s_2) = \frac{1}{n_{s_1}n_{s_2}}
\sum_{\substack{|k|\in s_1\\|p|\in s_2}}
V(k)\cdot V(p),
\end{equation}
where $n_s$ is the number of active modes in
shell~$s$.  For the NS solution, same-shell
correlations are positive ($C(s,s)>0$: modes at
the same scale tend to align) while cross-shell
correlations are \emph{negative}
($C(s_1,s_2)<0$ for $s_1\neq s_2$: modes at
different scales tend to anti-align), with a
ratio of $-263\times$.  For amplitude-matched
random fields, cross-shell correlations are
\emph{positive} (ratio $+144\times$).  The sign
of the cross-shell correlation is opposite in
the NS solution and random fields.

\begin{remark}[Spreading as a dynamical invariant]
\label{rem:dynamical}
The randomization tests establish that
$\Sigma^{(1)}<0$ is not a property of any static
function class accessible by randomizing a
developed NS solution.  It is a property of the
\emph{NS orbit} --- specifically, of the
inter-scale anti-correlation
$C(s_1,s_2)<0$ for $s_1\neq s_2$ that the
dynamics create and maintain.  The regularity
program therefore requires a dynamical statement:
the NS flow preserves negativity of
$\Sigma^{(1)}$ along trajectories with
$\cC^{(1)}>1/J$.  This is a contraction property
of the nonlinear term, not a static inequality on
a function space.
\end{remark}


\subsection{Dynamical creation of spreading (N21--N22)}
\label{sec:creation}

The randomization tests of N20 show that spreading
requires NS-specific correlations.  A stronger
question: does the NS flow \emph{create} these
correlations from random initial data, or must
they be inherited from a special initial condition?

\medskip\noindent\textbf{Construction.}
We take the amplitude spectrum $|V(k)|$ from the
developed TG solution at $t=10$ ($N=128$,
$\nu=0.005$) and construct 30 random initial
conditions: each has $\hat u(k) = iV_{\rm rand}(k)$
with $|V_{\rm rand}(k)| = |V_{\rm NS}(k)|$,
random direction $\perp k$, and the reality
constraint $V(-k)=-V(k)$.  These fields have
zero inter-scale correlation by construction
(N20: $\Sigma^{(1)}<0$ at $\approx 50\%$,
indistinguishable from the conjugation-symmetric
null).
Each is evolved under NS for $t\in[0,10]$ at
$\nu=0.005$, with $\Sigma^{(1)}$ recorded at
$t\in\{0, 0.1, 0.5, 1, 2, 3, 5, 7, 10\}$.

\begin{table}[h]
\centering
\caption{Dynamical creation of spreading: 30 random
ICs evolved under NS ($N=128$, $\nu=0.005$, $J=32$).
The fraction with $\Sigma^{(1)}<0$ rises from
$50\%$ (null) to $100\%$ within one time unit
and persists through $t=10$.}
\label{tab:creation}
\begin{tabular}{@{}ccccc@{}}
\toprule
$t$ & $\Sigma^{(1)}<0$ & mean $\Sigma^{(1)}$
  & mean $g^{(1)}$ \\
\midrule
$0.0$ & $15/30$ ($50\%$)
  & $-1.8\times 10^{-5}$ & $+0.004$ \\
$0.1$ & $18/30$ ($60\%$)
  & $-2.6\times 10^{-4}$ & $+0.060$ \\
$0.5$ & $29/30$ ($97\%$)
  & $-9.1\times 10^{-4}$ & $+0.213$ \\
$1.0$ & $30/30$ ($100\%$)
  & $-1.3\times 10^{-3}$ & $+0.296$ \\
$2.0$ & $29/30$ ($97\%$)
  & $-1.3\times 10^{-3}$ & $+0.302$ \\
$3.0$ & $28/30$ ($93\%$)
  & $-1.2\times 10^{-3}$ & $+0.262$ \\
$5.0$ & $27/30$ ($90\%$)
  & $-1.1\times 10^{-3}$ & $+0.205$ \\
$7.0$ & $27/30$ ($90\%$)
  & $-1.0\times 10^{-3}$ & $+0.145$ \\
$10.0$ & $29/30$ ($97\%$)
  & $-1.3\times 10^{-3}$ & $+0.127$ \\
\bottomrule
\end{tabular}
\end{table}

The transition is rapid: by $t=0.5$, $97\%$ of
seeds spread; by $t=1.0$, all 30 spread
(Table~\ref{tab:creation}).  The constitutive
coefficient $g^{(1)}$ rises from $+0.004$ (noise)
to $+0.30$ at $t=1$ --- a value comparable to the
surgical experiments of N19 ($g^{(1)}\approx 0.47$
for developed TG turbulence).  At late times
($t=10$), $g^{(1)}$ relaxes to $+0.13$ as the
field approaches isotropy
($\cC^{(1)}\to 1/J$), but spreading
persists: $\Sigma^{(1)}<0$ at $97\%$ of seeds
even at $t=10$.

\medskip\noindent\textbf{Contraction asymmetry (N22).}
Exactly 15 of 30 seeds start with $\Sigma^{(1)}>0$
and 15 start with $\Sigma^{(1)}<0$, consistent with
the conjugation-symmetric null (Theorem~\ref{thm:gaussian}).
Conditioning on initial sign reveals a sharp asymmetry:
\begin{itemize}[leftmargin=2em]
\item \textbf{Positive $\to$ negative:}
  All 15 seeds starting with $\Sigma^{(1)}>0$
  flip to $\Sigma^{(1)}<0$ by $t=1$ (15/15).
  By $t=0.5$, 14/15 have already flipped.  The
  positive basin is unstable.
\item \textbf{Negative persistence:}
  Of 15 seeds starting with $\Sigma^{(1)}<0$,
  all 15 persist through $t=1$.  Even at $t=7$,
  12/15 remain negative (80\%).  The negative
  basin is nearly absorbing.
\end{itemize}
The contraction is not symmetric: the NS
nonlinearity actively pulls $\Sigma^{(1)}>0$
states toward $\Sigma^{(1)}<0$ on the $O(1)$
timescale, while $\Sigma^{(1)}<0$ states resist
displacement.  The rare positive excursions
at late times ($\leq 3/30$ at any snapshot) are
brief transients, not persistent sign changes.

\subsection{High-resolution creation (N23)}
\label{sec:creation-highres}

To rule out resolution artifacts, the creation
experiment is repeated at $256^3$ with 10 random
initial conditions ($\nu=0.005$, $J=32$),
evolved to $t=5$.

\begin{table}[h]
\centering
\caption{High-resolution creation: 10 random ICs
at $256^3$ ($\nu=0.005$, $J=32$).  The
$50\%\to 100\%$ transition replicates at full
resolution.}
\label{tab:creation-256}
\begin{tabular}{@{}ccccc@{}}
\toprule
$t$ & $\Sigma^{(1)}<0$ & mean $\Sigma^{(1)}$
  & mean $g^{(1)}$ \\
\midrule
$0.0$ & $3/10$ ($30\%$)
  & $+2.0\times 10^{-4}$ & $-0.045$ \\
$0.1$ & $5/10$ ($50\%$)
  & $-9.5\times 10^{-5}$ & $+0.024$ \\
$0.5$ & $9/10$ ($90\%$)
  & $-9.2\times 10^{-4}$ & $+0.222$ \\
$1.0$ & $10/10$ ($100\%$)
  & $-1.6\times 10^{-3}$ & $+0.376$ \\
$2.0$ & $10/10$ ($100\%$)
  & $-1.7\times 10^{-3}$ & $+0.430$ \\
$3.0$ & $9/10$ ($90\%$)
  & $-1.2\times 10^{-3}$ & $+0.316$ \\
$5.0$ & $9/10$ ($90\%$)
  & $-8.8\times 10^{-4}$ & $+0.206$ \\
\bottomrule
\end{tabular}
\end{table}

The $256^3$ run (Table~\ref{tab:creation-256})
replicates the $128^3$ creation trajectory:
$100\%$ negative at $t=1$, stabilizing at $90\%$
at $t=5$.  The mean $g^{(1)}$ peaks at $+0.43$
at $t=2$, somewhat higher than the $128^3$ value of
$+0.30$.  The creation phenomenon is not a
resolution artifact.

\subsection{Viscosity independence of creation (N24)}
\label{sec:creation-viscosity}

The constitutive law's viscosity independence (N14)
suggests that spreading creation should also be
$\nu$-independent.  To test this, 10 random
initial conditions are evolved at each of
$\nu\in\{0.002,0.005,0.01,0.02\}$ ($N=128$,
$J=32$, $t\in[0,5]$), using the respective
developed TG templates as amplitude sources.

\begin{table}[h]
\centering
\caption{Viscosity independence of creation: fraction
with $\Sigma^{(1)}<0$ at four viscosities
($N=128$, $J=32$, 10 seeds each).  All four reach
$\geq 90\%$ spreading by $t=1$.}
\label{tab:creation-nu}
\begin{tabular}{@{}ccccc@{}}
\toprule
$t$ & $\nu=0.02$ & $\nu=0.01$ & $\nu=0.005$
  & $\nu=0.002$ \\
\midrule
$0.0$ & $50\%$ & $40\%$ & $60\%$ & $40\%$ \\
$0.5$ & $80\%$ & $90\%$ & $100\%$ & $100\%$ \\
$1.0$ & $90\%$ & $100\%$ & $100\%$ & $80\%$ \\
$2.0$ & $100\%$ & $100\%$ & $100\%$ & $100\%$ \\
$5.0$ & $100\%$ & $100\%$ & $90\%$ & $90\%$ \\
\bottomrule
\end{tabular}
\end{table}

The creation timescale is $O(1)$ at all four
viscosities (Table~\ref{tab:creation-nu}): by $t=1$,
all reach $\geq 80\%$ negative, and by $t=2$, all
reach $100\%$.  Since the viscous timescale
varies as $1/\nu$ while the nonlinear turnover
time is $O(1)$, the $\nu$-independent creation
timescale confirms that spreading creation is a
purely nonlinear phenomenon --- the same mechanism
responsible for the constitutive law (N14).

The magnitudes $|\Sigma^{(1)}|$ do vary with
$\nu$ (larger at higher $\nu$), reflecting
differences in the energy and enstrophy levels
of the respective TG templates.  But the
\emph{sign} behavior and \emph{timescale} are
$\nu$-independent.  The $\nu=0.002$ case (highest
Reynolds number) shows particularly dramatic
transients: $g^{(1)}$ spikes to $\approx 1.8$ at
$t=0.5$ before relaxing to $\approx 0.24$ at $t=5$,
indicating violent nonlinear redistribution during
the initial cascade development.

\begin{remark}[The NS attractor spreads and contracts]
\label{rem:attractor}
The creation and contraction tests (N21--N24)
establish that enstrophy spreading is not a property
of special initial conditions: it is a property
of the \emph{NS attractor}.  Random divergence-free
fields, evolved under NS, develop the inter-scale
anti-correlation (N20) within approximately one
nonlinear turnover time and maintain
$\Sigma^{(1)}<0$ thereafter.  The contraction
asymmetry (N22) reveals that the basin
$\{\Sigma^{(1)}<0\}$ is nearly absorbing:
the NS dynamics drive $\Sigma^{(1)}>0$ states
toward the spreading regime on the $O(1)$
timescale, while states already in the spreading
regime resist displacement.  The resolution
independence (N23) and viscosity independence (N24)
confirm that this contraction is a robust,
purely nonlinear property of the NS dynamics.
The regularity
conjecture (Conjecture~\ref{conj:enstrophy})
is therefore a statement about the attractor:
the NS dynamics contract \emph{all} generic
initial data toward the spreading basin.
\end{remark}


\subsection{The affine ratio and its dynamical contraction
(N29--N33)}
\label{sec:affine-experiments}

The exact identity~\eqref{eq:Sigma-via-c} reduces the
spreading question to the sign of $c^{(1)}-1$.  The
following experiments characterize $c^{(1)}$ along
the TG trajectory, under surgical concentration,
for random fields, and under NS evolution.

\medskip\noindent\textbf{N29 (Degree-0 verification).}
Rescaling the developed TG solution at $t=10$ by
$\lambda\in\{0.25,0.5,1,2,5\}$ gives $c^{(1)}=0.617$
at every $\lambda$ to all displayed digits
(14 significant figures), while $g^{(1)}$ scales as
$\lambda$ (from $0.079$ at $\lambda=0.25$ to $1.58$
at $\lambda=5$).
This confirms Proposition~\ref{prop:c-degree-zero}
numerically and establishes that $c^{(1)}$ is the
natural dimensionless diagnostic for spreading,
while the homogeneity obstruction
(Proposition~\ref{prop:homogeneity}) resides
entirely in the degree-1 prefactor $T_{\rm tot}/E$.

\medskip\noindent\textbf{N30 (Affine ratio along TG
trajectory).}
$c^{(1)}$ along the TG orbit:
\begin{center}
\renewcommand{\arraystretch}{1.2}
\begin{tabular}{@{}cccccc@{}}
\toprule
$t$ & $c^{(1)}$ & $b/T_{\rm tot}$
  & $\Sigma^{(1)}$ & $g^{(1)}$ & $\cC^{(1)}$ \\
\midrule
$0.5$ & $1.04$ & $+0.064$ & $-3.5\!\times\!10^{-2}$ & $+0.42$ & $0.116$ \\
$1.0$ & $0.81$ & $+0.057$ & $-4.7\!\times\!10^{-2}$ & $+0.75$ & $0.094$ \\
$2.0$ & $0.44$ & $+0.045$ & $-2.8\!\times\!10^{-2}$ & $+1.11$ & $0.057$ \\
$3.0$ & $0.18$ & $+0.026$ & $-8.7\!\times\!10^{-3}$ & $+0.78$ & $0.042$ \\
$5.0$ & $0.69$ & $+0.010$ & $-5.0\!\times\!10^{-3}$ & $+0.45$ & $0.042$ \\
$7.0$ & $0.57$ & $+0.013$ & $-3.9\!\times\!10^{-3}$ & $+0.43$ & $0.040$ \\
$10.0$ & $0.62$ & $+0.012$ & $-1.8\!\times\!10^{-3}$ & $+0.32$ & $0.037$ \\
\bottomrule
\end{tabular}
\end{center}
At $t=0.5$, $c^{(1)}\approx 1.04$ ($\Sigma^{(1)}<0$ already
holds because $\sum\varepsilon_j T_j<0$ at early times:
concentrated sectors actively lose transfer while
seeding the cascade).  By $t=1$, $c^{(1)}$ drops below~$1$
and remains there for all subsequent times.  The
normalized intercept $b/T_{\rm tot}$ is strictly positive
at every snapshot.

\medskip\noindent\textbf{N31 (Surgical concentration:
$c^{(1)}$ vs.\ $\cC^{(1)}$).}
Boosting the 8 most concentrated sectors by $\gamma$ at
$t=10$:
\begin{center}
\renewcommand{\arraystretch}{1.2}
\begin{tabular}{@{}ccccc@{}}
\toprule
$\gamma$ & $\cC^{(1)}$ & $c^{(1)}$
  & $b/T_{\rm tot}$ & $\Sigma^{(1)}$ \\
\midrule
$1$ & $0.037$ & $0.62$ & $+0.012$ & $-1.8\!\times\!10^{-3}$ \\
$1.5$ & $0.055$ & $0.64$ & $+0.011$ & $-8.7\!\times\!10^{-3}$ \\
$2$ & $0.073$ & $0.69$ & $+0.010$ & $-1.5\!\times\!10^{-2}$ \\
$3$ & $0.095$ & $0.76$ & $+0.007$ & $-2.2\!\times\!10^{-2}$ \\
$5$ & $0.113$ & $0.85$ & $+0.005$ & $-2.5\!\times\!10^{-2}$ \\
$10$ & $0.122$ & $0.92$ & $+0.003$ & $-2.4\!\times\!10^{-2}$ \\
$20$ & $0.124$ & $0.96$ & $+0.001$ & $-2.2\!\times\!10^{-2}$ \\
\bottomrule
\end{tabular}
\end{center}
$c^{(1)}$ increases monotonically with concentration but
remains strictly below~$1$ at all tested levels
(up to $\cC^{(1)}=0.124$).  The intercept $b/T_{\rm tot}$
shrinks but stays positive.  However, single-sector
boosts ($\gamma=50$--$500$ on one sector, pushing $p_{\max}$
above $0.99$) give $c^{(1)}>1$: the mechanism fails at
extreme concentration.  The NS dynamics correct this:
evolving the $\gamma=100$ single-sector state under NS,
$c^{(1)}$ drops from $1.035$ to $0.25$ within $t=0.25$
(one quarter turnover time).

\medskip\noindent\textbf{N32 (Kinematic test:
$c^{(1)}<1$ requires NS dynamics).}
Random divergence-free fields $\hat u = iV$ with
$T_{\rm tot}>0$ (50~seeds): $c^{(1)}<1$ in only $3/26$
cases (12\%), with $c^{(1)}$ ranging up to 294.
Random fields with sector~3 boosted $3\times$ and
$T_{\rm tot}>0$: $c^{(1)}<1$ in $3/22$ (14\%).
By contrast, the NS solution at $t=10$ has
$c^{(1)}=0.617$.  The bound $c^{(1)}<1$ is not a
kinematic consequence of incompressibility or
the trilinear structure; it requires the
correlations created by the NS cascade.

\medskip\noindent\textbf{N33 (Absorbing basin for
$c^{(1)}<1$).}
Five seeds starting with $c^{(1)}<1$ (perturbed NS
snapshots) and five seeds starting with $c^{(1)}>1$
(random $V$ with $T_{\rm tot}>0$, $c^{(1)}$ up to
$38$), all evolved under NS for $t=1.5$:
\begin{itemize}[leftmargin=2em]
\item All 5 low-$c$ seeds remain $c^{(1)}<1$
  through $t=1.5$ ($c^{(1)}\in[0.22, 0.67]$).
\item All 5 high-$c$ seeds contract to $c^{(1)}<1$
  by $t=0.5$ ($c^{(1)}\in[0.10, 0.83]$ at $t=0.5$,
  $c^{(1)}\in[0.03, 0.63]$ at $t=1.5$).
\end{itemize}
The basin $\{c^{(1)}<1\}$ is absorbing under NS
dynamics: random fields with $c^{(1)}\gg 1$ are
pulled below~$1$ within one half turnover time,
while NS-evolved fields with $c^{(1)}<1$ persist.
Combined with the degree-0 property
(Proposition~\ref{prop:c-degree-zero}), this
identifies $c^{(1)}<1$ as a \emph{dimensionless
dynamical invariant} of the NS attractor.

\begin{remark}[The affine ratio vs.\ the
constitutive coefficient]
\label{rem:c-vs-g}
The constitutive law $\Sigma^{(1)}\approx
-g^{(1)}(\cC^{(1)}-1/J)$ and the affine identity
$\Sigma^{(1)} = 2(c^{(1)}-1)\,T_{\rm tot}\,
\sum\varepsilon_j^2/E$ are exact restatements of
the same algebraic
identity~\eqref{eq:Sigma-linear}.  Their content
differs:
\begin{itemize}[leftmargin=2em]
\item $g^{(1)}$ is degree~1 (scales with $\|u\|$);
  $c^{(1)}$ is degree~0 (dimensionless).
\item $g^{(1)}$ absorbs all amplitude information;
  $c^{(1)}$ isolates the geometric spreading
  mechanism from the stretching magnitude.
\item The regularity obstruction
  (Remark~\ref{rem:obstruction}) resides in the
  degree-1 prefactor $T_{\rm tot}/E$, not in
  $c^{(1)}$.
\end{itemize}
The natural splitting is: (i)~prove $c^{(1)}<1$
(dimensionless, no homogeneity obstruction);
(ii)~bound $T_{\rm tot}/E$ (classical stretching
control).  The contribution of the angular
spreading program is~(i).
\end{remark}


\subsection{The spectral centroid inequality and
viscous balance (N34--N35)}
\label{sec:centroid-inequality}

The identity $c^{(1)} = c_{\rm cen}$ in spectral
steady state reduces the spreading question to a
statement about the angular distribution of enstrophy
across spectral shells.  We prove $c_{\rm cen}<1$
for smooth NS solutions with anisotropic initial
data, using only the NS equation, smoothness, and
the Gronwall inequality.

\subsubsection{Algebraic reduction}

\begin{proposition}[Centroid characterization]
\label{prop:c-centroid}
In exact spectral steady state
($T_j = 2\nu D_j$ for all~$j$, where
$D_j = \sum_{k\in\Omega_j}|k|^4|\hat u(k)|^2$),
the affine ratio equals the centroid ratio:
\begin{equation}\label{eq:c-centroid-exact}
c^{(1)} = c_{\rm cen} :=
\frac{\sum_j\varepsilon_j\,
\langle|k|^2\rangle_j\,p_j}
{\langle|k|^2\rangle\,
\sum_j\varepsilon_j^2},
\end{equation}
where
$\langle|k|^2\rangle_j := D_j/E_j^{(1)}$
is the enstrophy-weighted spectral centroid of
sector~$j$ and
$\langle|k|^2\rangle := D/E^{(1)}$ is the global
centroid.
\end{proposition}

\begin{proof}
$T_j = 2\nu D_j$ and $T_{\rm tot} = 2\nu D$ give
$c^{(1)} = \sum\varepsilon_j D_j/(D\sum\varepsilon_j^2)$.
Writing $D_j = \langle|k|^2\rangle_j\,p_j\,E^{(1)}$
and $D = \langle|k|^2\rangle\,E^{(1)}$
yields~\eqref{eq:c-centroid-exact}.
\end{proof}

\begin{proposition}[Two-shell identity]
\label{prop:two-shell}
Partition the spectral shells into two groups:
$L$ (low, centroid $\kappa_L$) and $H$ (high,
centroid $\kappa_H>\kappa_L$), with combined
weights $w_L,w_H=1-w_L$ and cross-correlations
$\Phi_G = \sum_j\varepsilon_j\varepsilon_j^{(G)}$.
Then:
\begin{equation}\label{eq:two-shell-id}
(\mathrm{Den}-\mathrm{Num}\text{ of }c_{\rm cen})
= w_L w_H(\kappa_H-\kappa_L)(\Phi_L-\Phi_H).
\end{equation}
Therefore $c_{\rm cen}<1$ if and only if $\Phi_L>\Phi_H$.
\end{proposition}

\begin{proof}
Let $N = w_L\kappa_L\Phi_L+w_H\kappa_H\Phi_H$
and $D = (w_L\kappa_L+w_H\kappa_H)(w_L\Phi_L+w_H\Phi_H)$.
Expand $D-N$ using $w_L+w_H=1$:
$D-N = w_Lw_H(\kappa_H\Phi_L+\kappa_L\Phi_H
-\kappa_L\Phi_L-\kappa_H\Phi_H)
= w_Lw_H(\kappa_H-\kappa_L)(\Phi_L-\Phi_H)$.
\end{proof}

\begin{remark}[Interpretation]
The two-shell condition is transparent: the
low-wavenumber group must have more angular
anisotropy than the high-wavenumber group.
This is exactly spectral separation.
\end{remark}

\subsubsection{Chebyshev sufficient condition}

\begin{theorem}[Chebyshev sufficient condition]
\label{thm:centroid}
If the shell-level cross-correlation
$\Phi(r) := \sum_j\varepsilon_j\varepsilon_j^{(r)}$
is non-increasing in $r$ (condition RTI), then
$c_{\rm cen}\leq 1$.
\end{theorem}

\begin{proof}
The shell decomposition gives
$c_{\rm cen} = \sum_r\kappa_r w_r\Phi(r)\,/\,
[(\sum_r\kappa_r w_r)(\sum_r w_r\Phi(r))]$.
With $(\kappa_r)$ non-decreasing and $(\Phi(r))$
non-increasing, the Chebyshev sum inequality
(Hardy--Littlewood--P\'olya, Theorem~43) gives
$c_{\rm cen}\leq 1$.
\end{proof}

\begin{remark}[RTI is sufficient but not necessary]
\label{rem:RTI-not-necessary}
The NS field at $t=10$ \emph{violates} RTI:
$\Phi(r)$ decreases through shell~2 but increases
in shells~3--5 (N35 below).  Yet $c_{\rm cen}=0.691<1$.
The Chebyshev theorem provides a clean sufficient
condition but is too strong.  The two-shell identity
(Proposition~\ref{prop:two-shell}) gives the correct
necessary-and-sufficient condition: $\Phi_L>\Phi_H$.
The following dynamical proof establishes this directly.
\end{remark}

\subsubsection{Dynamical proof via viscous balance}

The NS equation provides a mode-level enstrophy
balance.  For the sector-shell enstrophy
$E_j^{(r)} = \sum_{k\in\Omega_j\cap S_r}
|k|^2|\hat u(k)|^2$:
\begin{equation}\label{eq:Ej-evolution}
\frac{d}{dt}E_j^{(r)} = F_j^{(r)}
- 2\nu D_j^{(r)},
\end{equation}
where $F_j^{(r)}$ is the nonlinear forcing rate
and $D_j^{(r)}\approx\kappa_r E_j^{(r)}$ for
narrow shells.

\begin{proposition}[Gronwall bound on sector excess]
\label{prop:gronwall}
The sector excess
$\tilde E_j^{(r)} = E_j^{(r)}-E(r)/J$
satisfies:
\begin{equation}\label{eq:gronwall}
|\tilde E_j^{(r)}(t)| \leq
|\tilde E_j^{(r)}(0)|\,e^{-2\nu\kappa_r t}
+ \frac{M}{2\nu\kappa_r},
\end{equation}
where $M = \sup_{[0,T]}|\tilde F_j^{(r)}|$ is
bounded by $C\|u\|_{H^2}^3$ uniformly in $r$
and~$j$ (by standard bilinear estimates for smooth
fields).
\end{proposition}

\begin{proof}
$d|\tilde E_j|/dt \leq |\tilde F_j|
-2\nu\kappa_r|\tilde E_j|$ (from~\eqref{eq:Ej-evolution}
with $D_j^{(r)}\approx\kappa_r E_j^{(r)}$).
Gronwall with integrating factor $e^{2\nu\kappa_r t}$.
\end{proof}

\begin{theorem}[Spectral separation from viscous
balance]
\label{thm:viscous-sep}
Let $u(t)$ be a smooth NS solution on $\T^3$ with
$\nu>0$ and initial data satisfying
$\mathcal{V}^{(L)}(0) := \sum_j
(\varepsilon_j^{(L)}(0))^2 > 0$ (angular anisotropy
at low wavenumber).  Then there exist
$\kappa_H^*$ and $T^*>0$ (depending on $u_0,\nu,J$)
such that for $\kappa_H>\kappa_H^*$ and
$t\in(0,T^*]$:
\begin{equation}
c_{\rm cen}(u(t)) < 1.
\end{equation}
The time $T^*$ is of order
$1/(\nu\kappa_L)\cdot\ln(\kappa_H/\kappa_L)$
(the viscous timescale of the low-shell group).
No hypothesis on the cross-correlation $X =
\sum_j\varepsilon_j^{(L)}\varepsilon_j^{(H)}$
is needed.
\end{theorem}

\begin{proof}
\textit{Step 1 (Low-shell anisotropy persists).}
The Gronwall bound gives $|\varepsilon_j^{(L)}(t)|
\geq |\varepsilon_j^{(L)}(0)|e^{-2\nu\kappa_L t}
- M/(2\nu\kappa_L E_L)$.  For
$t\leq T_1 = 1/(4\nu\kappa_L)$:
$e^{-2\nu\kappa_L t}\geq e^{-1/2}\approx 0.61$,
so $\Phi_L(t)\geq c_0\Phi_L(0)>0$.

\textit{Step 2 (High-shell angular anisotropy is
viscously suppressed).}
For $t\geq 1/(2\nu\kappa_H)$, the transient decays
and the Gronwall bound gives
$|\varepsilon_j^{(H)}(t)|\leq
C\|u\|_{H^2}^3/(2\nu\kappa_H\,E_H)$ for each~$j$.
By Cauchy--Schwarz:
\begin{equation}\label{eq:PhiH-bound}
|\Phi_H(t)| = \Bigl|\sum_j\varepsilon_j\,
\varepsilon_j^{(H)}\Bigr|
\leq \sqrt{\textstyle\sum\varepsilon_j^2}\cdot
\frac{C\sqrt{J}\,\|u\|_{H^2}^3}{2\nu\kappa_H\,E_H}.
\end{equation}
This vanishes as $\kappa_H\to\infty$.

\textit{Step 3 ($\Phi_L>\Phi_H$ directly).}
From Steps 1--2: $\Phi_L\geq c_0\Phi_L(0)>0$
while $|\Phi_H|\leq C/(\nu\kappa_H)\to 0$.
For $\kappa_H>\kappa_H^* :=
C\sqrt{J}\|u\|_{H^2}^3\sqrt{\sum\varepsilon_j^2}
/(2\nu c_0\Phi_L(0)\,E_H)$:
$\Phi_L>|\Phi_H|\geq\Phi_H$.

By Proposition~\ref{prop:two-shell}: $c_{\rm cen}<1$.

The time window $[1/(2\nu\kappa_H),\,T^*]$
requires $T^*\geq 1/(2\nu\kappa_H)$, which holds
since $T^*\sim 1/(\nu\kappa_L)\gg 1/(\nu\kappa_H)$.
\end{proof}

\begin{remark}[The viscous mechanism]
\label{rem:viscous-mechanism}
The proof has a simple physical content: the
viscous term $-\nu|k|^2\hat u(k)$ damps the
angular anisotropy at shell~$r$ at rate
$2\nu\kappa_r$.  High-$k$ anisotropy equilibrates
on timescale $1/(2\nu\kappa_H)$ (fast), while
low-$k$ anisotropy decays on timescale
$1/(2\nu\kappa_L)$ (slow).  The ratio of
timescales is $\kappa_H/\kappa_L\gg 1$, creating
a window in which the field is anisotropic at
low~$k$ but isotropic at high~$k$.

The nonlinear forcing $B(u,u)$ can be (and is,
per experiment N35b) highly anisotropic at high~$k$.
But the viscous damping kills this anisotropy
$\kappa_H/\kappa_L$ times faster than the low-$k$
anisotropy.  The result: the \emph{field} is
spectrally separated even when the \emph{forcing}
is not.
\end{remark}

\subsubsection{Numerical verification (N34--N35)}

\medskip\noindent\textbf{N34 (Spectral centroid
verification).}
At $t=10$ ($128^3$, $\nu=0.005$):
$c^{(1)} = 0.617$ (direct), $c_{\rm cen} = 0.691$
(from centroids), steady-state departure
$\eta_{\rm mean}=0.29$.  Concentrated sectors
($\varepsilon_j>0$, the 8 TG sectors) have mean
spectral centroid $\langle|k|^2\rangle = 48.3$
(below the global mean $58.3$), while depleted sectors
have $60.8$ (above).

\medskip\noindent\textbf{N34b (Amplitude spectrum
dependence).}
Twenty random divergence-free fields with the
NS-evolved amplitude spectrum: $c_{\rm cen}<1$ for
20/20.  However, random fields with flat spectrum
give $c_{\rm cen}=1.39>1$; with K41 spectrum,
$c_{\rm cen}=1.59>1$.  The centroid inequality
requires the steep NS spectral decay.
Theorem~\ref{thm:viscous-sep} explains why: the
viscous damping that creates spectral separation
also steepens the spectrum, concentrating enstrophy
at moderate wavenumbers.

\medskip\noindent\textbf{N35 (RTI violation and
spectral separation).}
The shell-by-shell $\Phi(r)$ for the NS field:
\begin{center}
\renewcommand{\arraystretch}{1.2}
\begin{tabular}{@{}ccccc@{}}
\toprule
Shell $r$ & $\kappa_r$ & $w_r$ & $\Phi(r)$
  & $\Delta\Phi$ \\
\midrule
$0$ & $3.0$ & $0.068$ & $+2.2\!\times\!10^{-2}$ & --- \\
$1$ & $10.8$ & $0.140$ & $+1.5\!\times\!10^{-2}$
  & $-7.6\!\times\!10^{-3}$ \\
$2$ & $35.4$ & $0.537$ & $+2.4\!\times\!10^{-3}$
  & $-1.2\!\times\!10^{-2}$ \\
$3$ & $117.9$ & $0.227$ & $+3.4\!\times\!10^{-3}$
  & $\mathbf{+1.0\!\times\!10^{-3}}$ \\
$4$ & $374.8$ & $0.028$ & $+6.3\!\times\!10^{-3}$
  & $\mathbf{+2.8\!\times\!10^{-3}}$ \\
$5$ & $1238.6$ & $0.000$ & $+8.7\!\times\!10^{-3}$
  & $\mathbf{+2.5\!\times\!10^{-3}}$ \\
\bottomrule
\end{tabular}
\end{center}
(RTI) is violated (shells 3--5 increase), yet
$c_{\rm cen}=0.691<1$.  The two-shell analysis
($L=\{0,1\}$, $H=\{2,\ldots,5\}$) gives
$\Phi_L=0.017>\Phi_H=0.003$, confirming the
two-shell condition.

\medskip\noindent\textbf{N35b (Forcing angular
anisotropy).}
The nonlinear forcing $B(u,u)$ at $t=10$ is
\emph{more} anisotropic at high $k$: the forcing
angular variance $\mathcal{V}_B$ reaches $328\times$
the geometric baseline at shell~4 and $1596\times$ at
shell~5.  Despite this, the \emph{field} angular
variance $\mathcal{V}^{(H)}=0.0018$ is small,
because viscous damping at rate $2\nu\kappa_H$
suppresses the forcing anisotropy before it
accumulates.  This confirms the mechanism of
Theorem~\ref{thm:viscous-sep}: differential
damping, not forcing isotropy, creates spectral
separation.

\begin{corollary}[Conditional spreading from
viscous balance]
\label{cor:conditional-spreading}
For a smooth NS solution with anisotropic initial
data ($\mathcal{V}^{(L)}(0)>0$) in
$\eta$-approximate spectral steady state with
$\eta < 1-c_{\rm cen}$: $c^{(1)}<1$, hence
$\Sigma^{(1)}<0$ when $T_{\rm tot}>0$.
\end{corollary}

\begin{proof}
Theorem~\ref{thm:viscous-sep} gives
$c_{\rm cen}<1$.  The steady-state perturbation
$|c^{(1)}-c_{\rm cen}|\leq C\eta$ gives
$c^{(1)}<1$ when $\eta<(1-c_{\rm cen})/C$.
\end{proof}

\begin{remark}[Status of the proof]
\label{rem:centroid-status}
\textbf{Proved unconditionally:}
Theorem~\ref{thm:viscous-sep} ($c_{\rm cen}<1$
for smooth NS solutions with anisotropic initial
data, for $t\leq T^*=O(1/(\nu\kappa_L))$).
The proof uses only: the NS equation, smoothness,
the Gronwall inequality, the two-shell algebraic
identity, and the Cauchy--Schwarz inequality.
No hypothesis on the cross-correlation $X$, the
spectral slope, or condition (RTI) is needed.

\textbf{Conditional:}
the link $c^{(1)}\approx c_{\rm cen}$ (approximate
spectral steady state, $\eta\approx 0.29$ at $t=10$).
The following subsection eliminates this condition.
\end{remark}

\subsection{Concentration dynamics and the
spreading bound}
\label{sec:concentration-dynamics}

The identity $c_{\rm cen}<1$ (proved) can be
combined with the concentration evolution equation
to establish spreading results \emph{without}
assuming spectral steady state.

\begin{theorem}[Concentration--spreading identity]
\label{thm:conc-spread}
For a smooth NS solution on $\T^3$ with $\nu>0$,
define $\sigma(t) = \cC^{(1)}(t)-1/J
= \sum_j\varepsilon_j^2$.  Then:
\begin{equation}\label{eq:conc-spread}
\Sigma^{(1)} = \frac{d\sigma}{dt}
- \frac{4\nu\sigma D\,(1-c_{\rm cen})}{E^{(1)}},
\end{equation}
where $D = \sum_j D_j$ is the total palinstrophy.
\end{theorem}

\begin{proof}
From $\varepsilon_j = p_j - 1/J$ and
$p_j = E_j^{(1)}/E^{(1)}$:
$d\sigma/dt = 2\sum\varepsilon_j\,d\varepsilon_j/dt
= (2/E)\sum\varepsilon_j(T_j-2\nu D_j
- p_j(T_{\rm tot}-2\nu D))$.
Using $\sum\varepsilon_j p_j = \sigma$,
$\sum\varepsilon_j T_j = c\,T_{\rm tot}\sigma$,
$\sum\varepsilon_j D_j = c_{\rm cen}\,D\,\sigma$:
$d\sigma/dt = (2\sigma/E)[(c-1)T_{\rm tot}
+ 2\nu(1-c_{\rm cen})D]$.
Since $(E/2)\Sigma^{(1)} = (c-1)T_{\rm tot}\sigma$:
$d\sigma/dt = \Sigma^{(1)}
+ 4\nu\sigma D(1-c_{\rm cen})/E$.
\end{proof}

\begin{corollary}[Spreading bound]
\label{cor:spread-bound}
If $c_{\rm cen}(t)<1$, then
$\Sigma^{(1)}(t) < d\sigma/dt$.
In particular: if $d\sigma/dt\leq 0$ (angular
concentration is not increasing), then
$\Sigma^{(1)}<0$.
\end{corollary}

\begin{proof}
The term $4\nu\sigma D(1-c_{\rm cen})/E$ is
positive when $c_{\rm cen}<1$, $\sigma>0$, $D>0$.
\end{proof}

\begin{theorem}[Spreading at concentration
equilibrium]
\label{thm:equil}
If $c_{\rm cen}(t)<1$, $T_{\rm tot}(t)>0$,
$\sigma(t)>0$, and $d\sigma/dt = 0$, then:
\begin{equation}\label{eq:c-at-equil}
c^{(1)}(t) = 1 - \frac{2\nu D\,(1-c_{\rm cen})}
{T_{\rm tot}} < 1.
\end{equation}
\end{theorem}

\begin{proof}
At $d\sigma/dt = 0$:
$(c-1)T_{\rm tot} = -2\nu(1-c_{\rm cen})D$
(from the concentration dynamics).
Since $c_{\rm cen}<1$, $D>0$, $T_{\rm tot}>0$:
$c-1<0$.
\end{proof}

\begin{theorem}[Time-averaged spreading]
\label{thm:time-avg}
For a smooth NS solution with $c_{\rm cen}(t)<1$
on $[0,T]$:
\begin{equation}\label{eq:time-avg}
\frac{1}{T}\int_0^T\!\Sigma^{(1)}\,dt
\leq \frac{\sigma(T)-\sigma(0)}{T}
- \frac{4\nu}{T}\int_0^T\!
\frac{\sigma D(1-c_{\rm cen})}{E}\,dt.
\end{equation}
Since $\sigma\leq 1-1/J$ and the integral is
positive, the time-averaged spreading is strictly
negative for $T$ sufficiently large.
\end{theorem}

\begin{proof}
Integrate~\eqref{eq:conc-spread} and divide by~$T$.
\end{proof}

\begin{remark}[Self-reinforcing spreading]
\label{rem:self-reinforce}
The concentration--spreading identity reveals a
feedback loop.  When $c^{(1)}<1$ (spreading active),
$\Sigma^{(1)}<0$ drives $\sigma$ downward.
Decreasing $\sigma$ maintains $d\sigma/dt\leq 0$,
which by Corollary~\ref{cor:spread-bound} ensures
$\Sigma^{(1)}<0$ continues.  The only way this
cycle breaks is if the viscous concentration term
$4\nu\sigma D(1-c_{\rm cen})/E$ overwhelms the
spreading and reverses $d\sigma/dt$.  But at the
balance point $d\sigma/dt=0$:
$\Sigma^{(1)} = -4\nu\sigma D(1-c_{\rm cen})/E<0$,
so even at balance the spreading is active.

The implication: $c^{(1)}\geq 1$ cannot persist.
Any interval with $c\geq 1$ drives $\sigma$ upward
(since $(c-1)T_{\rm tot}>0$ and the viscous term
helps concentration).  But $\sigma$ is bounded by
$1-1/J$, so the increase terminates.  At the
turning point ($d\sigma/dt=0$):
Theorem~\ref{thm:equil} gives $c<1$.
\end{remark}

\subsubsection{The excursion mechanism (N37)}

\medskip\noindent\textbf{N37a (Excursions require
enstrophy growth).}
At $\nu=0.005$ (1000 snapshots, $\Delta t=0.01$):
$c^{(1)}\geq 1$ at 92 snapshots, and \emph{every one}
has $dE^{(1)}/dt>0$ (92/92, zero violations).
The excursions are confined to the enstrophy growth
phase ($t\approx 3.5$--$4.8$); after the enstrophy
peak ($t\approx 6.0$), $c<1$ at all 402 subsequent
snapshots ($\max c = 0.667$).

\medskip\noindent\textbf{N37b (Shell-resolved
nonlinear redistribution).}
Decompose $R_{\rm NL} = R_{\rm NL}^{(L)}
+R_{\rm NL}^{(H)}$ (low shells $|k|<4$, high
shells $|k|\geq 4$).  Key finding:
$R_{\rm NL}^{(L)}<0$ at \emph{all} 1000 snapshots
--- the low shells always spread.  The excursions
$c>1$ are driven entirely by $R_{\rm NL}^{(H)}>0$
(high-shell nonlinear concentration) during cascade
filling ($t\approx 3.7$--$4.6$).  The cascade
deposits enstrophy anisotropically at high~$k$
(inheriting the anisotropy of the source modes)
faster than the viscous damping can isotropize the
new modes.  After the enstrophy peak, the deposition
rate drops, viscous isotropization catches up, and
$R_{\rm NL}^{(H)}$ becomes permanently negative.

\medskip\noindent\textbf{N37c (Viscosity dependence).}
The enstrophy peak and lock-in times:
\begin{center}
\renewcommand{\arraystretch}{1.2}
\begin{tabular}{@{}ccccc@{}}
\toprule
$\nu$ & Enst.\ peak & Last $c\geq 1$
  & $c\geq 1$ after peak & $\max c$ \\
\midrule
$0.002$ & $t=8.6$ & $t=4.3$ & 0 & $1.34$ \\
$0.005$ & $t=6.0$ & $t=4.7$ & 0 & $1.23$ \\
$0.010$ & $t=4.9$ & $t=5.4$ & 49 & $1.08$ \\
$0.020$ & $t=2.7$ & --- & 0 & $0.59$ \\
\bottomrule
\end{tabular}
\end{center}
At $\nu=0.002$ and $0.005$: the enstrophy peaks
\emph{after} the last $c\geq 1$ excursion --- the
excursions end before the cascade has fully
developed.  At $\nu=0.01$: the enstrophy peaks
\emph{before} the last excursion --- secondary
enstrophy growth phases drive 49 post-peak
excursions.  At $\nu=0.02$: $c<1$ always.

\medskip\noindent\textbf{N38 (Trilinear angular
structure at low shells).}
The drain ratio $T_j^{(L)}/E_j^{(L)}$ (per-sector
low-shell transfer per unit low-shell enstrophy)
reveals an unexpected asymmetry at $t=10$:
concentrated sectors ($\varepsilon_j>0$) drain at
rate $-0.107$; depleted sectors at rate $-0.586$
(ratio $0.18$).  The depleted sectors drain
$5.5\times$ faster per unit enstrophy.
Concentrated sectors \emph{retain} their low-$k$
enstrophy, contrary to the proportional-drain
hypothesis.

Despite this, $R_{\rm NL}^{(L)}\approx
+2\times 10^{-6}$ at $t=10$ --- essentially zero.
The low shells contribute negligibly to the
spreading or concentration of the field.  The
entire mechanism is at high shells:
$R_{\rm NL}^{(H)} = -1.85\times 10^{-3}$
(spreading), which is $1000\times$ larger in
magnitude.

This identifies the precise proof target: show
that $R_{\rm NL}^{(H)}<0$ (high-shell nonlinear
redistribution spreads) once the Gronwall transient
from cascade development has decayed.  In the
equilibrated regime: $T_j^{(H)}\approx 2\nu D_j^{(H)}$,
giving $R_{\rm NL}^{(H)}\approx
(4\nu D^{(H)}\sigma/E)(c_{\rm cen}^{(H)}-1)<0$
(since $c_{\rm cen}^{(H)}<1$ by viscous balance).
The excursions $c>1$ occur precisely when the
Gronwall residual from anisotropic cascade
deposition at high~$k$ exceeds this steady-state
spreading term.

\medskip\noindent\textbf{N39 (High-shell residual
correlation).}
Define the high-shell residual
$\delta_j^{(H)} = T_j^{(H)}-2\nu D_j^{(H)}$
(the departure of sector~$j$'s high-shell transfer
from Gronwall equilibrium).  The correlation
$\corr(\varepsilon_j,\delta_j^{(H)})$ measures
whether the departure concentrates ($>0$) or
spreads ($<0$).  Data ($128^3$, $\nu=0.005$):
\begin{center}
\renewcommand{\arraystretch}{1.2}
\begin{tabular}{@{}ccccc@{}}
\toprule
$t$ & $\corr(\varepsilon,\delta^{(H)})$
  & $\sum\varepsilon_j\delta_j^{(H)}$ & $c$ & $c_{\rm cen}$ \\
\midrule
$2.0$ & $+0.32$ & $+5.6\!\times\!10^{9}$ & $-0.44$ & --- \\
$4.0$ & $-0.95$ & $-8.1\!\times\!10^{10}$ & $1.19$ & --- \\
$4.2$ & $-0.98$ & $-1.1\!\times\!10^{11}$ & $1.23$ & $1.00$ \\
$5.0$ & $-0.99$ & $-1.2\!\times\!10^{11}$ & $0.69$ & --- \\
$7.0$ & $-0.99$ & $-6.7\!\times\!10^{10}$ & $0.57$ & --- \\
$10.0$ & $-0.88$ & $-2.4\!\times\!10^{10}$ & $0.62$ & $0.69$ \\
\bottomrule
\end{tabular}
\end{center}
The correlation is \emph{strongly negative}
($-0.88$ to $-0.99$) at every time $t\geq 4$.
The high-shell residual always spreads: concentrated
sectors have more negative residual (their high-$k$
enstrophy drains faster than the cascade replenishes).
The gap $c-c_{\rm cen}$ is \emph{negative} (helpful)
whenever the high-shell residual dominates.

The excursion $c>1$ at $t=4.2$ occurs despite the
negative high-shell residual because
$c_{\rm cen}=1.004>1$: the \emph{centroid itself}
exceeds~$1$ during explosive enstrophy growth.
Shell-resolved data confirms this: at $t=4.2$,
shell~$0$ ($|k|<2$) has a positive residual
$\sum\varepsilon_j\delta_j^{(0)} = +7.2\times 10^9$
(the TG modes are being amplified by vortex
stretching).  This low-$k$ concentration pushes
$c_{\rm cen}$ above~$1$.  After cascade development
($t\geq 5$), the low-$k$ contribution weakens,
$c_{\rm cen}$ drops to $0.69$, and the negative
high-shell residual ensures $c<c_{\rm cen}<1$.

\emph{Implication:} the obstruction to pointwise
$c<1$ is not the gap $c-c_{\rm cen}$ (which is
always helpful for $t\geq 4$) but the centroid
$c_{\rm cen}$ itself exceeding~$1$ during the
brief interval of explosive enstrophy growth.
The viscous balance theorem
(Theorem~\ref{thm:viscous-sep}) proves
$c_{\rm cen}<1$ with a margin that depends on
the Sobolev norms of~$u$.  During explosive growth,
$\|u\|_{H^2}$ increases and the margin shrinks;
the excursion $c_{\rm cen}>1$ reflects a
\emph{temporary} loss of the viscous balance
margin, not a failure of the spreading mechanism.

\begin{remark}[Summary of the proved results]
\label{rem:proved-summary}
The complete chain of proved theorems:
\begin{enumerate}[nosep]
\item $c^{(1)}$ is degree~$0$
  (Proposition~\ref{prop:c-degree-zero}).
\item $c_{\rm cen}<1$ for smooth NS solutions
  with anisotropic initial data, for
  $t\leq O(1/(\nu\kappa_L))$
  (Theorem~\ref{thm:viscous-sep}).
\item $\Sigma^{(1)} < d\sigma/dt$ whenever
  $c_{\rm cen}<1$
  (Corollary~\ref{cor:spread-bound}).
\item $c^{(1)}<1$ at every concentration
  equilibrium ($d\sigma/dt=0$) when
  $c_{\rm cen}<1$ and $T_{\rm tot}>0$
  (Theorem~\ref{thm:equil}).
\item Time-averaged $\Sigma^{(1)}<0$ for
  $T$ large enough
  (Theorem~\ref{thm:time-avg}).
\item $c^{(1)}\geq 1$ cannot persist: any
  excursion is self-correcting
  (Remark~\ref{rem:self-reinforce}).
\end{enumerate}
No steady-state hypothesis is needed.
High-resolution sampling (N36--N39) confirms
that $c^{(1)}>1$ occurs transiently during
cascade development ($\max c = 1.23$ at
$\nu=0.005$, 92/1000 snapshots).  The high-shell
residual correlation (N39) is strongly negative
($\corr(\varepsilon,\delta^{(H)})=-0.88$ to
$-0.99$ for $t\geq 4$): the gap $c-c_{\rm cen}$
is always helpful.  The excursion $c>1$ is driven
entirely by $c_{\rm cen}>1$ (the viscous balance
margin temporarily overwhelmed by explosive
low-$k$ enstrophy growth), not by the high-shell
nonlinear transfer.  After cascade development
($t\geq 5$): $c<c_{\rm cen}<1$ permanently.
\end{remark}


%======================================================================
\section{A Heuristic Mechanism: the $-i$ Rotation}
\label{sec:mechanism}
%======================================================================

The exact results of Sections~\ref{sec:general}--\ref{sec:TG}
and the numerical findings of Section~\ref{sec:results}
suggest a general mechanism.  \textbf{The arguments in this
section and the next are heuristic: they rely on
concentration and slow-variation approximations that are
not rigorously justified for general NS solutions.}
We present them as a proposed explanation, not as proved
theorems.

\subsection{Polarization}

\begin{definition}[Polarization]
\label{def:polarization}
For a mode $\hat u(k)\neq 0$, the polarization is
$\hat e(k) = \hat u(k)/|\hat u(k)|\in\C^3$.  When $u$ is
angularly concentrated in sector $\Omega_{j^*}$ with
slowly varying polarization, we write
$\hat u(q)\approx A(q)\hat e$ for
$\hat q\in\Omega_{j^*}$, where $A(q)=|\hat u(q)|$ and
$\hat e$ is the (approximately) common polarization
direction.  For a cross-sector mode
$p\notin\Omega_{j^*}$, we write
$\hat e_p = \hat u(p)/|\hat u(p)|$.
\end{definition}

\subsection{Cross-sector forcing under concentration}

Under the concentration hypothesis, the cross-sector
forcing at a mode $k\in\Omega_{j^*}$ takes the
approximate form
\begin{equation}\label{eq:Bcross-approx}
\hat B^{\rm cross}(k) \approx -iF_{\rm cross}(k)\hat e,
\end{equation}
where
\begin{equation}
F_{\rm cross}(k) = \sum_{\substack{p+q=k\\
\hat p\notin\Omega_{j^*},\;\hat q\in\Omega_{j^*}}}
(q\cdot\hat u(p))\,A(q).
\end{equation}
The key step: since $\hat u(k)\approx A_k\hat e$ and
$P(k)\hat e\approx\hat e$ (both from the narrow-sector
hypothesis), the transfer becomes
$T^{\rm cross}(k) \approx 2|k|A_k\,\im[F_{\rm cross}(k)]$.

Spreading of the dominant sector therefore requires
$\im[F_{\rm cross}(k)] < 0$ for most modes
$k\in\Omega_{j^*}$.

\subsection{Why the TG case is exact}

For TG, the concentration and narrow-sector hypotheses
are not approximations: the 8 modes have identical
amplitudes and the purely imaginary structure makes
$P(k)\hat e = \hat e$ exact (since $V(k)\perp k$
for every $k\in\cK$).  The approximate mechanism becomes
exact, and $\im[F_{\rm cross}]<0$ is established by
Theorem~\ref{thm:TG}.

\begin{remark}[Rigorous content of the mechanism
for purely imaginary fields]
\label{rem:mechanism-rigorous}
The ``$-i$ rotation'' described heuristically in
this section has a rigorous counterpart for
purely imaginary fields.  By
Proposition~\ref{prop:real-transfer}, when
$\hat u = iV$ with $V\in\R^3$, the transfer
$T(k) = 2|k|\sum_{p+q=k}(q\cdot V(p))(V(k)\cdot V(q))$
is a real sum: no complex phases appear.  The
``phase rotation'' becomes a statement about the
geometry of the real vectors: the
factor $(q\cdot V(p))$ measures the alignment of
the advecting velocity with the wavevector $q$,
while $(V(k)\cdot V(q))$ measures the alignment
of the advected velocity with the mode being
transferred to.  The cross-sector depletion
observed numerically amounts to a systematic
tendency for these two dot products to have
opposite signs when $k$ and $p$ belong to
different angular sectors.  By
Corollary~\ref{cor:TG-real-orbit}, this
reformulation holds along the entire TG orbit.
\end{remark}

\subsection{The $-i$ rotation in general}

The factor $-i$ in the advection operator
$\widehat{(u\cdot\nabla)u}(k) = \sum_{p+q=k}
iq_l\hat u_l(p)\hat u(q)$ creates a $\pi/2$ phase
rotation between the ``natural'' phase of a triadic
interaction and the resulting forcing.  For concentrated
fields, this means cross-sector forcing (where the
advecting velocity comes from a different angular
direction) arrives with a systematic phase offset
relative to the existing velocity in the concentrated
sector.  The offset biases $\cos\alpha$ toward negative
values (depletion rather than amplification).

For TG, the offset is exactly $\pi$ (perfect
anti-alignment), because the purely imaginary structure
locks the accumulated phase to $i^2=-1$.  For general
fields, the offset is partial, producing
$\cos\alpha\approx-0.6$ on average in high-$p$ sectors
(N8).

The triad phase locking observed in N10 provides
direct microscopic evidence for this mechanism.
Equation~\eqref{eq:triad-imaginary} shows that the
purely imaginary structure forces every individual
triad integrand $z(k;p,q)$ onto the imaginary axis,
so the nonlinear transfer $T(k) = 2|k|\,\im[\sum z]$
is entirely controlled by the signs of the real
products $(q\cdot V(p))(V(k)\cdot V(q))$.  For
random-phase fields, these products are complex
with uniformly distributed arguments, and the
$\pm\pi/2$ locking disappears.  The cascade phase
structure required by Conjecture~\ref{conj:enstrophy}
is therefore precisely the persistence of
$\hat u\approx iV$, which the NS dynamics maintain
but arbitrary initial data need not possess.

The extension to $\Sigma^{(s)}$ for general~$s$ (N12)
follows because the factor $|k|^{2s}$ in $T^{(s)}(k)$
is a positive real weight that preserves the triad
phase structure.  The $\pm\pi/2$ locking drives
spreading at every Sobolev level simultaneously: the
nonlinear term redistributes the \emph{entire regularity
structure} of the field --- including the enstrophy
($s=1$) and higher derivatives ($s=3/2$) --- away from
angular concentration.

The depletion transition of N16 reveals a threshold
effect: the $-i$ mechanism engages only when angular
concentration exceeds $\cC^{(1)}\approx 0.15$.  Below
this threshold, the same NS phases amplify rather than
deplete, because there is no dominant sector for
cross-sector transfer to target.  The spreading
robustness of N17 shows that 50\% of the phase
structure can be destroyed before the mechanism fails,
establishing that the $-i$ rotation is not fragile.

%======================================================================
\section{A Conditional Sign Criterion}
\label{sec:conditional}
%======================================================================

\textbf{This section presents a conditional result:
depletion under an explicit sign hypothesis.  It does
not constitute a rigorous reduction to Kolmogorov's
forward cascade.}

\subsection{Quasi-static cross-sector modes}

If cross-sector modes equilibrate viscously
($\nu|k'|^2$ large compared to the forcing variation
rate), the NS equation gives
$\hat u(k')\approx\hat B(k')/(\nu|k'|^2)$.  Since
$\hat B(k')\approx -iG(k')$ for modes created by
within-sector triads (with
$G(k')=P(k')\sum_{p,q\in\Omega_{j^*}}
(q\cdot\hat u(p))\hat u(q)$), the cross-sector mode
carries a phase approximately $\pi/2$ behind its source.

This approximation is rigorous in the dissipation range
where $\nu|k'|^2\gg\tau_{\rm NL}^{-1}$.  In the
inertial range, where the nonlinear turnover time
competes with viscous damping, the quasi-static
assumption softens, but the $\pi/2$ bias from the $-i$
factor persists as a leading-order effect.

\subsection{The sign criterion}

When quasi-static cross-sector modes feed back to
$k\in\Omega_{j^*}$, the feedback transfer involves
\begin{equation}
\im[\gamma(k,p')] = \im\Bigl[q\cdot P(p')
\sum_{\substack{r+s=p'\\r,s\in\Omega_{j^*}}}
(s\cdot\hat u(r))\hat u(s)\Bigr],
\end{equation}
where $q=k-p'$.  If this quantity is negative for
the dominant intermediate wavenumbers $p'$, the feedback
depletes the concentrated sector.

\begin{remark}[Not a reduction to K41]
We emphasize that $\im[\gamma]<0$ is a sign condition
on the same trilinear functional whose negativity
one wants to establish.  It is \emph{not} equivalent
to the standard forward cascade (constant energy flux
through spectral shells).  The connection to Kolmogorov's
theory is suggestive --- both involve the sign of
triadic phase interactions --- but the condition here is
sector-projected and geometry-dependent, not a direct
consequence of K41.  We state it as a conditional
criterion, not as a proved reduction.
\end{remark}

%======================================================================
\section{Conjecture and Discussion}
\label{sec:conjecture}
%======================================================================

\begin{conjecture}[Enstrophy spreading for developed
turbulence]
\label{conj:enstrophy}
For smooth NS solutions on $\T^3$ that have developed
cascade phase structure (in the sense of having evolved
from generic initial data for a time sufficient to
establish a forward energy cascade): if
$\cC^{(1)}(u(t))>1/J$, then $\Sigma^{(1)}(u(t))<0$.
\end{conjecture}

The evidence for this conjecture consists of:
the exact algebraic extension of N12 (the
$|k|^{2s}$ weight preserves the triad phase
structure for all~$s$);
the TG result (N12: $\Sigma^{(1)}<0$ at every
snapshot);
the viscosity independence (N14);
the multi-IC evidence (N15: all random broadband
ICs spread at $s=1$, 30/30 negative snapshots);
the jet counterexample (N15: cascade phase structure
remains necessary);
the depletion transition (N16: the mechanism engages
above $\cC^{(1)}\approx 0.15$);
and the phase robustness (N17: spreading survives
50\% phase destruction);
the surgical concentration test (N18: $\Sigma^{(1)}<0$
at all 44 snapshots under adversarial concentration up
to $\cC^{(1)}=0.50$, with monotone relaxation);
and the constitutive law (N19: $g^{(1)}\approx
0.47\pm 0.04$ across the full concentration range).
The randomization tests (N20) further establish
that this spreading is a dynamical property of
the NS flow, not a consequence of any static
feature of the velocity field.
The dynamical creation test (N21) confirms that
NS evolution drives random initial data to
$\Sigma^{(1)}<0$ within one turnover time
($50\%\to 100\%$), establishing spreading as a
property of the NS attractor.
The contraction asymmetry (N22) reveals that the
basin $\{\Sigma^{(1)}<0\}$ is nearly absorbing:
all seeds starting with $\Sigma^{(1)}>0$ flip
to negative by $t=1$, while seeds starting
negative persist at $\geq 93\%$ through $t=10$.
This is confirmed at $256^3$ (N23) and across
$\nu\in[0.002,0.02]$ (N24), ruling out resolution
artifacts and confirming a purely nonlinear
creation mechanism.

Conjecture~\ref{conj:enstrophy} implies the $s=1/2$
version, since the triad phase structure that drives
$\Sigma^{(1)}<0$ simultaneously drives
$\Sigma^{(s)}<0$ for all $s\geq 0$ (N12).

\begin{conjecture}[$\Hdot{1/2}$ spreading]
\label{conj:main}
Under the same hypotheses: if $\cC(u(t))>1/J$,
then $\Sigma(u(t))<0$.
\end{conjecture}

The experiments of Section~\ref{sec:affine-experiments}
identify a sharper, dimensionless form of the
spreading conjecture:

\begin{remark}[Affine spreading: proved and open]
\label{conj:affine}
The dimensionless form of the spreading question
asks: does $c^{(1)}(u(t))<1$ hold along smooth
NS trajectories with $T_{\rm tot}>0$?
High-resolution sampling (N36: 1000 snapshots,
$\Delta t=0.01$) reveals that $c^{(1)}$ can
\emph{transiently} exceed~$1$ during cascade
development: $\max c^{(1)} = 1.23$ for TG at
$\nu=0.005$ (at $t\approx 4.2$, 92/1000 snapshots),
and larger excursions occur at lower viscosity
($\max c = 1.34$ at $\nu=0.002$).  At $\nu=0.02$,
$c^{(1)}<1$ at all 1000 snapshots.

The \emph{proved} results
(Section~\ref{sec:concentration-dynamics}) are:
\begin{enumerate}[nosep,leftmargin=2em]
\item $c_{\rm cen}<1$ always
  (Theorem~\ref{thm:viscous-sep});
\item $c^{(1)}<1$ at every concentration
  equilibrium (Theorem~\ref{thm:equil});
\item time-averaged $\Sigma^{(1)}<0$
  (Theorem~\ref{thm:time-avg});
\item $c^{(1)}\geq 1$ cannot persist
  (Remark~\ref{rem:self-reinforce}).
\end{enumerate}
These establish that spreading is the dominant
regime; the transient excursions are self-correcting
and confined to the cascade development phase.
\end{remark}

The proved results imply a weaker but still
useful version of the spreading conjecture:

\begin{conjecture}[Time-averaged enstrophy spreading]
\label{conj:time-avg}
For smooth NS solutions on $\T^3$ with
$\cC^{(1)}>1/J$ and anisotropic initial data:
$\frac{1}{T}\int_0^T\Sigma^{(1)}\,dt < 0$
for all $T$ sufficiently large.
\end{conjecture}

This is \emph{proved} by Theorem~\ref{thm:time-avg}
under the condition $c_{\rm cen}<1$
(itself proved by Theorem~\ref{thm:viscous-sep}).
The original pointwise conjecture
(Conjecture~\ref{conj:enstrophy}) remains open,
with the transient excursions of N36 showing that
the cascade development phase requires separate
treatment.

\begin{remark}[Connection to regularity]
\label{rem:regularity-final}
The regularity question reduces to three coupled
statements about the NS nonlinearity.

\smallskip\noindent\textbf{Step (i): Enstrophy
spreading.}
Conjecture~\ref{conj:enstrophy}:
$\Sigma^{(1)}(u)<0$ whenever $\cC^{(1)}>1/J$.
Evidence: N12 (all Sobolev levels), N14 (viscosity
independence), N15 (generality), N17 (robustness),
N18 (44/44 surgical snapshots).

\smallskip\noindent\textbf{Step (ii):
Stretching--spreading coupling.}
The surgical concentration tests of N18 provide
direct evidence: under adversarial enstrophy
concentration ($\cC^{(1)}$ up to $0.50$), the ratio
$\mathcal{R} = |\mathcal{W}|/(|\Sigma^{(1)}|
\cdot\|\omega\|^2)$ remains in $[0.6, 2.2]$ and
\emph{decreases} with concentration
(Table~\ref{tab:surgical}).

\smallskip\noindent\textbf{Step (iii):
Dynamical contraction.}
The randomization tests of N20 show that
$\Sigma^{(1)}<0$ is not a property of the
amplitude spectrum, the purely imaginary
structure, or the divergence-free constraint:
all randomizations give $48\%$--$50\%$.
The creation test (N21) shows the
converse: evolving random ICs under NS drives
$\Sigma^{(1)}<0$ from $50\%$ to $100\%$ within
one turnover time.  The contraction asymmetry
(N22) establishes that this is not mere drift
but active contraction: all 15/15 positive seeds
are pulled negative by $t=1$, while
15/15 negative seeds persist through $t=1$ (and
$\geq 93\%$ through $t=10$).
The basin $\{\Sigma^{(1)}<0\}$ is nearly absorbing,
and this holds at $256^3$ (N23) and across a
decade of viscosity (N24).
The remaining analytical challenge is
to prove that the NS flow is a contraction
toward $\Sigma^{(1)}<0$ for fields with
$\cC^{(1)}>1/J$.
\end{remark}

The elevation of the conjecture from $s=1/2$ to $s=1$
has a structural consequence.  At $s=1/2$, the
connection to regularity required an intermediate
``angular Bernstein inequality'' relating concentration
of $\Hdot{1/2}$ energy to pointwise norms.  At $s=1$,
this intermediate step is bypassed: enstrophy
concentration is already the regularity-relevant
quantity through the vortex-stretching term
in~\eqref{eq:enstrophy} and the Beale--Kato--Majda
criterion~\cite{BKM}.  The regularity problem
reduces to questions about the trilinear form and its
angular geometry --- precisely the domain where the
algebraic identities of
Sections~\ref{sec:general}--\ref{sec:TG} provide
leverage.  However, the homogeneity analysis of
Section~\ref{sec:homogeneity} shows that proving
the sign conjecture alone does not suffice for
regularity: the constitutive coefficient is degree~1,
and a separate argument (connecting angular isotropy
to magnitude bounds on vortex stretching) would be
needed to close the gap.

\begin{remark}[Reformulation using the proved tools]
\label{rem:reformulation}
The algebraic results of Section~\ref{sec:general}
reformulate the conjecture target as follows.
By Proposition~\ref{prop:linear-response}, the
conjecture $\Sigma^{(1)}<0$ when $\cC^{(1)}>1/J$
is equivalent (at leading order near isotropy) to
\begin{equation}\label{eq:target}
\sum_j \varepsilon_j^{(1)}\, T_j^{(1)} < 0
\quad\text{whenever}\quad
\sum_j (\varepsilon_j^{(1)})^2 > 0,
\end{equation}
where $\varepsilon_j^{(1)} = p_j^{(1)}-1/J$ is
the enstrophy sector excess and
$T_j^{(1)} = \sum_{\hat k\in\Omega_j}
T^{(1)}(k;u)$.  By
Corollary~\ref{cor:universal-s}, the transfer
$T^{(1)}(k)$ has the same sign as $T(k)$, and
by Proposition~\ref{prop:real-transfer}, for
fields with $\hat u=iV$ it reduces to a real
bilinear form in the sector excess and the
triad sums.  The contraction statement (N22) ---
that the NS dynamics pull $\Sigma^{(1)}>0$ states
toward $\Sigma^{(1)}<0$ --- is
condition~\eqref{eq:target} asserted for NS
trajectories that have developed cascade phase
structure.
\end{remark}

\subsection{The homogeneity obstruction}
\label{sec:homogeneity}

The proved algebraic tools and numerical findings
suggest a natural path toward regularity:
if the constitutive law $\Sigma^{(1)}\leq
-g_0(\cC^{(1)}-1/J)$ and the stretching--spreading
coupling $|\mathcal{W}|\leq R_0|\Sigma^{(1)}|
\|\omega\|^2$ held with \emph{dimensionless}
constants $g_0,R_0$, a Gr\"onwall argument would
give at-most-exponential enstrophy growth and
hence global regularity.  However, the
homogeneity structure of the NS equations
prevents this.

\begin{proposition}[Homogeneity of the
constitutive coefficient]
\label{prop:homogeneity}
Under $u\mapsto\lambda u$ ($\lambda>0$):
\begin{enumerate}[leftmargin=2em,label=\textup{(\roman*)}]
\item $\Sigma^{(1)}(\lambda u) = \lambda\,\Sigma^{(1)}(u)$
  \textup{(}degree~$1$, by
  Proposition~\textup{\ref{prop:homog})}.
\item $\cC^{(1)}(\lambda u) = \cC^{(1)}(u)$
  \textup{(}degree~$0$\textup{)}.
\item The constitutive coefficient
  $g^{(1)}(u) := -\Sigma^{(1)}(u)/
  (\cC^{(1)}(u)-1/J)$ satisfies
  $g^{(1)}(\lambda u) = \lambda\,g^{(1)}(u)$
  \textup{(}degree~$1$\textup{)}.
\end{enumerate}
In particular, $g^{(1)}$ is not a dimensionless
constant: it scales with the amplitude of $u$.
\end{proposition}

\begin{proof}
Part (i) is Proposition~\ref{prop:homog} at $s=1$.
Part (ii): $p_j^{(1)}(\lambda u)
= E_j^{(1)}(\lambda u)/E^{(1)}(\lambda u)
= \lambda^2 E_j^{(1)}(u)/(\lambda^2 E^{(1)}(u))
= p_j^{(1)}(u)$.  Part (iii) follows from (i) and (ii).
\end{proof}

\begin{remark}[Why angular spreading does not
directly imply regularity]
\label{rem:obstruction}
The constitutive law $\Sigma^{(1)}\leq
-g^{(1)}(\cC^{(1)}-1/J)$ with degree-1
coefficient $g^{(1)}$ controls the
\emph{direction} of nonlinear enstrophy transfer
(from concentrated to dispersed sectors) but not
the \emph{magnitude} of vortex stretching.  To see
this, note that the stretching--spreading coupling
(H2 in the notation of N18--N19) gives
$|\mathcal{W}| \leq R_0\,|\Sigma^{(1)}|\,
\|\omega\|^2$, where both sides are degree~3.
Substituting $|\Sigma^{(1)}| \leq
C\,\|\omega\|$ (degree~1), we obtain
$|\mathcal{W}| \leq C'\|\omega\|^3$ --- the
standard cubic bound, which does not improve
the scaling exponent.  Finite-time blow-up is
controlled by whether $|\mathcal{W}|$ grows
faster than $\nu\|\nabla\omega\|^2$, a
question about exponents (degree~3 vs.\
degree~2), not about angular signs (degree~0).

The analogy with Bakry--\'Emery curvature
(Section~\ref{sec:bakry-emery}) illuminates
the gap: in the BE framework, $\Gamma_2\geq
\kappa\Gamma$ involves a degree-0 constant
$\kappa$ because both $\Gamma_2$ and $\Gamma$
are degree~2 in $f$.  In the NS setting,
$\Sigma^{(1)}$ is degree~1 and
$\cC^{(1)}-1/J$ is degree~0, so the
``curvature'' is degree~1 --- fundamentally
different from a Bakry--\'Emery constant.
Bridging this gap would require a
Fourier-to-physical-space argument showing that
angular isotropy of enstrophy (controlled by
spreading) implies reduced vorticity alignment
(controlling stretching magnitude), potentially
connecting to the Constantin--Fefferman
direction~\cite{CF93}.

The affine ratio $c^{(1)}$
(Definition~\ref{def:affine-ratio}) provides a
precise decomposition of the obstruction.
The identity $\Sigma^{(1)} = 2(c^{(1)}-1)\,
T_{\rm tot}\sum\varepsilon_j^2/E$ separates
the spreading mechanism into two factors:
(a)~the affine ratio $c^{(1)}<1$, which is
degree~0 and carries no homogeneity obstruction
(Proposition~\ref{prop:c-degree-zero}); and
(b)~the prefactor $T_{\rm tot}/E$, which is
degree~1 and encodes the classical stretching
control problem.  The angular spreading program
addresses~(a); the gap~(b) is the standard
regularity question of bounding vortex stretching
relative to enstrophy.
\end{remark}

The numerical data (N25) confirm the degree-1
nature of $g^{(1)}$: across viscosities at $t=5$,
$g^{(1)}$ ranges from $0.14$ ($\nu=0.02$) to
$0.24$ ($\nu=0.002$), reflecting different energy
levels of the respective TG templates.  At
$256^3$ vs.\ $128^3$ with the same $\nu=0.005$,
$g^{(1)}$ is $27\%$ larger at the higher
resolution (more active modes, more energy).
Within a fixed ensemble (same amplitude spectrum),
$g^{(1)}$ varies only with angular structure
(seeds at the same $\nu$ and resolution), but
\emph{between} ensembles the amplitude dependence
is clear.

\subsection{Viscous--nonlinear balance in
concentration dynamics (N25)}
\label{sec:visc-balance}

The full time derivative of $\cC^{(1)}$ receives
contributions from both the nonlinear term
($\Sigma^{(1)}$, which spreads) and the viscous
term (which concentrates by selectively damping
small-scale modes).  At $256^3$ ($\nu=0.005$,
seed-averaged):
\begin{center}
\renewcommand{\arraystretch}{1.2}
\begin{tabular}{@{}ccccc@{}}
\toprule
$t$ & $\cC^{(1)}-1/J$
  & $\Sigma^{(1)}$ & visc.\ contr.
  & $|\text{visc}/\text{NL}|$ \\
\midrule
$1.0$ & $4.4\times 10^{-3}$
  & $-1.2\times 10^{-3}$ & $+1.1\times 10^{-3}$
  & $0.87$ \\
$2.0$ & $4.1\times 10^{-3}$
  & $-1.6\times 10^{-3}$ & $+1.3\times 10^{-3}$
  & $0.81$ \\
$3.0$ & $4.1\times 10^{-3}$
  & $-1.4\times 10^{-3}$ & $+1.4\times 10^{-3}$
  & $1.00$ \\
$5.0$ & $5.3\times 10^{-3}$
  & $-1.0\times 10^{-3}$ & $+1.6\times 10^{-3}$
  & $1.57$ \\
\bottomrule
\end{tabular}
\end{center}
The viscous and nonlinear contributions are comparable
at all times ($|\text{visc}/\text{NL}|\approx 1$).
For $t\leq 2$, the nonlinear spreading narrowly wins
and $\cC^{(1)}$ decreases; for $t\geq 3$, the viscous
concentration wins (as small-scale modes are damped)
and $\cC^{(1)}$ increases.  The nonlinear spreading
$\Sigma^{(1)}<0$ persists throughout, but it is not
the sole determinant of the concentration dynamics.

\subsection{Connection to the Bakry--\'Emery
curvature framework}
\label{sec:bakry-emery}

The spreading conjecture has a structural
interpretation in the language of the Bakry--\'Emery
$\Gamma$-calculus~\cite{Fontes2026b,Fontes2026c},
with a key disanalogy at the level of homogeneity.

In the standard Bakry--\'Emery theory for a Markov
generator $L$, the curvature condition
$\Gamma_2(f,f)\geq\kappa\,\Gamma(f,f)$
implies exponential decay of the Dirichlet energy:
$\mathcal{E}(P_t f) \leq e^{-2\kappa t}\mathcal{E}(f)$.
The constant $\kappa$ is degree~0 (dimensionless)
because both $\Gamma_2$ and $\Gamma$ are degree~2
in $f$.

The NS spreading conjecture has a parallel structure:
\begin{center}
\renewcommand{\arraystretch}{1.2}
\begin{tabular}{@{}ll@{}}
\toprule
\textbf{Bakry--\'Emery} & \textbf{NS spreading} \\
\midrule
$\Gamma(f,f)$ (carr\'e du champ)
  & $\cC^{(1)}-1/J$ (concentration excess) \\
$\Gamma_2(f,f)$ (iterated CdC)
  & $-\Sigma^{(1)}$ (spreading force) \\
$\Gamma_2\geq\kappa\,\Gamma$ (curvature)
  & $-\Sigma^{(1)}\geq g^{(1)}(\cC^{(1)}-1/J)$
  (constitutive law) \\
$\kappa$ degree~0 (dimensionless)
  & $g^{(1)}$ \textbf{degree~1} (scales with $\|u\|$) \\
$\rho^*<1$ (off-diagonal decay)
  & inter-scale anti-correlation (N20) \\
Gaussian universality at high chaos
  & enstrophy spreading for all $s\geq 0$ (N12) \\
\bottomrule
\end{tabular}
\end{center}
The dictionary is exact at the level of
\emph{sign} structure: $\Gamma_2\geq\kappa\Gamma$
with $\kappa>0$ is the BE analogue of
$\Sigma^{(1)}<0$ when $\cC^{(1)}>1/J$.
The inter-scale anti-correlation (N20) plays the
role of the curvature deficit at $H\neq 1/2$
in the fBM covariance theory.

The analogy breaks at \emph{magnitude}: BE
curvature $\kappa$ is degree~0 and directly
controls semigroup decay, while the NS
``curvature'' $g^{(1)}$ is degree~1 and controls
concentration dynamics but not enstrophy growth.
This homogeneity mismatch (Proposition~\ref{prop:homogeneity})
is the fundamental obstruction to obtaining
regularity from angular spreading alone.

The numerical finding N20 (inter-scale
anti-correlation ratio $-263\times$ for NS vs.\
$+144\times$ for random fields) is the direct
analogue of the Bakry--\'Emery curvature deficit:
NS solutions have negative off-diagonal
structure (like $H=1/2$, independent increments),
while random fields have positive correlations
(like $H\neq 1/2$).  Proving
Conjecture~\ref{conj:enstrophy} is
analogous to proving positive Bakry--\'Emery
curvature for the NS nonlinearity --- but
the trilinear (degree~3) nature of the
advection term means this curvature lives in
a different homogeneity class than the standard
theory.

\subsection{The alignment deviation and
vortex stretching}
\label{sec:alignment}

The homogeneity obstruction
(Proposition~\ref{prop:homogeneity}) shows that
angular spreading alone cannot bound the
\emph{magnitude} of vortex stretching.  However,
the following identity reveals the precise
geometric factor that spreading controls.

At any point $x$, the local vortex stretching rate
is $\alpha(x) = \hat\omega\cdot S\hat\omega$, where
$S = \frac{1}{2}(\nabla u + \nabla u^T)$ is the
strain tensor with trace-free eigenvalues
$\lambda_1+\lambda_2+\lambda_3=0$ and $\theta_i$ is
the angle between the vorticity $\omega$ and the
$i$-th eigenvector of $S$.  Then:
\begin{equation}\label{eq:alpha-deviation}
\alpha(x) = \sum_{i=1}^3 \lambda_i\cos^2\theta_i
= \sum_{i=1}^3 \lambda_i
\Bigl(\cos^2\theta_i - \tfrac{1}{3}\Bigr),
\end{equation}
using $\sum\lambda_i=0$.  Define the
\emph{alignment deviation}
\begin{equation}\label{eq:Dalign}
D(x) := \Bigl(\sum_{i=1}^3
\bigl(\cos^2\theta_i(x)-\tfrac{1}{3}\bigr)^2
\Bigr)^{1/2} \in \bigl[0,\sqrt{2/3}\,\bigr].
\end{equation}
By Cauchy--Schwarz:
$|\alpha(x)| \leq |S(x)|\,D(x)$.

The identity~\eqref{eq:alpha-deviation} has a
transparent meaning:
\emph{isotropic alignment kills stretching}.
If $\cos^2\theta_i = 1/3$ for all $i$ (vorticity
equally distributed among the strain eigendirections),
then $\alpha(x)=0$ exactly, regardless of $|S(x)|$
and $|\omega(x)|$.  The vortex stretching term
$\mathcal{W} = \int\omega\cdot S\omega\,dx$
therefore satisfies
\begin{equation}\label{eq:W-alignment}
|\mathcal{W}| \leq \int |\omega(x)|^2\,|S(x)|\,
D(x)\,dx
\leq \|D\|_{L^\infty}\int |\omega|^2\,|S|\,dx.
\end{equation}
The factor $\|D\|_{L^\infty}\in[0,\sqrt{2/3}\,]$
is degree~0 --- purely geometric.  The remaining
integral $\int|\omega|^2|S|\,dx$ is degree~3, giving
the standard cubic bound.  Spreading reduces the
\emph{constant} (via $\|D\|_{L^\infty}<\sqrt{2/3}$)
but not the \emph{exponent}.

The connection to the Fourier-space results is:
\begin{itemize}[leftmargin=2em]
\item Angular spreading of enstrophy
  ($\cC^{(1)}\approx 1/J$, $\Sigma^{(1)}<0$) means
  vorticity is a superposition of contributions from
  $\sim J$ incoherent angular sectors in Fourier space.
\item Physical-space vorticity-strain alignment is
  controlled by the Fourier-space angular distribution:
  concentrated enstrophy ($\cC^{(1)}\approx 1$) allows
  coherent alignment ($D\approx\sqrt{2/3}$), while
  spread enstrophy forces partial cancellation.
\item The numerical findings N20 (inter-scale
  anti-correlation) and N22 (contraction toward
  $\Sigma^{(1)}<0$) provide Fourier-space evidence that
  NS dynamics push $D$ toward zero.
\end{itemize}
The precise relationship between Fourier-space
concentration and physical-space alignment is
sharpened by the following identity, which
connects the paper's sector transfer framework
directly to vortex stretching.

\begin{proposition}[Stretching--spreading identity]
\label{prop:stretching-spreading}
The total enstrophy-level sector transfer equals
twice the vortex stretching:
\begin{equation}\label{eq:Ttot-W}
T_{\rm tot}^{(1)} := \sum_j T_j^{(1)} = 2\mathcal{W}.
\end{equation}
Combined with the linear response
(Proposition~\textup{\ref{prop:linear-response}}),
this gives the exact identity:
\begin{equation}\label{eq:W-Sigma}
2\mathcal{W}\,\bigl(\cC^{(1)}-1/J\bigr)
= \sum_j \varepsilon_j^{(1)} T_j^{(1)}
- \frac{E^{(1)}}{2}\,\Sigma^{(1)},
\end{equation}
valid whenever $\cC^{(1)}>1/J$.
\end{proposition}

\begin{proof}
The enstrophy $E^{(1)}=\sum_k|k|^2|\hat u(k)|^2
=\|\omega\|_{L^2}^2$ evolves as
$\frac{d}{dt}E^{(1)} = -2\nu\|\nabla\omega\|^2
+ 2\mathcal{W}$
(from the enstrophy
equation~\eqref{eq:enstrophy}).  In the paper's
Fourier framework,
$\frac{d}{dt}E^{(1)} = -2\nu\sum_k|k|^4|\hat u|^2
+ T_{\rm tot}^{(1)}$.
Equating the nonlinear contributions
gives~\eqref{eq:Ttot-W}.
For~\eqref{eq:W-Sigma}, substitute
$T_{\rm tot}^{(1)}=2\mathcal{W}$ into the
linear response formula
(Proposition~\ref{prop:linear-response}):
$(E^{(1)}/2)\Sigma^{(1)}
= \sum_j\varepsilon_j T_j^{(1)}
- T_{\rm tot}^{(1)}(\cC^{(1)}-1/J)$
and rearrange.
\end{proof}

\begin{remark}[What the identity reveals]
\label{rem:W-Sigma}
The identity~\eqref{eq:W-Sigma} decomposes vortex
stretching into two competing contributions:
\begin{enumerate}[leftmargin=2em,label=\textup{(\roman*)}]
\item The sector-excess correlation
  $\sum\varepsilon_j T_j^{(1)}$: this is negative
  when high-enstrophy sectors receive below-average
  transfer (the spreading mechanism of N13).  It
  \emph{reduces} $\mathcal{W}$.
\item The spreading force $-(E^{(1)}/2)\Sigma^{(1)}$:
  when $\Sigma^{(1)}<0$ (spreading), this is
  \emph{positive} and contributes \emph{to}
  stretching.  Enstrophy redistribution across sectors
  does not reduce total enstrophy growth --- it
  makes the growth more uniform.
\end{enumerate}
The Pearson correlation $\rho=\corr(\varepsilon_j,
T_j^{(1)})=-0.43$ (N13) is degree~0 and measures
how efficiently the negative sector-excess
correlation offsets the spreading contribution.
The identity makes precise why spreading controls
the \emph{partition} of stretching across sectors
but not its \emph{total magnitude}.
\end{remark}

\begin{remark}[Combining the two factorizations]
\label{rem:two-factorizations}
The stretching $\mathcal{W}$ is now factorized
from two independent directions:
\begin{itemize}[leftmargin=2em]
\item \textbf{Physical space} (eq.~\eqref{eq:W-alignment}):
  $|\mathcal{W}|\leq\|D\|_{L^\infty}\int|\omega|^2|S|\,dx$.
  Geometry ($D$, degree~0) times magnitude
  ($\int|\omega|^2|S|$, degree~3).
\item \textbf{Fourier space} (eq.~\eqref{eq:W-Sigma}):
  $\mathcal{W}=[\sum\varepsilon_j T_j
  -(E/2)\Sigma]/(2(\cC-1/J))$.
  Sector correlation (degree~3) minus spreading
  (degree~3), divided by concentration excess
  (degree~0).
\end{itemize}
Equating gives a constraint between Fourier-space
quantities and physical-space alignment:
\begin{equation}\label{eq:D-constraint}
\|D\|_{L^\infty} \geq
\frac{|\sum\varepsilon_j T_j^{(1)}
  -(E^{(1)}/2)\Sigma^{(1)}|}
{2(\cC^{(1)}-1/J)\int|\omega|^2|S|\,dx}.
\end{equation}
The left side is physical-space geometry; the right
is computed entirely from Fourier-space angular data.
This implicit \emph{lower} bound on $D$ from Fourier
data is proved.  The conjectured \emph{upper} bound
(transfer inequality: $\|D\|\leq h(\cC^{(1)})$)
was tested experimentally and found to be
\emph{false}: $D$ increases with $\cC^{(1)}$, not
decreases (Section~\ref{sec:alignment-experiment}).
The actual depletion mechanism is eigenvalue selection,
not alignment isotropy.
\end{remark}

The following lemma provides the analytical basis
for the transversality suppression of within-sector
triads, which the divergence-free constraint imposes.

\begin{lemma}[Transversality depletion]
\label{lem:transversality}
Let $u$ be a smooth divergence-free field on $\T^3$
and $\Omega_j$ a sector of half-angle $\theta_j$.
For any within-sector triad ($\hat p, \hat q,
\hat k\in\Omega_j$ with $p+q=k$), the advection
coupling satisfies
\begin{equation}\label{eq:transversality}
|q\cdot\hat u(p)| \leq
2\theta_j\,|q|\,|\hat u(p)|.
\end{equation}
Consequently, the within-sector contribution to
the per-mode transfer at $k\in\Omega_j$ satisfies
\begin{equation}\label{eq:Twithin-bound}
\bigl|T_j^{\rm within}(k)\bigr|
\leq 4\theta_j\,|k|\!\!\!
\sum_{\substack{p+q=k\\\hat p,\hat q\in\Omega_j}}
\!\!|q|\,|\hat u(p)|\,|\hat u(k)|\,|\hat u(q)|,
\end{equation}
compared with the generic (unsuppressed) bound
$|T(k)| \leq 2|k|\sum_{p+q=k}
|q|\,|\hat u(p)|\,|\hat u(k)|\,|\hat u(q)|$.
The ratio is at most $2\theta_j$.
\end{lemma}

\begin{proof}
Decompose $q = (q\cdot\hat p)\hat p + q_{\perp p}$
with $q_{\perp p}\perp p$.  Since $\hat u(p)\perp p$
(divergence-free):
$q\cdot\hat u(p) = q_{\perp p}\cdot\hat u(p)$.
For $\hat p,\hat q\in\Omega_j$:
$\angle(p,q)\leq 2\theta_j$, so
$|q_{\perp p}| = |q|\sin\angle(p,q)
\leq 2\theta_j|q|$.
Therefore
$|q\cdot\hat u(p)| \leq |q_{\perp p}|\,|\hat u(p)|
\leq 2\theta_j|q||\hat u(p)|$.
Substituting into the factored form
(Proposition~\ref{prop:factored}) with the
$|k|^{2s}$ weight at $s=1$ gives~\eqref{eq:Twithin-bound}.
\end{proof}

\begin{remark}[The self-regulating feedback]
\label{rem:feedback}
Lemma~\ref{lem:transversality} reveals a
self-regulating mechanism inherent in the
divergence-free constraint:
\begin{itemize}[leftmargin=2em]
\item \textbf{Concentrated configurations}
  ($\cC^{(1)}\to 1$): most triads are within-sector.
  By Lemma~\ref{lem:transversality}, each is
  suppressed by the sector angle $\theta_j$.
  The dominant contribution to stretching is
  geometrically depleted.
\item \textbf{Spread configurations}
  ($\cC^{(1)}\approx 1/J$): most triads are
  cross-sector.  These have no transversality
  suppression, but they are precisely the
  interactions controlled by the spreading mechanism
  (N7: cross-sector transfer dominates spreading
  by $830\times$; N20: inter-scale anti-correlation
  drives $\Sigma^{(1)}<0$).
\end{itemize}
In both regimes, the trilinear form is constrained
relative to its generic bound.  The
within-sector suppression factor $\theta_j$ is the
analytical counterpart of the numerical finding N7
(cross-sector dominance of spreading): cross-sector
transfer dominates \emph{because} within-sector
transfer is transversality-depleted.

For an adaptive partition aligned with the
concentration direction, the sector angle $\theta_j$
can be taken proportional to the concentration
aperture $\delta$ (the smallest cone half-angle
containing half the enstrophy).  In this case
the within-sector suppression improves as
concentration narrows.
This provides a structural link between
Fourier-space angular concentration and the
suppression of the dominant trilinear interactions,
which is the mechanism that eigenvalue selection
(Section~\ref{sec:alignment-experiment}) quantifies.
\end{remark}

The transversality suppression is not limited to
within-sector triads.  The divergence-free constraint
imposes a \emph{global} angular weight on every
trilinear interaction, which the standard estimates
discard.

\begin{corollary}[Global transversality weight]
\label{cor:global-transversality}
For any smooth divergence-free $u$ on $\T^3$:
\begin{equation}\label{eq:W-weighted}
|\mathcal{W}| \leq \sum_k |k|^2
\sum_{p+q=k} |q|\,\sin\angle(p,q)\,
|\hat u(p)|\,|\hat u(k)|\,|\hat u(q)|
=: \cB_{\rm w}(u).
\end{equation}
The standard (unweighted) trilinear bound is
\begin{equation}\label{eq:W-unweighted}
|\mathcal{W}| \leq \sum_k |k|^2
\sum_{p+q=k} |q|\,|\hat u(p)|\,|\hat u(k)|\,|\hat u(q)|
=: \cB_0(u).
\end{equation}
The ratio $\cB_{\rm w}/\cB_0 \leq 1$ always,
with equality only when all triads have
$p\perp q$.  For a field with enstrophy
concentrated in a cone of half-angle~$\theta$:
$\cB_{\rm w}/\cB_0 \leq \sin(2\theta)
\leq 2\theta$.
\end{corollary}

\begin{proof}
From Proposition~\ref{prop:factored},
$T(k) = 2|k|\,\im[\sum_{p+q=k}
(q\cdot\hat u(p))(\hat u(k)^*\cdot\hat u(q))]$.
Since $\hat u(p)\perp p$, write
$q = (q\cdot\hat p)\hat p + q_{\perp p}$;
then $q\cdot\hat u(p) = q_{\perp p}\cdot\hat u(p)$,
so $|q\cdot\hat u(p)| \leq |q_{\perp p}|\,|\hat u(p)|
= |q|\sin\angle(p,q)\,|\hat u(p)|$.
Summing with the $|k|^{2s}$ weight at $s=1$ and
using $T_{\rm tot}^{(1)}=2\mathcal{W}$
gives~\eqref{eq:W-weighted}.
The cone bound follows from
$\angle(p,q)\leq 2\theta$ for all $p,q$ in the cone.
\end{proof}

\begin{remark}[Why the standard estimates miss
the angular structure]
\label{rem:standard-miss}
The passage from $\cB_{\rm w}$ to $\cB_0$ (dropping
the $\sin\angle(p,q)$ factor) is the step where
angular information is lost.  The standard trilinear
estimates
($|\mathcal{W}|\leq C\|\omega\|\,\|\nabla\omega\|^2$)
use $\cB_0$ and therefore cannot distinguish between
angularly concentrated and angularly spread fields.
Any improvement to the standard enstrophy estimate
must exploit the $\sin\angle(p,q)$ weight --- it is
the only place where the angular geometry of
the velocity field enters the stretching bound.

The self-regulating structure is now visible at the
level of individual triads:
\begin{itemize}[leftmargin=2em]
\item \textbf{Aligned triads} ($\angle(p,q)\ll 1$):
  suppressed by $\sin\angle(p,q)\approx\angle(p,q)$.
  These dominate in concentrated configurations.
\item \textbf{Transverse triads} ($\angle(p,q)\sim 1$):
  no suppression ($\sin\angle(p,q)\sim 1$), but the
  inter-scale anti-correlation (N20) and spreading
  mechanism drive their net contribution toward
  angular redistribution rather than coherent
  stretching.
\end{itemize}
The spreading conjecture asserts that the NS dynamics
push the enstrophy toward isotropic angular
distributions where the transverse triads dominate.
But transverse triads contribute to the
\emph{weighted} bound $\cB_{\rm w}$ at full strength
--- so spreading does not suppress the trilinear
form through the sin weight alone.  The additional
ingredient is the \emph{sign structure}: transverse
triads have $\sin\angle(p,q)\sim 1$ but their
contributions are mutually incoherent (random phases),
producing cancellations in the \emph{signed} sum
$\mathcal{W}$ that the absolute-value bound
$\cB_{\rm w}$ does not capture.
This phase cancellation is the content of N20
(anti-correlation ratio $-263\times$) and is the
mechanism that eigenvalue selection
(Section~\ref{sec:alignment-experiment}) quantifies.
\end{remark}

\subsection{Experimental test of the alignment
mechanism (N26--N28)}
\label{sec:alignment-experiment}

We tested the alignment deviation framework
numerically ($128^3$, $\nu=0.005$, 10~random ICs
evolved to $t=5$, plus surgical concentration
at 10~levels~$\times$~3~sectors).

\textbf{N26 (Effective transversality).}
The triad-level effective transversality
$\Theta(u)=\cB_{\rm w}/\cB_0$
(Definition~\ref{def:Theta}), estimated by
Monte Carlo sampling of $10^6$ triads, rises from
$0.83$ (random phases at $t=0$) to
$0.95$--$0.99$ within one turnover time and
remains there.  For developed NS turbulence,
$\Theta\approx 0.97$: the transversality weight
provides negligible improvement over the standard
trilinear bound.  The triads that dominate
developed turbulence are \emph{not} nearly parallel.

\textbf{N27 (Eigenvalue selection).}
The enstrophy-weighted normalized stretching rate
$\langle\alpha/|S|\rangle_{\rm enst}
=\langle\hat\omega\cdot S\hat\omega/|S|\rangle$
rises from $\approx 0$ (random phases) to
$\approx 0.19$ by $t\approx 2$ and saturates,
consistently across all 10~seeds
(range: $0.177$--$0.195$ at $t=5$).
The theoretical maximum is $1$ (vorticity aligned
with the most extensional eigenvector).
NS turbulence depletes $\mathbf{81\%}$ of the
potential vortex stretching through eigenvalue
selection: vorticity preferentially aligns with
the intermediate strain eigenvector
($\cos^2\theta_2$ increases from $0.34$
at $t=0$ to $0.52$--$0.59$ at $t=5$), which
carries the smallest eigenvalue
($\lambda_2/|S|\approx 0.13$).

\textbf{N28 (Stretching depletion under
concentration).}
Under surgical angular concentration at fixed
energy (developed TG field, enstrophy pushed into
a single sector):
\begin{center}
\renewcommand{\arraystretch}{1.2}
\begin{tabular}{@{}cccc@{}}
\toprule
$\cC^{(1)}$ & $\alpha/|S|$
  & $\lambda_2/|S|$ & $\cos^2\theta_2$ \\
\midrule
$0.036$ & $0.134$ & $0.114$ & $0.52$ \\
$0.127$ & $0.121$ & $0.102$ & $0.58$ \\
$0.282$ & $0.104$ & $0.088$ & $0.61$ \\
$0.513$ & $0.081$ & $0.070$ & $0.67$ \\
$0.819$ & $0.044$ & $0.041$ & $0.79$ \\
\bottomrule
\end{tabular}
\end{center}
The normalized stretching rate $\alpha/|S|$
\emph{decreases monotonically} with angular
concentration~$\cC^{(1)}$.  The functional form
is approximately $\alpha/|S|\approx
c_0\sqrt{1-\cC^{(1)}}$ with $c_0\approx 0.14$.
The mechanism is eigenvalue selection: as
$\cC^{(1)}\to 1$, $\cos^2\theta_2\to 1$
(vorticity fully aligned with the intermediate
eigenvector), and $\lambda_2/|S|\to 0$ (the
intermediate eigenvalue is the smallest), so
$\alpha\to\lambda_2\to 0$ even though $D\to
\sqrt{2/3}$ (maximal alignment deviation).
Consistent results across three target sectors.

\begin{remark}[Why $D$ was the wrong target]
\label{rem:wrong-target}
The alignment deviation
$D=(\sum(\cos^2\theta_i-1/3)^2)^{1/2}$ measures
\emph{how anisotropic} the vorticity--strain
alignment is, not \emph{how much stretching}
it produces.  The NS dynamics drive
$D\to\sqrt{2/3}$ (maximal anisotropy) by aligning
$\omega$ with a single eigenvector --- but it is
the \emph{weakest} eigenvector, so stretching
is depleted despite maximal $D$.  The
correct depletion mechanism is eigenvalue
selection ($\alpha/|S|\ll 1$), not alignment
isotropy ($D\approx 0$).  The transfer inequality
conjecture ($\|D\|\leq h(\cC^{(1)})$) was tested
and found to be false: $D$ \emph{increases}
with~$\cC^{(1)}$.
\end{remark}

\begin{remark}[The eigenvalue selection ratio
as the degree-0 depletion]
\label{rem:eigenvalue-selection}
The ratio $\rho(u):=\langle\alpha/|S|\rangle_{\rm enst}
\in[-1,1]$ is degree~0 (dimensionless) and
measures the fraction of potential stretching
that is realized.  The experimental values:
\begin{itemize}[leftmargin=2em]
\item Random phases: $\rho\approx 0$ (no
  coherent stretching).
\item Developed NS turbulence: $\rho\approx 0.19$
  (eigenvalue selection depletes $81\%$).
\item Under surgical concentration:
  $\rho\approx 0.14\sqrt{1-\cC^{(1)}}$
  (additional depletion from concentration).
\end{itemize}
The vortex stretching satisfies
$\mathcal{W} = \int|\omega|^2\alpha\,dx
\approx\rho\int|\omega|^2|S|\,dx$, so the
enstrophy equation becomes
\begin{equation}\label{eq:enst-depleted}
\frac{d}{dt}\|\omega\|^2 \leq
2\rho\int|\omega|^2|S|\,dx
- 2\nu\|\nabla\omega\|^2.
\end{equation}
With $\rho\approx 0.19$ (a factor $5\times$ below
the worst case), the effective nonlinearity is
substantially weaker than the standard bound ---
but the exponent remains cubic ($\int|\omega|^2|S|$
is degree~3), so regularity does not follow from
eigenvalue selection alone.

The open question is whether
$\rho$ remains bounded below~$1$ along all NS
trajectories.  The numerical evidence
(N27: saturation at $0.19$ across 10~seeds)
and the eigenvalue selection mechanism
(intermediate alignment is a well-known
universal feature of NS turbulence) both
support this, but a proof would require
understanding \emph{why} the NS dynamics
select the intermediate eigenvector --- an
open problem first identified
experimentally by Ashurst, Kerstein, Kerr,
and Gibson (1987) and studied analytically
by Galanti and
Gibbon~\cite{GalantiGibbon1997}.
\end{remark}

This connects to Miller~\cite{Miller2024}, who
proved global regularity for the strain-vorticity
interaction model (which retains the eigenvalue
selection mechanism) and derived regularity criteria
for the full NS equation involving the alignment
of advection with strain self-interaction.
The experimental finding $\rho\approx 0.19$ is a
quantitative input to Miller's framework.

\subsection{Quantitative consequence: improved
enstrophy bound}
\label{sec:cone-improvement}

The transversality weight yields a concrete,
\emph{unconditional} improvement on the classical
enstrophy estimate.

\begin{definition}[Effective transversality]
\label{def:Theta}
For a smooth divergence-free field $u$ on $\T^3$
with $\cB_0(u)>0$, define
\begin{equation}\label{eq:Theta-def}
\Theta(u) := \frac{\cB_{\rm w}(u)}{\cB_0(u)}
\in [0,1],
\end{equation}
the ratio of the
weighted~\eqref{eq:W-weighted} to
unweighted~\eqref{eq:W-unweighted}
trilinear bounds.  This is degree~0
(dimensionless), depends only on the angular
structure of $u$, and satisfies $\Theta<1$
for any field that is not isotropically
distributed in direction.
\end{definition}

\begin{proposition}[Unconditional improved
stretching bound]
\label{prop:general-stretching}
For any smooth divergence-free $u$ on $\T^3$:
\begin{equation}\label{eq:W-general}
|\mathcal{W}(u)| \leq
C\,\Theta(u)\,\|\omega\|_{L^2}^{3/2}\,
\|\nabla\omega\|_{L^2}^{3/2},
\end{equation}
where $C$ is the constant in the standard
trilinear estimate
$|\mathcal{W}|\leq C\|\omega\|^{3/2}
\|\nabla\omega\|^{3/2}$.
\end{proposition}

\begin{proof}
$|\mathcal{W}| \leq \cB_{\rm w}$
(Corollary~\ref{cor:global-transversality})
$= \Theta\cdot\cB_0
\leq \Theta\cdot C\|\omega\|^{3/2}
\|\nabla\omega\|^{3/2}$
(standard trilinear: H\"older with
$\|\omega\|_{L^3}^2\|S\|_{L^3}$ and
Gagliardo--Nirenberg).
\end{proof}

\begin{corollary}[Cone-supported fields]
\label{cor:cone}
For fields with Fourier support in a cone of
half-angle~$\theta$: $\Theta(u)\leq\sin(2\theta)
\leq 2\theta$ for $\theta\leq\pi/4$.
\end{corollary}

\begin{proof}
Every triad has $\angle(p,q)\leq 2\theta$, so
$\sin\angle(p,q)\leq\sin(2\theta)$ in every term
of $\cB_{\rm w}$.
\end{proof}

\begin{proposition}[Unconditional enstrophy
improvement]
\label{prop:enst-general}
For any smooth solution of~\eqref{eq:NS} on
$\T^3$:
\begin{equation}\label{eq:enst-general}
\frac{d}{dt}\|\omega\|_{L^2}^2 \leq
-\nu\|\nabla\omega\|_{L^2}^2
+ \frac{C'\,\Theta(u)^4}{\nu^3}\,
\|\omega\|_{L^2}^6,
\end{equation}
where $C'$ depends only on the domain.
Enstrophy is decreasing whenever
\begin{equation}\label{eq:threshold-general}
\|\omega(t)\|_{L^2}^2 <
\frac{c\,\nu^2}{\Theta(u(t))^2}.
\end{equation}
This reduces to the standard threshold
$\|\omega\|^2<c\nu^2$ when $\Theta=1$
(isotropic), and improves to
$\|\omega\|^2<c\nu^2/\sin^2(2\theta)$
for cone-supported fields.
\end{proposition}

\begin{proof}
From~\eqref{eq:W-general} and the enstrophy
equation, by Young's inequality
($ab\leq\varepsilon\,a^{4/3}+C_\varepsilon\,b^4$
with $a=\|\nabla\omega\|^{3/2}$,
$b=\Theta\|\omega\|^{3/2}$):
$2C\Theta\|\omega\|^{3/2}\|\nabla\omega\|^{3/2}
\leq\nu\|\nabla\omega\|^2
+ C'\Theta^4\nu^{-3}\|\omega\|^6$.
Combining with Poincar\'e gives the threshold.
\end{proof}

\begin{proposition}[Unconditional blow-up
concentration budget]
\label{prop:budget}
If a smooth solution of~\eqref{eq:NS} on
$\T^3$ blows up at time $T^*<\infty$, then:
\begin{equation}\label{eq:budget}
\int_0^{T^*}\!\Theta\bigl(u(t)\bigr)^4\,dt
\geq \frac{\nu^3}{2C'\,\|\omega_0\|_{L^2}^4}.
\end{equation}
No assumption on angular support is required.
\end{proposition}

\begin{proof}
From~\eqref{eq:enst-general}, dropping dissipation:
$y' \leq C'\Theta^4\nu^{-3}y^3$ where
$y=\|\omega\|^2$.  Then
$\frac{d}{dt}(y^{-2})
\geq -2C'\Theta^4\nu^{-3}$.
Integrate to $T^*$ (where $y^{-2}\to 0$) and
rearrange.
\end{proof}

\begin{remark}[The effective transversality
$\Theta$ as the degree-0 obstruction]
\label{rem:Theta-meaning}
The functional $\Theta(u)\in[0,1]$ is the
dimensionless, degree-0 quantity that the standard
trilinear estimates discard.  It captures the
angular geometry of the velocity field at the
level of individual triads --- the only place where
the divergence-free constraint enters the
stretching bound.

Properties of $\Theta$:
\begin{itemize}[leftmargin=2em]
\item $\Theta=0$ iff $u=0$.
\item $\Theta=\sin(2\theta)$ for cone-supported
  fields of half-angle $\theta$.
\item $\Theta(\lambda u)=\Theta(u)$ for all
  $\lambda>0$ (degree~0).
\item $\Theta$ is computable from the Fourier data.
\end{itemize}

For cone-supported fields with
$\theta=0.1$\,rad: $\Theta\leq 0.20$,
giving a blow-up time improvement of
$\Theta^{-4}\approx 625$ over the standard bound.
For $\theta=0.01$\,rad: $\Theta\leq 0.020$,
improvement $\approx 6.25\times 10^6$.
\end{remark}

\begin{remark}[Concentration--compactness
and the pincer]
\label{rem:conc-compact}
By the concentration-compactness framework
(Gallagher--Koch--Planchon~\cite{GKP2013}), a
minimal blow-up element of NS must be almost
periodic, which forces angular concentration.
But the budget~\eqref{eq:budget}
says that blow-up requires
$\int\Theta^4\,dt$ to be sufficiently large,
and $\Theta$ is small precisely when
the field is angularly concentrated.
The spreading conjecture
(which asserts that the NS dynamics resist
angular concentration) would make $\Theta$
bounded away from zero, filling the budget
immediately and giving regularity.

The pincer: blow-up requires sustained
small $\Theta$ (concentration-compactness),
the budget requires sustained large $\Theta$
(integral constraint), and the spreading
dynamics resist small $\Theta$ (dynamical
instability of concentration).  The three
constraints are mutually opposing.
\end{remark}

%======================================================================
\section{Summary}
%======================================================================

\begin{center}
\renewcommand{\arraystretch}{1.2}
\begin{tabular}{@{}lll@{}}
\toprule
\textbf{Result} & \textbf{Status} & \textbf{Section} \\
\midrule
$T(k;\bar u) = -T(k;u)$
  & Proved & \ref{sec:general} \\
$\E_{\bar u=u}[T(k)]=0$, $\E[\Sigma]=0$
  & Proved & \ref{sec:general} \\
$T(k)$ is pressure-invariant
  & Proved & \ref{sec:general} \\
$\Sigma(\lambda u)=\lambda\Sigma(u)$
  & Proved & \ref{sec:general} \\
$V(k)\cdot\widetilde R(k)=-1/512$ (TG)
  & Proved & \ref{sec:TG} \\
$T^{(s)}(k)=|k|^{2(s-1/2)}T(k)$ (Sobolev scaling)
  & Proved & \ref{sec:general} \\
All algebraic results extend to $\Sigma^{(s)}$, all $s\geq 0$
  & Proved & \ref{sec:general} \\
$T(k)$ real sum for $\hat u=iV$ (real transfer formula)
  & Proved & \ref{sec:general} \\
TG orbit: transfer is real for all $t\geq 0$
  & Proved & \ref{sec:TG} \\
$\Sigma = (2/E)\sum\varepsilon_j T_j
  - (2T_{\rm tot}/E)\sum\varepsilon_j^2$ (exact)
  & Proved & \ref{sec:general} \\
$\Sigma<0$ along TG orbit
  & Numerical & \ref{sec:results} \\
$\nu$-independence of spreading
  & Numerical & \ref{sec:results} \\
Cross-sector dominance ($830\times$)
  & Numerical & \ref{sec:results} \\
Jet concentrates despite mode-level depletion
  & Numerical & \ref{sec:multi-ic} \\
Triad phases locked to $\pm\pi/2$ (NS only)
  & Numerical & \ref{sec:triad-phase} \\
$\Sigma/(\cC-1/J)$ independent of $\nu$
  & Numerical & \ref{sec:constitutive} \\
$\Sigma^{(s)}<0$ for $s\in[0,3/2]$
  & Numerical & \ref{sec:Hs-spreading} \\
Stretching vs.\ spreading ($+0.14$ vs.\ $-0.43$)
  & Numerical & \ref{sec:stretching} \\
$\Sigma^{(1)}/(\cC^{(1)}-1/J)$ indep.\ of $\nu$
  & Numerical & \ref{sec:enst-visc} \\
$\Sigma^{(1)}<0$ for TG + all 3 random ICs
  & Numerical & \ref{sec:enst-multi-ic} \\
$\kappa_{\rm align}$ sign transition at $\cC^{(1)}\approx 0.15$
  & Numerical & \ref{sec:Phi-landscape} \\
$\Sigma^{(1)}$ robust to 50\% phase randomization
  & Numerical & \ref{sec:robustness} \\
Surgical: $\Sigma^{(1)}<0$ at all 44 snapshots
  & Numerical & \ref{sec:surgical} \\
$\mathcal{R}\in[0.6,2.2]$ under adversarial concentration
  & Numerical & \ref{sec:surgical} \\
$g^{(1)}\approx 0.47\pm 0.04$ (constitutive law)
  & Numerical & \ref{sec:constitutive-surgical} \\
Randomizations give 50\% ($\Sigma^{(1)}<0$ not static)
  & Numerical & \ref{sec:dynamical} \\
Inter-scale anti-correlation ($-263\times$, NS only)
  & Numerical & \ref{sec:dynamical} \\
Spreading is a dynamical invariant of the NS flow
  & Numerical & \ref{sec:dynamical} \\
$50\%\to 100\%$: NS creates spreading (30 seeds, $128^3$)
  & Numerical & \ref{sec:creation} \\
Contraction asymmetry: $\Sigma^{(1)}>0$ basin unstable
  & Numerical & \ref{sec:creation} \\
High-res creation ($256^3$, 10 seeds) confirms $128^3$
  & Numerical & \ref{sec:creation-highres} \\
Creation timescale $\nu$-independent ($\times 4$ viscosities)
  & Numerical & \ref{sec:creation-viscosity} \\
$-i$ rotation mechanism
  & Heuristic & \ref{sec:mechanism} \\
$\im[\gamma]<0$ sign criterion
  & Conditional & \ref{sec:conditional} \\
Conjecture~\ref{conj:enstrophy} ($s=1$ spreading)
  & Conjecture & \ref{sec:conjecture} \\
Stretching--spreading coupling ($\mathcal{R}$ bounded)
  & Numerical & \ref{sec:surgical} \\
$g^{(1)}\geq c > 0$ (dynamical invariant)
  & Open & \ref{sec:dynamical} \\
$g^{(1)}$ is degree~1, not dimensionless (N25)
  & Proved & \ref{sec:homogeneity} \\
Viscous--nonlinear balance $|\text{visc}/\text{NL}|\approx 1$
  & Numerical & \ref{sec:visc-balance} \\
Homogeneity obstruction for direct regularity
  & Proved & \ref{sec:homogeneity} \\
Bakry--\'Emery dictionary (with degree mismatch)
  & Structural & \ref{sec:bakry-emery} \\
$\alpha=\sum\lambda_i(\cos^2\theta_i-1/3)$
  (isotropic alignment kills stretching)
  & Proved & \ref{sec:alignment} \\
$T_{\rm tot}^{(1)}=2\mathcal{W}$
  (sector transfer $=$ vortex stretching)
  & Proved & \ref{sec:alignment} \\
Stretching--spreading identity~\eqref{eq:W-Sigma}
  & Proved & \ref{sec:alignment} \\
Implicit transfer inequality~\eqref{eq:D-constraint}
  & Proved & \ref{sec:alignment} \\
Within-sector depletion by $\theta_j$
  (transversality)
  & Proved & \ref{sec:alignment} \\
$|\mathcal{W}|\leq C\Theta\|\omega\|^{3/2}
  \|\nabla\omega\|^{3/2}$ (unconditional)
  & Proved & \ref{sec:cone-improvement} \\
Enstrophy threshold $\|\omega\|^2<c\nu^2/\Theta^2$
  (unconditional)
  & Proved & \ref{sec:cone-improvement} \\
Blow-up budget $\int\Theta^4\,dt\geq c\nu^3/\|\omega_0\|^4$
  (unconditional)
  & Proved & \ref{sec:cone-improvement} \\
$\Theta\approx 0.97$ for developed turbulence (N26)
  & Numerical & \ref{sec:alignment-experiment} \\
$\alpha/|S|\approx 0.19$ (eigenvalue selection, N27)
  & Numerical & \ref{sec:alignment-experiment} \\
$\alpha/|S|\approx 0.14\sqrt{1-\cC^{(1)}}$ (surgical, N28)
  & Numerical & \ref{sec:alignment-experiment} \\
\midrule
\multicolumn{3}{@{}l@{}}{\emph{Affine ratio $c^{(1)}$
  (Section~\ref{sec:affine-experiments})}} \\
$\Sigma^{(1)}=2(c^{(1)}\!-\!1)\,T_{\rm tot}\sum\varepsilon_j^2/E$
  & Proved & \ref{sec:general} \\
$c^{(1)}(\lambda u)=c^{(1)}(u)$ (degree~0)
  & Proved & \ref{sec:general} \\
$c^{(1)}<1$ along TG for $t\geq 1$ (N30)
  & Numerical & \ref{sec:affine-experiments} \\
$c^{(1)}<1$ under surgical concentration (N31)
  & Numerical & \ref{sec:affine-experiments} \\
$c^{(1)}<1$ not kinematic: only 12\% of random $V$ (N32)
  & Numerical & \ref{sec:affine-experiments} \\
$\{c^{(1)}<1\}$ absorbing basin ($5/5$, N33)
  & Numerical & \ref{sec:affine-experiments} \\
(RTI) $\Rightarrow$ $c_{\rm cen}\leq 1$ (Chebyshev)
  & Proved & \ref{sec:centroid-inequality} \\
Two-shell identity: $c_{\rm cen}<1$ iff $\Phi_L>\Phi_H$
  & Proved & \ref{sec:centroid-inequality} \\
$c_{\rm cen}<1$ from viscous balance
  (Theorem~\ref{thm:viscous-sep})
  & Proved & \ref{sec:centroid-inequality} \\
NS violates (RTI) but satisfies $\Phi_L>\Phi_H$ (N35)
  & Numerical & \ref{sec:centroid-inequality} \\
Forcing $B(u,u)$ anisotropic at high $k$ (N35b)
  & Numerical & \ref{sec:centroid-inequality} \\
$\Sigma^{(1)}=d\sigma/dt-4\nu\sigma D(1-c_{\rm cen})/E$
  & Proved & \ref{sec:concentration-dynamics} \\
$\Sigma^{(1)}<0$ at concentration equilibrium
  & Proved & \ref{sec:concentration-dynamics} \\
Time-averaged $\Sigma^{(1)}<0$
  & Proved & \ref{sec:concentration-dynamics} \\
$c^{(1)}\geq 1$ cannot persist (self-correcting)
  & Proved & \ref{sec:concentration-dynamics} \\
Transient $c>1$ during cascade development (N36)
  & Numerical & \ref{sec:concentration-dynamics} \\
$c\geq 1$ $\Rightarrow$ $dE/dt>0$ (92/92, N37a)
  & Numerical & \ref{sec:concentration-dynamics} \\
$R_{\rm NL}^{(L)}<0$ always; excursions from
  $R_{\rm NL}^{(H)}>0$ (N37b)
  & Numerical & \ref{sec:concentration-dynamics} \\
$R_{\rm NL}^{(L)}\approx 0$ at late times;
  high shells dominate $1000{:}1$ (N38)
  & Numerical & \ref{sec:concentration-dynamics} \\
$\corr(\varepsilon,\delta^{(H)})<0$ for $t\geq 4$;
  gap $c-c_{\rm cen}$ always helpful (N39)
  & Numerical & \ref{sec:concentration-dynamics} \\
Time-averaged spreading (Conjecture~\ref{conj:time-avg})
  & Proved & \ref{sec:concentration-dynamics} \\
\bottomrule
\end{tabular}
\end{center}

\subsection*{The road to regularity}

The paper establishes a chain of proved algebraic
results, unconditional analytical bounds, and
numerical measurements.  We state what is proved,
what the experiments reveal, and where the gap is.

\medskip\noindent\textbf{Proved (analytical):}
\begin{enumerate}[leftmargin=2em,label=\textbf{(\arabic*)}]
\item Complete algebraic framework for angular
  enstrophy transfer: sector decomposition,
  spreading functional $\Sigma^{(1)}$, linear
  response, Sobolev scaling, real transfer formula
  (Section~\ref{sec:general}).
\item Stretching--spreading identity:
  $T_{\rm tot}^{(1)}=2\mathcal{W}$, connecting
  Fourier-space sector transfer to physical-space
  vortex stretching
  (Proposition~\ref{prop:stretching-spreading}).
\item Alignment deviation identity:
  $\alpha=\sum\lambda_i(\cos^2\theta_i-1/3)$;
  isotropic alignment kills stretching
  (eq.~\eqref{eq:alpha-deviation}).
\item Global transversality weight: every trilinear
  coupling carries $\sin\angle(p,q)$ from the
  divergence-free constraint
  (Corollary~\ref{cor:global-transversality}).
\item Homogeneity obstruction: $g^{(1)}$ is degree~1,
  not dimensionless; angular spreading controls
  direction but not magnitude
  (Section~\ref{sec:homogeneity}).
\item Affine ratio decomposition:
  $\Sigma^{(1)}=2(c^{(1)}-1)\,T_{\rm tot}
  \sum\varepsilon_j^2/E$ with $c^{(1)}$ degree~0
  (Proposition~\ref{prop:c-degree-zero}), isolating
  the homogeneity obstruction in the prefactor
  $T_{\rm tot}/E$ (Section~\ref{sec:general}).
\item Spectral centroid inequality
  (Theorem~\ref{thm:centroid}): (RTI) implies
  $c_{\rm cen}\leq 1$ via Chebyshev.  Two-shell
  identity (Proposition~\ref{prop:two-shell}):
  $c_{\rm cen}<1$ iff $\Phi_L>\Phi_H$.
\item Spectral separation from viscous balance
  (Theorem~\ref{thm:viscous-sep}): $c_{\rm cen}<1$
  for all smooth NS solutions with
  $\mathcal{V}^{(L)}(0)>0$, for
  $t\leq O(1/(\nu\kappa_L))$.  Uses only the NS
  equation, Gronwall, Cauchy--Schwarz, and the
  two-shell identity.
\item Concentration--spreading identity
  (Theorem~\ref{thm:conc-spread}):
  $\Sigma^{(1)} = d\sigma/dt
  - 4\nu\sigma D(1-c_{\rm cen})/E$.
  Consequences:
  $c^{(1)}<1$ at every concentration equilibrium
  (Theorem~\ref{thm:equil});
  time-averaged $\Sigma^{(1)}<0$
  (Theorem~\ref{thm:time-avg});
  $c^{(1)}\geq 1$ cannot persist
  (Remark~\ref{rem:self-reinforce}).
  (Section~\ref{sec:concentration-dynamics}).
\end{enumerate}

\noindent\textbf{Measured (numerical):}
\begin{enumerate}[leftmargin=2em,label=\textbf{(N\arabic*$'$)}]
\item $\Sigma^{(1)}<0$ is a dynamical attractor:
  100\% of seeds reach spreading by $t=1$, robust
  across resolutions and viscosities (N21--N24).
\item Effective transversality $\Theta\approx 0.97$
  for developed turbulence: the triad-level
  sin-weight provides negligible improvement over
  the standard bound (N26).
\item Eigenvalue selection ratio $\rho=\alpha/|S|
  \approx 0.19$: NS depletes $81\%$ of potential
  stretching by aligning vorticity with the
  intermediate strain eigenvector (N27).
\item Stretching depletion under concentration:
  $\alpha/|S|\approx 0.14\sqrt{1-\cC^{(1)}}$,
  monotonically decreasing (N28).
\item Affine ratio $c^{(1)}\in[0.18, 0.69]$ for
  developed TG turbulence ($t\geq 5$); $c^{(1)}<1$
  under surgical concentration up to
  $\cC^{(1)}=0.124$; the basin $\{c^{(1)}<1\}$ is
  absorbing under NS dynamics (5/5 each direction,
  N29--N33).
\item High-resolution sampling (N36: 1000 snapshots,
  $\Delta t=0.01$, $t\in[0,10]$): transient
  excursions $c^{(1)}>1$ during cascade development
  ($t\approx 3.5$--$4.8$), with $\max c = 1.23$ at
  $\nu=0.005$ (92/1000 snapshots above~$1$).
  At $\nu=0.02$: $c<1$ at all 1000 snapshots.
  The excursions are confined to the cascade
  development phase and do not persist, consistent
  with the proved self-correction
  (Remark~\ref{rem:self-reinforce}).
\end{enumerate}

\noindent\textbf{The gap:}
\begin{enumerate}[leftmargin=2em,label=\textbf{(G\arabic*)}]
\item The pointwise bound $c^{(1)}<1$ is false
  during cascade development (N36).  The high-shell
  residual correlation (N39) reveals the precise
  mechanism: $\corr(\varepsilon,\delta^{(H)})<0$
  for all $t\geq 4$, so the gap $c-c_{\rm cen}$
  is \emph{always helpful} (the high-shell transfer
  spreads more than the dissipation).  The excursion
  $c>1$ at $t\approx 4.2$ occurs because
  $c_{\rm cen}$ itself exceeds~$1$ (the viscous
  balance margin is temporarily overwhelmed by
  explosive enstrophy growth at low~$k$).  After
  cascade development ($t\geq 5$), $c_{\rm cen}$
  drops to~$0.69$ and $c<c_{\rm cen}<1$ permanently.

  The remaining obstruction is purely in the
  viscous balance margin: proving that
  $c_{\rm cen}(t)<1$ holds \emph{during} explosive
  enstrophy growth, not just before and after.
  This requires controlling the growth of
  $\|u\|_{H^2}$ during cascade development ---
  a quantitative regularity condition.
\item Prove that $\rho(u(t))=
  \langle\alpha/|S|\rangle_{\rm enst}$ remains
  bounded below~$1$ along all smooth NS
  trajectories.
\end{enumerate}

The $\Theta$-route is closed: $\Theta\approx 0.97$
for developed turbulence (N26), so the blow-up
budget gives no improvement over the standard
bound.  The $D$-route (transfer inequality) is
closed: $D$ \emph{increases} with concentration
(first experiment), the opposite of what was
conjectured.

The actual depletion mechanism is eigenvalue
selection (N27--N28): vorticity aligns with the
intermediate strain eigenvector, which carries the
smallest eigenvalue.  This gives
$\mathcal{W}\approx\rho\int|\omega|^2|S|\,dx$ with
$\rho\approx 0.19$ --- a $5\times$ reduction in the
effective nonlinearity.  But $\rho<1$ does not change
the degree-3 exponent, so regularity does not follow
from $\rho<1$ alone.

The remaining question is whether the eigenvalue
selection mechanism can be combined with the
spreading mechanism to produce a \emph{dynamic}
bound that closes.  The spreading conjecture
(N21--N24) controls the angular distribution; the
eigenvalue selection (N27--N28) controls the
stretching efficiency.  Whether these two degree-0
properties can together overcome the degree-3
scaling is the open problem.

The affine ratio $c^{(1)}$
(Definition~\ref{def:affine-ratio}) offers a
structural advance on this question.
Conjecture~\ref{conj:affine} ($c^{(1)}<1$) is a
\emph{dimensionless} statement: it carries no
homogeneity obstruction
(Proposition~\ref{prop:c-degree-zero}).  The
full spreading rate decomposes as
$\Sigma^{(1)} = 2(c^{(1)}-1)\,T_{\rm tot}\,
\sum\varepsilon_j^2/E$, where $c^{(1)}<1$ is
the contribution of the angular spreading
program, and the degree-1 prefactor
$T_{\rm tot}/E$ is the classical stretching
control problem.  The experimental evidence
(N29--N33: $c^{(1)}$ is degree~0, $c^{(1)}<1$
is an absorbing basin under NS dynamics, and
the bound is not kinematic) identifies proving
$c^{(1)}<1$ as the natural analytical target ---
a target that is, for the first time in this
program, free of the homogeneity obstruction.

\subsection*{Acknowledgments}

AI tools were used as interactive assistants for drafting
and mathematical exploration.
Simulations were performed on an NVIDIA RTX~5090.

%======================================================================
\begin{thebibliography}{99}

\bibitem{Fefferman}
C.\,L.~Fefferman,
Existence and smoothness of the Navier--Stokes equation,
\emph{Clay Mathematics Institute Millennium Prize Problems}
(2000), 1--5.

\bibitem{Leray}
J.~Leray,
Sur le mouvement d'un liquide visqueux emplissant l'espace,
\emph{Acta Math.}~\textbf{63} (1934), 193--248.

\bibitem{ESS}
L.~Escauriaza, G.\,A.~Seregin, V.~\v{S}ver\'ak,
$L_{3,\infty}$-solutions of Navier--Stokes equations and backward
uniqueness, \emph{Russian Math.\ Surveys}~\textbf{58} (2003),
211--250.

\bibitem{BKM}
J.\,T.~Beale, T.~Kato, A.~Majda,
Remarks on the breakdown of smooth solutions for the 3D Euler
equations, \emph{Comm.\ Math.\ Phys.}~\textbf{94} (1984), 61--66.

\bibitem{Tao}
T.~Tao,
Finite time blowup for an averaged three-dimensional
Navier--Stokes equation,
\emph{J.\ Amer.\ Math.\ Soc.}~\textbf{29} (2016), 601--674.

\bibitem{Kolmogorov}
A.\,N.~Kolmogorov,
The local structure of turbulence in incompressible viscous fluid
for very large Reynolds numbers,
\emph{Dokl.\ Akad.\ Nauk SSSR}~\textbf{30} (1941), 299--303.

\bibitem{Rotta}
J.\,C.~Rotta,
Statistische Theorie nichthomogener Turbulenz,
\emph{Z.\ Physik}~\textbf{129} (1951), 547--572.

\bibitem{CF93}
P.~Constantin and C.~Fefferman,
Direction of vorticity and the problem of global
regularity for the Navier--Stokes equations,
\emph{Indiana Univ.\ Math.\ J.}~\textbf{42} (1993),
775--789.

\bibitem{Miller2024}
E.~Miller,
On the interaction of strain and vorticity for
solutions of the Navier--Stokes equation,
\emph{arXiv preprint} arXiv:2407.02691 [math.AP], 2024.

\bibitem{GKP2013}
I.~Gallagher, G.~S.~Koch, and F.~Planchon,
A profile decomposition approach to the
$L^\infty_t(L^3_x)$ Navier--Stokes regularity
criterion,
\emph{Math.\ Ann.}~\textbf{355} (2013), 1527--1559.

\bibitem{GalantiGibbon1997}
B.~Galanti, J.~D.~Gibbon, and M.~Heritage,
Vorticity alignment results for the three-dimensional
Euler and Navier--Stokes equations,
\emph{Nonlinearity}~\textbf{10} (1997), 1675--1694.

\bibitem{Fontes2026b}
R.~Fontes,
Bakry--\'Emery curvature of the fractional Laplacian
via fractional Brownian covariance,
\emph{Preprint}, 2026.

\bibitem{Fontes2026c}
R.~Fontes,
The iterated $\Gamma$-calculus for the fractional
Laplacian, II: operator realization and Wiener chaos,
\emph{Preprint}, 2026.

\end{thebibliography}

\end{document}
